
\documentclass[11pt]{book}
\usepackage[a4paper,margin=1.15in]{geometry}
\usepackage[T1]{fontenc}
\usepackage[utf8]{inputenc}
\usepackage{lmodern}
\usepackage{microtype}
\usepackage{amsmath,amssymb,mathtools}
\usepackage{longtable,array,booktabs}
\usepackage{enumitem}
\usepackage{hyperref}
\hypersetup{
  hypertexnames=false,
  hidelinks,
  pdfauthor={Luigi Bianchi},
  pdftitle={Lezioni sulla teoria dei gruppi continui finiti di trasformazioni -- Italiano -- pagine manoscritte 1--116},
  pdfsubject={Trascrizione italiana cumulativa da un volume manoscritto pubblicato nel 1903}
}
\setlength{\parindent}{1.3em}
\setlength{\parskip}{0.2em}
\setlength{\emergencystretch}{4em}
\hbadness=10000
\newcommand{\sen}{\operatorname{sen}}
\newcommand{\dif}{\,d}
\newcommand{\entry}[3]{#1 & #2 & #3\\}
\begin{document}

\begin{titlepage}
\centering
{\Large LUIGI BIANCHI\par}
\vspace{0.4em}
{\small PROFESSORE DELLA REGIA UNIVERSITÀ DI PISA\par}
\vfill
{\Huge LEZIONI\par}
\vspace{1.5em}
{\large SULLA\par}
\vspace{1.5em}
{\Large Teoria dei gruppi continui finiti di trasformazioni\par}
\vspace{2em}
{\large Anno 1902--1903\par}
\vfill
{\Large PISA\par}
{\large ENRICO SPOERRI\par}
{\small LIBRAIO-EDITORE\par}
\vspace{0.5em}
{\large 1903\par}
\end{titlepage}

\frontmatter
\chapter*{Indice dei capitoli}
\footnotesize
\begin{longtable}{@{}>{\raggedright\arraybackslash}p{0.42in} >{\raggedright\arraybackslash}p{4.60in} >{\raggedleft\arraybackslash}p{0.32in}@{}}
\toprule
\textbf{§} & \textbf{Titolo} & \textbf{pag.}\\
\midrule
\endfirsthead
\toprule
\textbf{§} & \textbf{Titolo} & \textbf{pag.}\\
\midrule
\endhead
\multicolumn{3}{c}{\textsc{Capitolo I}}\\
\entry{}{\textit{Generalità sui gruppi continui finiti di trasformazioni. -- Equazioni differenziali fondamentali. -- Trasformazioni infinitesime come generatrici del gruppo.}}{}
\entry{1}{Trasformazioni e loro composizione}{1}
\entry{2}{Serie $\infty^r$ di trasformazioni}{5}
\entry{3}{Criterio per riconoscere il numero dei parametri essenziali}{10}
\entry{4}{Definizione dei gruppi continui finiti}{16}
\entry{5}{Esempî di gruppi continui}{20}
\entry{6}{Risolubilità delle equazioni $c_i=\varphi_i(a,b)$ rispetto alle $a$ o alle $b$}{24}
\entry{7}{Equazioni differenziali fondamentali}{27}
\entry{8}{Alcune proprietà delle funzioni $\psi_{sh}(a)$, $\xi_{si}(x')$}{31}
\entry{9}{Gruppi ad un parametro $G_1$}{34}
\entry{10}{Trasformazione infinitesima di un $G_1$}{38}
\entry{11}{Cangiamento di variabili nelle trasformazioni infinitesime}{42}
\entry{12}{Serie $\infty^r$ di trasformazioni e loro trasformazioni infinitesime}{43}
\entry{13}{Caso particolare di un gruppo}{47}
\entry{14}{Nuova deduzione delle equazioni fondamentali (A)}{50}
\entry{15}{Le $r$ trasformazioni infinitesime $X_\lambda f=\sum\limits_{i=1}^n \xi_{\lambda i}(x)\dfrac{\partial f}{\partial x_i}$}{52}
\entry{16}{Teorema fondamentale per le serie di trasformazioni soddisfacenti le (A)}{55}
\addlinespace[0.5em]
\multicolumn{3}{c}{\textsc{Capitolo II}}\\
\entry{}{\textit{I tre teoremi fondamentali della teoria dei gruppi. -- I gruppi misti.}}{}
\entry{17}{Il primo teorema fondamentale}{58}
\entry{18}{Casi ulteriori}{61}
\entry{19}{Sistema completo di equazioni a derivate parziali equivalente al sistema (A)}{63}
\entry{20--21}{Condizioni d'integrabilità $(\alpha)$ e $(\beta)$}{65}
\entry{22}{Dimostrazione di Schur del primo teorema fondamentale}{69}
\entry{23}{Il secondo teorema fondamentale (teorema principale)}{72}
\entry{24}{Le relazioni fra le costanti di composizione $C_{iks}$}{77}
\entry{25}{Proprietà del simbolo {\scriptsize $|UV|:=\sum\limits_{i,k,s=1}^{r}C_{iks}y_s\,\dfrac{\partial U}{\partial y_i}\,\dfrac{\partial V}{\partial y_k}$}}{80}
\entry{26}{Cangiamento di variabili}{83}
\entry{27}{Il terzo teorema fondamentale}{86}
\entry{28}{Le equazioni di Maurer per le $\varphi_{ik}$}{89}
\entry{29}{Dimostrazione di Schur del secondo e terzo teorema fondamentale}{92}
\entry{30}{Composizione di due serie $\infty^r$ di trasformazioni}{95}
\entry{31}{Serie $\infty^r$ di trasformazioni nelle quali i prodotti di $m$ loro trasformazioni appartengono alla serie}{100}
\entry{32}{I gruppi misti (o complessi)}{102}
\addlinespace[0.5em]
\multicolumn{3}{c}{\textsc{Capitolo III}}\\
\entry{}{\textit{Equazioni di definizione di un gruppo. -- Ordine delle trasformazioni infinitesime.}}{}
\entry{33}{Eliminazione di costanti arbitrarie da un sistema di funzioni}{107}
\entry{34}{Sistemi di equazioni a derivate parziali di Mayer e loro proprietà}{111}
\entry{35}{Le equazioni di definizione di un gruppo $G_r$}{116}
\entry{36}{Altro modo di definire un gruppo continuo. -- Cenno sui gruppi infiniti}{119}
\entry{37}{Ordine di una trasformazione infinitesima}{124}
\entry{38}{Distribuzione degli ordini delle trasformazioni infinitesime in un gruppo}{127}
\entry{39}{Calcolo dei numeri $r$ dalle equazioni di definizione del gruppo}{132}
\entry{40}{Sottogruppi generati dalle trasformazioni degli ordini superiori}{137}
\entry{41}{Esempio del gruppo di movimenti}{144}
\addlinespace[0.5em]
\multicolumn{3}{c}{\textsc{Capitolo IV}}\\
\entry{}{\textit{Funzioni e varietà invarianti rispetto ad un gruppo. -- Gruppi transitivi ed intransitivi. -- Gruppi imprimitivi. -- Gruppi sistatici ed asistatici.}}{}
\entry{42}{Funzioni invarianti rispetto ad un gruppo}{148}
\entry{43}{Digressione sulle varietà in uno spazio $S_n$ a $n$ dimensioni}{150}
\entry{44}{Varietà invarianti rispetto ad un gruppo $G_1$}{154}
\entry{45}{Varietà invarianti rispetto ad un gruppo $G_r$}{159}
\entry{46}{Dimensione delle varietà minime invarianti}{161}
\entry{47}{Determinazione di tutte le varietà minime invarianti}{164}
\entry{48}{Gruppi transitivi}{169}
\entry{49}{Invarianti dei gruppi intransitivi}{171}
\entry{50}{Esempî varii}{174}
\entry{51}{Gruppi imprimitivi}{178}
\entry{52}{Sistemi completi che ammettono un gruppo}{182}
\entry{53}{Gruppi sistatici ed asistatici}{186}
\entry{54}{Imprimitività di un gruppo sistatico}{190}
\entry{55}{Criterio per la sistaticità dedotto dalle equazioni di definizione}{193}
\addlinespace[0.5em]
\multicolumn{3}{c}{\textsc{Capitolo V}}\\
\entry{}{\textit{Serie invarianti di trasformazioni infinitesime. -- Sottogruppi invarianti e serie di composizione di un gruppo.}}{}
\entry{56}{Serie invarianti di trasformazioni infinitesime}{198}
\entry{57}{Serie invarianti rispetto ad un gruppo}{201}
\entry{58}{Trasformazioni finite della serie $\sum\limits_h^{1,\ldots,r} e_h X_h f$}{207}
\entry{59}{Trasformazioni permutabili}{209}
\entry{60}{Sottogruppi invarianti. -- Gruppi semplici}{212}
\entry{61}{Serie di composizione. -- Gruppi integrabili}{215}
\entry{62}{Gruppo derivato}{218}
\entry{63}{Condizioni d'integrabilità di un gruppo}{221}
\entry{64}{Prodotto di due gruppi permutabili}{224}
\entry{65}{Teorema relativo a due sottogruppi invarianti massimi}{226}
\entry{66}{Teorema di Jordan-Lie sugli indici delle serie di composizione}{229}
\addlinespace[0.5em]
\multicolumn{3}{c}{\textsc{Capitolo VI}}\\
\entry{}{\textit{Il gruppo aggiunto e i sottogruppi di un gruppo. -- Sottogruppi a due parametri. -- Equazione caratteristica. -- Composizione o struttura di un gruppo. -- Genere di un gruppo.}}{}
\entry{67}{Permutazioni fra le trasformazioni di un gruppo indotte dalle trasformazioni stesse}{232}
\entry{68}{Forma canonica delle equazioni di un gruppo}{235}
\entry{69}{Il gruppo aggiunto e le sue trasformazioni infinitesime}{239}
\entry{70}{Numero dei parametri del gruppo aggiunto}{243}
\entry{71}{Ricerca dei sottogruppi di un gruppo}{247}
\entry{72}{Sottogruppi affini}{249}
\entry{73}{Rappresentazione nello spazio ad $r$ dimensioni}{252}
\entry{74}{Esempio del gruppo di movimenti nel piano}{256}
\entry{75}{Applicazione al gruppo proiettivo sopra una variabile}{262}
\entry{76}{Sottogruppi a due parametri di un gruppo $G_r$}{266}
\entry{77}{L'equazione caratteristica}{271}
\entry{78}{Proprietà dell'equazione caratteristica}{276}
\entry{79}{I coefficienti di $\Delta(\rho)$ come invarianti del gruppo aggiunto}{279}
\entry{80}{Composizione (o struttura) di un gruppo}{282}
\entry{81}{Le relazioni fra le $C_{iks}$, $C'_{iks}$ come individuanti un gruppo}{285}
\entry{82}{Genere (Rang) di un gruppo}{287}
\addlinespace[0.5em]
\multicolumn{3}{c}{\textsc{Capitolo VII}}\\
\entry{}{\textit{Prime ricerche sulla similianza dei gruppi. -- Trasformazioni permutabili con tutte le trasformazioni di un gruppo. -- Gruppi semplicemente transitivi reciproci.}}{}
\entry{83}{Simiglianza dei gruppi}{291}
\entry{84}{Caso in cui la condizione d'eguale composizione è sufficiente}{296}
\entry{85}{Caso dei gruppi Abeliani e dei gruppi semplicemente transitivi}{302}
\entry{86}{Condizioni generali necessarie per la simiglianza}{303}
\entry{87}{Caso dei gruppi asistatici}{306}
\entry{88}{Simiglianza di un gruppo con sè stesso}{309}
\entry{89}{Trasformazioni permutabili colle singole trasformazioni di un gruppo}{312}
\entry{90}{Gruppi semplicemente transitivi reciproci}{317}
\entry{91}{Proprietà delle coppie di gruppi reciproci}{320}
\entry{92}{Costruzione diretta del gruppo reciproco $\Gamma_n$}{325}
\entry{93}{Imprimitività di un gruppo semplicemente transitivo}{328}
\addlinespace[0.5em]
\multicolumn{3}{c}{\textsc{Capitolo VIII}}\\
\entry{}{\textit{Isomorfismo dei gruppi riferito alle loro trasformazioni infinitesime. -- Gruppi parametrici. -- Isomorfismo dei gruppi in relazione colle trasformazioni finite. -- Risoluzione del problema generale di riconoscere la simiglianza di due gruppi.}}{}
\entry{94}{Isomorfismo di due gruppi riferito alle loro trasformazioni infinitesime}{331}
\entry{95}{Proprietà dell'isomorfismo}{337}
\entry{96}{Gruppo complementare rispetto ad un sottogruppo invariante}{340}
\entry{97}{Serie di composizione di due gruppi isomorfi. -- Teorema di Hölder-Engel}{345}
\entry{98}{Costruzione di gruppi isomorfi a gruppi primitivi}{350}
\entry{99}{Definizione dei gruppi parametrici}{354}
\entry{100}{Trasformazioni infinitesime del (primo) gruppo parametrico}{359}
\entry{101}{Gruppi parametrici equivalenti}{363}
\entry{102}{Isomorfismo oloedrico in relazione colle trasformazioni finite}{366}
\entry{103}{Isomorfismo meriedrico}{372}
\entry{104}{Ritorno alle condizioni di simiglianza per due gruppi}{376}
\entry{105}{Simiglianza fra gruppi transitivi}{379}
\entry{106}{Simiglianza fra gruppi intransitivi}{384}
\entry{107}{Nuovo enunciato per le condizioni di simiglianza}{391}
\entry{108}{Trasformazioni permutabili con un gruppo}{394}
\entry{109}{Teoremi sui gruppi asistatici e sistatici}{397}
\addlinespace[0.5em]
\multicolumn{3}{c}{\textsc{Capitolo IX}}\\
\entry{}{\textit{Gruppi ampliati e prolungati. -- Invarianti differenziali.}}{}
\entry{110}{Transitività multipla}{401}
\entry{111}{Invarianti di gruppi di punti}{404}
\entry{112}{Gruppi prolungati agli elementi lineari}{408}
\entry{113}{Trasformazioni infinitesime del gruppo prolungato}{412}
\entry{114}{Gruppo indotto dal sottogruppo di stabilità sugli elementi lineari per un punto}{415}
\entry{115}{Definizione degli invarianti differenziali}{419}
\entry{116}{Trasformazioni prolungate}{422}
\entry{117}{Gruppi prolungati}{426}
\entry{118}{Esistenza e calcolo degli invarianti differenziali}{431}
\entry{119}{Invarianti differenziali dei gruppi su due variabili}{433}
\entry{120}{Esempî di calcolo d'invarianti differenziali}{439}
\entry{121}{Invarianti differenziali di superficie rispetto ai movimenti}{443}
\addlinespace[0.5em]
\multicolumn{3}{c}{\textsc{Capitolo X}}\\
\entry{}{\textit{I gruppi continui finiti sopra una e sopra due variabili.}}{}
\entry{122}{Numero dei parametri di un gruppo sopra una variabile}{446}
\entry{123}{Riduzione dei gruppi sopra una variabile a gruppi proiettivi}{449}
\entry{124}{Gruppo lineare omogeneo $G_4$ su due variabili}{453}
\entry{125}{Sottogruppi invarianti e sottogruppi $G_1$ di $G_4$}{456}
\entry{126}{Sottogruppi $G_2$ a due parametri in $G_4$}{460}
\entry{127}{Sottogruppi $G_3$ di $G_4$ e riepilogo}{464}
\entry{128}{Criterio per distinguere i gruppi primitivi del piano dagli imprimitivi}{466}
\entry{129}{Ordine massimo delle trasformazioni infinitesime in un gruppo primitivo}{471}
\entry{130}{Numero dei parametri di un gruppo primitivo nel piano}{474}
\entry{131}{Riduzione dei gruppi primitivi a forma normale nel caso A)}{478}
\entry{132}{Riduzione a forma normale nei casi B) e C)}{482}
\entry{133}{I gruppi $G_8$, $G_6$, $G_5$ come gruppi proiettivi}{486}
\entry{134}{Gruppi intransitivi del piano. -- Casi in cui il gruppo $G_r$ è ad uno o a due parametri}{490}
\entry{135}{Caso in cui il gruppo ridotto $G'_r$ è a tre parametri}{495}
\entry{136}{Gruppi imprimitivi transitivi. -- Caso in cui il gruppo accorciato $\Gamma'$ ha un parametro}{500}
\entry{137}{Caso in cui $\Gamma'$ ha due parametri}{505}
\entry{138}{Caso in cui $\Gamma'$ ha tre parametri}{511}
\addlinespace[0.5em]
\multicolumn{3}{c}{\textsc{Capitolo XI}}\\
\entry{}{\textit{Prime applicazioni della teoria dei gruppi continui finiti ai problemi d'integrazione delle equazioni differenziali.}}{}
\entry{139}{Preliminari}{518}
\entry{140}{Equazioni differenziali del 1.º ordine che ammettono un gruppo $Xf$}{520}
\entry{141}{Relazioni fra le trasformazioni infinitesime ed i fattori integranti}{524}
\entry{142}{Determinazione delle trasformazioni infinitesime ammesse da un'equazione differenziale del 1.º ordine}{527}
\entry{143}{Interpretazione geometrica dei fattori integranti}{532}
\entry{144}{Caso di equazioni del 1.º ordine non risolute rapporto alla derivata}{536}
\entry{145}{Determinazione delle equazioni del 1.º ordine che ammettono un dato $G_1$}{538}
\entry{146}{Sistemi di equazioni differenziali e loro trasformazioni infinitesime}{542}
\entry{147}{Proprietà delle trasformazioni infinitesime ammesse da un sistema differenziale}{544}
\entry{148}{Integrazione di una $Af=0$ con note trasformazioni infinitesime}{549}
\entry{149}{Equazioni $Af=0$ con un gruppo $G_n$ integrabile}{553}
\entry{150}{Riduzione di un'equazione $Af=0$ con gruppo composto a successive equazioni con gruppi semplici}{556}
\entry{151}{I moltiplicatori di Jacobi per un'equazione $Af=0$}{563}
\entry{152}{Relazioni fra i moltiplicatori e le trasformazioni infinitesime}{568}
\entry{153}{Integrazione delle equazioni $Af=0$ a tre variabili con una trasformazione infinitesima nota}{572}
\entry{154}{Caso in cui della $Af=0$ si conoscono due trasformazioni infinitesime}{574}
\entry{155}{Criterio per riconoscere se un'equazione del 2.º ordine $y''=\omega(x,y,y')$ ammette un gruppo $Xf$}{580}
\entry{156}{Massimo numero di trasformazioni infinitesime ammesse da un'equazione del 2.º ordine}{584}
\entry{157}{Equazioni differenziali del 2.º ordine con un gruppo $G_3$}{589}
\entry{158}{Equazioni differenziali del 2.º ordine o d'ordine superiore con un dato $G_1$}{592}
\entry{159}{Integrazione di un'equazione differenziale del 2.º ordine con un noto $G_1$}{595}
\entry{160}{Riduzione a forma normale dei gruppi Abeliani $G_2$ su due variabili}{598}
\entry{161}{Caso di un $G_2$ non Abeliano}{602}
\entry{162}{Integrazione di un'equazione differenziale del 2.º ordine con un noto $G_2$}{605}
\addlinespace[0.5em]
\multicolumn{3}{c}{\textsc{Capitolo XII}}\\
\entry{}{\textit{Applicazioni alla teoria delle equazioni differenziali lineari. -- Estensione delle teorie di Galois secondo Picard e Vessiot.}}{}
\entry{163}{Criterio per la dipendenza lineare di più funzioni}{609}
\entry{164}{Sistemi fondamentali d'integrali di un'equazione differenziale lineare}{614}
\entry{165}{Invarianti differenziali razionali del gruppo lineare omogeneo totale $L$}{617}
\entry{166}{Invarianti dei sottogruppi di un gruppo $G_r$}{623}
\entry{167}{Risolventi delle equazioni differenziali lineari}{628}
\entry{168}{La risolvente di Picard e le sue proprietà fondamentali}{634}
\entry{169}{Campo di razionalità ed equazioni differenziali algebriche}{642}
\entry{170}{Gruppo di razionalità di un'equazione differenziale lineare}{648}
\entry{171}{Proprietà fondamentale del gruppo di razionalità}{654}
\entry{172}{Altra definizione del gruppo di razionalità. -- Gruppo dell'equazione generale}{658}
\entry{173}{Riduzione del gruppo di razionalità per l'aggiunta di una funzione razionale}{663}
\entry{174}{Riduzione del gruppo ad un gruppo puro $G_r$ di Lie}{668}
\entry{175}{Condizioni necessarie per l'integrabilità per quadrature}{672}
\entry{176}{Equazione differenziale lineare generale. -- Semplicità del gruppo $L_1$}{677}
\entry{177}{La seconda parte del teorema di Vessiot}{684}
\addlinespace[0.5em]
\entry{}{\textsc{Nota I.} Le trasformazioni finite di un gruppo $G_r=[X_1f,X_2f,\ldots,X_rf]$ composto con $r$ trasformazioni prese rispettivamente negli $r$ sottogruppi $[X_1f]$, $[X_2f]$, $\ldots$, $[X_rf]$}{688}
\entry{}{\textsc{Nota II.} Le equazioni differenziali ordinarie che posseggono sistemi di soluzioni fondamentali (Sistemi di Lie)}{692}
\bottomrule
\end{longtable}
\normalsize



\mainmatter
\chapter*{Capitolo I}
\addcontentsline{toc}{chapter}{Capitolo I}

\begin{center}
\begin{minipage}{0.88\linewidth}
\centering
Trasformazioni e loro composizione -- Serie di trasformazioni e parametri essenziali -- Concetto di gruppo continuo finito secondo Lie -- Equazioni differenziali fondamentali.\par
Gruppi ad un parametro -- Trasformazioni infinitesime -- Le $\infty^r$ trasformazioni infinitesime indipendenti per un gruppo a $r$ parametri.
\end{minipage}
\end{center}

\section*{\S{} 1. Trasformazioni e loro composizione}

Siano $x_1,x_2,\ldots,x_n$ $n$ variabili indipendenti, che in generale supporremo complesse. Chiameremo \emph{trasformazione} il passaggio da queste $n$ variabili ad $n$ nuove variabili indipendenti $x'_1,x'_2,\ldots,x'_n$, legate alle primitive da un sistema di formole:
\begin{equation*}\tag{1}
\left\{
\begin{aligned}
 x'_1&=f_1(x_1,x_2,\ldots,x_n),\\
 x'_2&=f_2(x_1,x_2,\ldots,x_n),\\
 &\hspace{1.6em}\cdots\cdots\cdots\\
 x'_n&=f_n(x_1,x_2,\ldots,x_n),
\end{aligned}
\right.
\end{equation*}
o come scriveremo più brevemente
\begin{equation*}\tag{$1^*$}
 x'_i=f_i(x_1,x_2,\ldots,x_n)\qquad (i=1,2,\ldots,n).
\end{equation*}
Rispetto alle funzioni $f_1,f_2,\ldots,f_n$ supporremo sempre, ove non si dica esplicitamente il contrario, che siano funzioni analitiche delle variabili $x$, definite dunque in un campo primitivo da una serie assolutamente convergente, procedente secondo le potenze intere e positive degli argomenti
\[
 x_1-x_1^{(0)},\quad x_2-x_2^{(0)},\quad \ldots,\quad x_n-x_n^{(0)},
\]
dove con $x_i^{(0)}$ indichiamo valori fissi delle $x$. Le funzioni $f_i$, definite così dapprima solo nell'intorno di un punto $(x_1^{(0)},\ldots,x_n^{(0)})$, si potranno poi prolungare analiticamente secondo i principî della teoria generale delle funzioni analitiche.

L'ipotesi che le $x'_i$ siano fra loro indipendenti porta che il determinante funzionale
\[
 \frac{\partial(f_1,f_2,\ldots,f_n)}{\partial(x_1,x_2,\ldots,x_n)}
\]
sia in generale diverso da zero. È ben noto allora che, restringendo convenientemente l'intorno del punto $x_i^{(0)}$, si possono risolvere le (1) rispetto alle $x$, esprimendo così le $x$ per le $x'$ colle formole della trasformazione inversa
\begin{equation*}\tag{2}
 x_i=F_i(x'_1,x'_2,\ldots,x'_n).
\end{equation*}
Come si è detto, si riterranno in generale le variabili come complesse. Per altro dallo sviluppo stesso delle ricerche risulterà che esse sono ancora applicabili a variabili e trasformazioni reali, quando si supponga che le funzioni $f_i$, da cui la trasformazione dipende, siano ad un sol valore, finite e continue in un certo campo insieme colle loro derivate prime, seconde ecc., fino all'ordine che conviene considerare.

Una trasformazione come la (1) si indicherà spesso con una sola lettera, o simbolo, $S$, nel qual caso la trasformazione inversa (2) si denoterà con $S^{-1}$. Ciò che vi ha di essenziale nella trasformazione $S$ rappresentata dalle (1) è naturalmente la specie delle funzioni $f_i$, e non già il modo come si indicano le variabili primitive e le trasformate; e così la trasformazione
\[
 y'_i=f_i(y_1,y_2,\ldots,y_n)
\]
sarà ancora la $S$ stessa.

Consideriamo una prima trasformazione
\[
 S)\qquad x'_i=f_i(x_1,x_2,\ldots,x_n)
\]
ed una seconda
\[
 T)\qquad x''_i=\varphi_i(x'_1,x'_2,\ldots,x'_n).
\]
Se supponiamo che le variabili $x$ si mantengano in un tal campo che i risultanti valori delle $x'$ appartengano al campo in cui sono definite le $\varphi_i$, potremo eseguire le due trasformazioni successivamente, e la trasformazione che ne risulta
\[
 x''_i=\varphi_i\{f_1(x_1,\ldots,x_n),f_2(x_1,\ldots,x_n),\ldots,f_n(x_1,\ldots,x_n)\},
\]
o, come spesso scriveremo per brevità,
\[
 x''_i=\varphi_i(f_1,f_2,\ldots,f_n),
\]
si dice la trasformazione composta o il prodotto delle due nell'ordine stabilito. Essa si indicherà col simbolo $ST$, ponendo a sinistra la trasformazione eseguita prima.

Si osservi che componendo una trasformazione $S$ colla sua inversa $S^{-1}$, il risultato è la sostituzione identica o l'identità
\[
 x'_i=x_i,
\]
che si indica col simbolo $1$, sicché
\[
 SS^{-1}=1.
\]
Colla definizione del prodotto di due trasformazioni risulta altresì stabilito il significato del prodotto di più trasformazioni $S_1,S_2,\ldots,S_m$, ritenendo
\[
 S_1S_2\cdots S_m=(S_1S_2\cdots S_{m-1})S_m.
\]
Si osservi subito che in un prodotto di più trasformazioni non vale in generale la legge permutativa, ma vale però quella associativa. Questa proprietà, che risulta immediatamente dal significato stesso delle trasformazioni,\footnote{Basta infatti dimostrare che sussiste la legge associativa elementare espressa dalla formula $(S_1S_2)S_3=S_1(S_2S_3)$. Ora se $S_1$ porta le $x_i$ nelle $x'_i$, la $S_2$ le $x'_i$ nelle $x''_i$, e infine la $S_3$ le $x''_i$ nelle $x'''_i$, l'effetto della trasformazione a sinistra è di portare le $x_i$ nelle $x'''_i$. D'altronde, siccome la $S_2S_3$ porta le $x'_i$ nelle $x'''_i$, così anche quella a destra trasforma le $x_i$ nelle $x'''_i$, e però le due trasformazioni sono identiche.} è importante perché permette di eseguire sopra un prodotto di trasformazioni quelle riduzioni che per la legge associativa sono permesse in un prodotto di ordinarie quantità, così per es.:
\[
 (S_1S_2)\cdot(S_3S_4)=S_1S_2S_3S_4,
\]
\[
 (S_1S)\cdot(S^{-1}S_2)=S_1S_2,\quad\hbox{etc.}
\]

Quando sopra le antiche variabili $x$ e le nuove $x'$ in una trasformazione
\[
 S)\qquad x'_i=f_i(x_1,x_2,\ldots,x_n)
\]
si eseguisca una medesima trasformazione
\[
T\left\{
\begin{aligned}
 y_i&=\varphi_i(x_1,x_2,\ldots,x_n),\\
 y'_i&=\varphi_i(x'_1,x'_2,\ldots,x'_n),
\end{aligned}
\right.
\]
la trasformazione che ne nasce pel passaggio dalle $y$ alle $y'$ è precisamente la
\[
 T^{-1}ST,
\]
che dicesi la trasformata della $S$ per mezzo della $T$. Invero colla $T^{-1}$ si passa dalle $y$ alle $x$, colla $S$ dalle $x$ alle $x'$, e infine colla $T$ dalle $x'$ alle $y'$; dunque colla $T^{-1}ST$ dalle $y$ alle $y'$.

\section*{\S{} 2. Serie $\infty^r$ di trasformazioni}

Supponiamo che nelle espressioni di una trasformazione $S$ entrino $r$ parametri $a_1,a_2,\ldots,a_r$, sicché si abbia
\begin{equation*}\tag{3}
 x'_i=f_i(x_1,x_2,\ldots,x_n; a_1,a_2,\ldots,a_r).
\end{equation*}
Questi parametri, come le variabili $x$, si supporranno in generale complessi, e le $f_i$ saranno funzioni analitiche anche dei parametri $a$.

Diremo che gli $r$ parametri $a$ sono essenziali nelle $f$, quando queste funzioni non possono ridursi a dipendere da un numero minore di parametri $b$ funzioni delle $a$.\footnote{Se per es. i due parametri $a_1,a_2$ compariscono nelle $f$ solo nella combinazione $a_1+a_2$, il numero dei parametri effettivi si abbasserebbe evidentemente almeno a $r-1$.}
Facendo percorrere ai parametri $a$ tutti i valori di un certo campo, la trasformazione $S$ percorrerà una serie continua di trasformazioni. Ma si noti che, ove i parametri $a_1,a_2,\ldots,a_r$ non siano tutti essenziali, ogni trasformazione della serie si otterrà infinite volte da diversi sistemi particolari di valori dei parametri, mentre quando sono essenziali, a quella trasformazione corrisponderà un solo sistema di valori dei parametri $a_1,a_2,\ldots,a_r$, ovvero un numero discreto di tali sistemi.

La prima asserzione è d'immediata evidenza, e la seconda risulta dimostrata osservando che negli sviluppi in serie che definiscono le $f_i$ come funzioni delle $x$ i coefficienti $\alpha$ sono certe funzioni delle $a$, e il numero $s$ di queste $\alpha$ indipendenti è certo $\leq r$. Quando $s<r$, introducendo i parametri $\alpha$ al posto dei parametri $a$, si vede che questi non sono essenziali; e inversamente, se i parametri $a$ non sono essenziali, sarà $s<r$, poiché le $a$ potranno scriversi come funzioni di quei parametri $b$ in numero minore di $r$ che possono sostituirsi alle $a$.

Quando i parametri $a$ sono essenziali, e dunque $s=r$, e le eguaglianze
\[
 f_i(x,a)=f_i(x,a')
\]
richiedono che si abbia
\[
\begin{aligned}
 \alpha_1(a)&=\alpha'_1(a'),\\
 \alpha_2(a)&=\alpha'_2(a'),\\
 &\cdots\cdots\cdots\\
 \alpha_r(a)&=\alpha'_r(a'),
\end{aligned}
\]
e quindi i valori ammissibili per $a'_1,\ldots,a'_r$ o si riducono all'unico sistema $a_1,a_2,\ldots,a_r$, o formano una serie discreta, c. d. d.

Per le circostanze ora osservate, allorquando nelle (3) i parametri $a_1,a_2,\ldots,a_r$ siano essenziali, si dirà che la serie (3) è una serie $\infty^r$ di trasformazioni. Propriamente quando i parametri $a$ siano complessi, la serie è $\infty^{2r}$, ma la locuzione introdotta ed ormai in uso in casi simili non può generare confusione quando si rammenti che, nel caso di parametri complessi, il grado effettivo d'infinità della serie è sempre il doppio di quello adottato.

È importante stabilire un criterio per riconoscere, data una serie (3) di trasformazioni, se gli $r$ parametri $a$ sono o no essenziali. Per questo incominciamo dall'osservare che, se $a_1,a_2,\ldots,a_r$ non sono essenziali, ma si possono ridurre agli $s<r$ parametri
\[
 b_1,b_2,\ldots,b_s,
\]
funzioni delle $a$, potremo determinare tali $r$ funzioni $A$ delle $a$ che le $n$ funzioni $f_i$, considerate come funzioni delle $a$, soddisfino tutte all'equazione lineare omogenea alle derivate parziali
\begin{equation*}\tag{$\alpha$}
 A_1\frac{\partial\psi}{\partial a_1}+A_2\frac{\partial\psi}{\partial a_2}+\cdots+A_r\frac{\partial\psi}{\partial a_r}=0.
\end{equation*}
Basta infatti determinare le $A$ in guisa che soddisfino alle $s<r$ equazioni lineari omogenee
\begin{equation*}\tag{$\beta$}
 A_1\frac{\partial b_i}{\partial a_1}+A_2\frac{\partial b_i}{\partial a_2}+\cdots+A_r\frac{\partial b_i}{\partial a_r}=0\qquad (i=1,2,\ldots,s),
\end{equation*}
e allora, poiché le $a$ entrano nelle $f_i$ solo nelle combinazioni $b_1,\ldots,b_s$, si vede subito che tutte le $f$ soddisferanno alla $(\alpha)$.

Si osservi di più che le $s$ equazioni $(\beta)$, lineari omogenee nelle $A$, hanno una caratteristica precisamente $=s$, poiché $b_1,b_2,\ldots,b_s$ si suppongono indipendenti (altrimenti il numero dei parametri essenziali sarebbe ancora $<s$); perciò $r-s$ delle $A$ restano arbitrarie, cioè in luogo della sola $(\alpha)$ potremo scrivere un sistema di $r-s$ tali equazioni indipendenti, che forniranno un sistema completo avendo le $r-(r-s)$ soluzioni $b$, e di questo sistema $(\alpha)$ saranno altrettante soluzioni le $f_1,f_2,\ldots,f_n$.

Inversamente se esiste una equazione $(\alpha)$, o un sistema di tali equazioni colle $A$ funzioni delle sole $a$, di cui le $f_i$ siano soluzioni, in queste funzioni i parametri $a_1,a_2,\ldots,a_r$ non saranno tutti essenziali, ma si ridurranno ad $s$ soltanto supposto che esistano $r-s$ equazioni indipendenti, come le $(\alpha)$, di cui le $f_i$ siano soluzioni. È infatti in tale ipotesi il sistema delle $r-s$ equazioni $(\alpha)$, supposto completo, ammetterà solo $s=r-(r-s)$ soluzioni indipendenti
\[
 b_1,b_2,\ldots,b_s,
\]
e le $f_i$ quindi, in quanto dipendono dalle $a$ ed essendo per ipotesi soluzioni del sistema $(\alpha)$, si ridurranno a funzioni delle sole $b$.

Concludiamo di qui. Condizione necessaria e sufficiente affinché nelle $n$ funzioni
\[
 f_i(x_1,x_2,\ldots,x_n; a_1,a_2,\ldots,a_r)
\]
gli $r$ parametri $a$ non siano essenziali è che esista una o più equazioni della forma
\begin{equation*}\tag{$\alpha$}
 A_1\frac{\partial f}{\partial a_1}+A_2\frac{\partial f}{\partial a_2}+\cdots+A_r\frac{\partial f}{\partial a_r}=0
\end{equation*}
colle $A$ funzioni delle $a$, di cui tutte le $f_i$ siano soluzioni. Precisamente se esistono $r-s$ equazioni $(\alpha)$ indipendenti, formanti un sistema completo di cui le $f_i$ siano soluzioni, i parametri essenziali nelle $f$ si riducono ad $s$ soltanto.

\section*{\S{} 3. Criterio per riconoscere il numero dei parametri essenziali}

Resta da trasformarsi il criterio teorico testé stabilito in un altro di pratica applicazione. Ma prima osserviamo che le considerazioni del \S{} precedente si applicano non solo al caso di un numero $n$ di funzioni $f_i$ eguale al numero delle variabili indipendenti $x_1,x_2,\ldots,x_n$, ma più in generale per riconoscere, dato un numero qualunque $m$ di funzioni
\[
 f_i(x_1,x_2,\ldots,x_n;a_1,a_2,\ldots,a_r)\qquad (i=1,2,\ldots,m)
\]
delle $n$ variabili indipendenti $x$ e di $r$ parametri $a$, a quanti parametri essenziali si possono ridurre i parametri $a$.

Tutto si riduce a trovare il numero $r-s$ delle equazioni lineari indipendenti
\[
 A_1\frac{\partial f}{\partial a_1}+A_2\frac{\partial f}{\partial a_2}+\cdots+A_r\frac{\partial f}{\partial a_r}=0,
\]
colle $A$ funzioni delle sole $a$, cui soddisfano identicamente tutte le funzioni $f_1,f_2,\ldots,f_m$. Dovendosi sussistere per ipotesi le $m$ equazioni
\begin{equation*}\tag{4}
 A_1\frac{\partial f_i}{\partial a_1}+A_2\frac{\partial f_i}{\partial a_2}+\cdots+A_r\frac{\partial f_i}{\partial a_r}=0\qquad (i=1,2,\ldots,m),
\end{equation*}
cominciamo dal costruire la matrice
\[
M=\left\|\begin{array}{cccc}
\dfrac{\partial f_1}{\partial a_1} & \dfrac{\partial f_2}{\partial a_1} & \cdots & \dfrac{\partial f_m}{\partial a_1}\\[0.65em]
\dfrac{\partial f_1}{\partial a_2} & \dfrac{\partial f_2}{\partial a_2} & \cdots & \dfrac{\partial f_m}{\partial a_2}\\
\vdots & \vdots & & \vdots\\
\dfrac{\partial f_1}{\partial a_r} & \dfrac{\partial f_2}{\partial a_r} & \cdots & \dfrac{\partial f_m}{\partial a_r}
\end{array}\right\|
\]
e sia $\mu$ la sua caratteristica. Si ha certamente $\mu\le r$; ma se $\mu=r$, allora nelle (4) si annulleranno forzatamente tutte le $A$, cioè gli $r$ parametri $a$ sono tutti essenziali e la questione è risoluta. Il numero sopra indicato con $s$ è allora $=r$.

Se poi le (4), derivando una prima volta rispetto a $x_1,x_2,\ldots,x_n$, si formano le nuove $nm$ equazioni lineari omogenee nelle $A$
\begin{equation*}\tag{$4_1$}
 A_1\frac{\partial^2 f_i}{\partial a_1\partial x_k}+A_2\frac{\partial^2 f_i}{\partial a_2\partial x_k}+\cdots+A_r\frac{\partial^2 f_i}{\partial a_r\partial x_k}=0,
 \qquad\begin{gathered} i=1,2,\ldots,m,\\ k=1,2,\ldots,n,\end{gathered}
\end{equation*}
e si costruisca la matrice completa delle equazioni lineari (4), $(4_1)$ nelle $A$, che indicheremo con $M_1$,
\[
M_1=\left\|\begin{array}{ccccccc}
\dfrac{\partial f_1}{\partial a_1}&\dfrac{\partial f_2}{\partial a_1}&\cdots&\dfrac{\partial f_m}{\partial a_1}&
\dfrac{\partial^2 f_1}{\partial a_1\partial x_1}&\cdots&\dfrac{\partial^2 f_m}{\partial a_1\partial x_n}\\[0.65em]
\dfrac{\partial f_1}{\partial a_2}&\dfrac{\partial f_2}{\partial a_2}&\cdots&\dfrac{\partial f_m}{\partial a_2}&
\dfrac{\partial^2 f_1}{\partial a_2\partial x_1}&\cdots&\dfrac{\partial^2 f_m}{\partial a_2\partial x_n}\\
\vdots&\vdots&&\vdots&\vdots&&\vdots\\
\dfrac{\partial f_1}{\partial a_r}&\dfrac{\partial f_2}{\partial a_r}&\cdots&\dfrac{\partial f_m}{\partial a_r}&
\dfrac{\partial^2 f_1}{\partial a_r\partial x_1}&\cdots&\dfrac{\partial^2 f_m}{\partial a_r\partial x_n}
\end{array}\right\|,
\]
e che si forma dalla $M$ aggiungendo alle $f_1,f_2,\ldots,f_m$, come nuove funzioni, tutte le loro derivate prime rapporto all'$x$. Se indichiamo con $\mu_1$ la caratteristica della matrice $M_1$, sarà evidentemente
\[
 \mu\le \mu_1\le r.
\]
Quando fosse $\mu_1=r$, le (4) colle $(4_1)$ mostrano che le $A$ sono tutte nulle e gli $r$ parametri $a$ sono tutti essenziali. Supposto ancora $\mu_1<r$, si derivino nuovamente le $(4_1)$ rispetto a tutte le $x$ e si formino le nuove $n^2m$ equazioni
\begin{equation*}\tag{$4_2$}
 A_1\frac{\partial^3 f_i}{\partial a_1\partial x_k\partial x_l}+A_2\frac{\partial^3 f_i}{\partial a_2\partial x_k\partial x_l}+\cdots+A_r\frac{\partial^3 f_i}{\partial a_r\partial x_k\partial x_l}=0.
\end{equation*}
E si indichi con $\mu_2$ la caratteristica della matrice $M_2$ formata aggiungendo alle colonne di $M_1$ a destra quelle relative alle equazioni $(4_2)$, cioè associando alle
\[
 f_i,\qquad \frac{\partial f_i}{\partial x_k},
\]
anche le derivate seconde $\dfrac{\partial^2 f_i}{\partial x_k\partial x_l}$. Sarà $\mu_1\le\mu_2\le r$ e, se $\mu_2=r$, gli $r$ parametri saranno essenziali. Così proseguendo costruiremo le successive matrici
\[
 M,
 M_1,
 M_2,
 \ldots
\]
di caratteristiche rispettive
\[
 \mu\le\mu_1\le\mu_2\le\mu_3\le\cdots .
\]
Poiché tutte le $\mu$ sono $\le r$, non potrà la serie
\[
 \mu,\mu_1,\mu_2,\ldots
\]
essere sempre crescente, e ad un certo punto sarà necessariamente una $\mu$ eguale alla precedente, diciamo
\[
 \mu_k=\mu_{k-1}.
\]
Allora si vede subito che tutte le successive $\mu$ conserveranno lo stesso valore $\mu_{k-1}$. E infatti, essendo $\mu_k=\mu_{k-1}$, le colonne aggiunte alla matrice $M_{k-1}$ per formare la matrice $M_k$ sono combinazioni lineari omogenee di $\mu_{k-1}$ colonne fisse in $M_{k-1}$; queste relazioni, derivate rispetto alle $x$, dimostrano che le nuove colonne in $M_{k+1}$ sono pure combinazioni lineari omogenee delle dette $\mu_{k-1}$ colonne di $M_{k-1}$; così dicasi per $M_{k+2},M_{k+3},\ldots$, come risulterà successivamente derivando. Dunque la serie di numeri
\[
 \mu,\mu_1,\mu_2,\ldots,\mu_{k-1},\ldots
\]
è sempre crescente fino ad un massimo valore $\mu_{k-1}$, dopo il quale tutte le $\mu$ eguagliano $\mu_{k-1}$, onde basterà calcolare le $\mu$ fino a $\mu_{k-1}$.

Se diciamo $s=\mu_{k-1}$ questo massimo valore delle $\mu$, sarà certamente $s\le r$, talché in ogni caso, essendo almeno $\mu=1$, e la serie di numeri
\[
 \mu,\mu_1,\mu_2,\ldots,\mu_{k-1}
\]
crescendo ad ogni passo almeno di un'unità, sarà anche $k\le r$, sicché colla costruzione di $r$ al più delle nostre matrici si arriverà al termine del calcolo.

Dimostriamo ora che questo numero $s$, che dà il massimo valore delle caratteristiche delle matrici
\[
 M,M_1,M_2,\ldots,
\]
eguaglia precisamente il numero dei parametri essenziali nelle $m$ funzioni $f_i(x,a)$. Perciò osserviamo che le equazioni
\[
 (4),(4_1),(4_2),\ldots,(4_{k-1})
\]
si riducono precisamente ad $s$ equazioni indipendenti nelle $A$, mentre delle successive $(4_k),(4_{k+1}),\ldots$ è inutile tener conto, perché sono conseguenza delle $s$ già scritte.

Dunque, se $s=r$, sono necessariamente nulle tutte le $A$ e, non esistendo alcuna equazione
\[
 A_1\frac{\partial f}{\partial a_1}+\cdots+A_r\frac{\partial f}{\partial a_r}=0,
\]
di cui le $f_i$ siano soluzioni, gli $r$ parametri $a$ sono in effetto essenziali.

Se invece $s<r$, è certo possibile soddisfare le dette equazioni
\[
 (4),(4_1),(4_2),\ldots,(4_{k-1})
\]
con dei valori delle $A$ in generale funzioni non solo delle $a$, ma anche delle $x$ stesse. E in ciò si possono prendere ad arbitrio $r-s$ delle $A$. Ora facciamo vedere che, se si prendono per queste $r-s$ $A$ arbitrarie delle funzioni dei parametri $a$ soltanto, anche le rimanenti $s$ $A$ conterranno le sole $a$. Suppongasi, per fissare le idee, che le $r-s$ $A$ arbitrarie siano $A_1,A_2,\ldots,A_{r-s}$ e le dette equazioni diano poi
\[
 A_{r-s+1},\ldots,A_r
\]
come funzioni lineari omogenee delle prime; diciamo p. e.
\begin{equation*}\tag{5}
 A_{r-s+i}=\lambda_{i1}A_1+\lambda_{i2}A_2+\cdots+\lambda_{i,r-s}A_{r-s},
\end{equation*}
dove le $\lambda$ sono funzioni note delle $x$ e delle $a$. Se deriviamo le
\[
 (4),(4_1),\ldots,(4_{k-1})
\]
rispetto ad una qualunque $x$, sia $x_\rho$, ed osserviamo che le $A$ soddisfano anche le equazioni $(4_1)$, ne deduciamo che le dette equazioni sono ancora soddisfatte sostituendo ad $A_1,A_2,\ldots,A_r$ le loro derivate
\[
 \frac{\partial A_1}{\partial x_\rho},\quad
 \frac{\partial A_2}{\partial x_\rho},\quad\ldots,\quad
 \frac{\partial A_r}{\partial x_\rho}
\]
rispetto ad una medesima $x$. Perciò, insieme alla (5), sussiste l'altra
\[
 \frac{\partial A_{r-s+i}}{\partial x_\rho}=\lambda_{i1}\frac{\partial A_1}{\partial x_\rho}+\lambda_{i2}\frac{\partial A_2}{\partial x_\rho}+\cdots+\lambda_{i,r-s}\frac{\partial A_{r-s}}{\partial x_\rho};
\]
e ne risulta appunto che prendendo $A_1,A_2,\ldots,A_{r-s}$ funzioni dei soli parametri $a$, anche $A_{r-s+1},\ldots,A_r$ conterranno questi soli parametri. Manifestamente otteniamo così precisamente $r-s$ equazioni indipendenti del tipo
\[
 A_1\frac{\partial f}{\partial a_1}+A_2\frac{\partial f}{\partial a_2}+\cdots+A_r\frac{\partial f}{\partial a_r}=0,\tag{*}
\]
colle $A$ funzioni dei soli parametri $a$, di cui tutte le $f_i(x;a)$ sono soluzioni, e quindi (\S{}2) il numero dei parametri essenziali si abbassa appunto ad $s$.\footnote{Per avere queste $r-s$ equazioni indipendenti basterà fare per es. successivamente $A_1=1,A_2=0,A_3=0,\ldots,A_{r-s}=0$; $A_1=0,A_2=1,A_3=0,\ldots,A_{r-s}=0$; e calcolare le rimanenti $A$ ogni volta dalle (5).}

Riepiloghiamo i risultati ottenuti nella regola seguente. Per vedere quanti sono i parametri essenziali $a$ in $m$ date funzioni
\[
 f_i(x_1,x_2,\ldots,x_n;a_1,a_2,\ldots,a_r)\qquad (i=1,2,\ldots,m)
\]
delle $x$ e dei parametri $a$, si cominci dal costruire la matrice funzionale $M$ delle derivate delle $f$ rapporto alle $a$; poi la matrice analoga $M_1$, aggiungendo alle $f_i$ tutte le loro prime derivate rapporto alle $x$; indi la matrice $M_2$, ottenuta aggiungendo ancora le derivate seconde, e così di seguito. Si prolunghi la costruzione di queste matrici $M,M_1,M_2,\ldots$, finché la serie $\mu,\mu_1,\mu_2,\ldots$ delle loro caratteristiche cessi di crescere, dopo di che le $\mu$ successive conserveranno un valore costante $s\le r$. Questo numero $s$ dà precisamente il numero dei parametri essenziali.

\section*{\S{} 4. Definizione dei gruppi continui finiti di trasformazioni}

Le trasformazioni, essendo operazioni suscettibili di composizione, sono loro applicabili i concetti generali dei gruppi di operazioni. Diremo che una serie finita o infinita di trasformazioni costituisce un gruppo quando il prodotto di due trasformazioni qualunque della serie (distinte od anche coincidenti) è nuovamente contenuto nella serie stessa.

Quei gruppi che constano di un numero discreto, finito o infinito di trasformazioni diconsi discontinui. Tali sono, p. e., i gruppi di sostituzioni sopra lettere, la cui teoria è diventata, per opera specialmente di Galois, di fondamentale importanza per lo studio delle equazioni algebriche. Altri gruppi discontinui, ma infiniti, s'incontrano in importanti rami di analisi (funzioni ellittiche modulari, funzioni automorfe).

I gruppi di cui ci occuperemo nel presente corso sono invece i gruppi continui, nelle cui trasformazioni entrano dei parametri arbitrarii, suscettibili di assumere con continuità tutti i valori di un certo campo. La teoria di questi gruppi continui, quale oggi la possediamo, è stata costituita quasi per intero da Sophus Lie; essa ha acquistato nella scienza matematica una grande importanza tanto nel campo analitico come nel campo geometrico, in particolare per l'analisi nella teoria della integrazione delle equazioni differenziali ordinarie e a derivate parziali, dove l'ufficio della teoria dei gruppi continui è comparabile a quello che la teoria dei gruppi di sostituzioni compie nello studio delle equazioni algebriche.

I gruppi continui diconsi finiti, secondo Lie, quando le loro trasformazioni dipendono da un numero finito di parametri arbitrarii. I gruppi continui finiti più semplici sono quelli le cui trasformazioni sono definite da un unico sistema di equazioni
\begin{equation*}\tag{6}
 x'_i=f_i(x_1,x_2,\ldots,x_n;a_1,a_2,\ldots,a_r)\qquad (i=1,2,\ldots,n),
\end{equation*}
con $r$ parametri che possono supporsi senz'altro essenziali. La proprietà della serie $\infty^r$ di trasformazioni (6) di costituire un gruppo si traduce nella seguente proprietà funzionale per le $f_i$. Poiché due successive trasformazioni della serie (6)
\[
 x'_i=f_i(x_1,x_2,\ldots,x_n;a_1,a_2,\ldots,a_r),
\]
\[
 x''_i=f_i(x'_1,x'_2,\ldots,x'_n;b_1,b_2,\ldots,b_r)
\]
debbono comporsi in una terza trasformazione
\[
 x''_i=f_i(x_1,x_2,\ldots,x_n;c_1,c_2,\ldots,c_r),
\]
i parametri $c$ della trasformazione composta, dipendendo unicamente da quelli $a,b$ delle componenti, le funzioni $f_i$ debbono soddisfare alle relazioni funzionali
\[
 f_i\{f_1(x,a),f_2(x,a),\ldots,f_n(x,a);b_1,b_2,\ldots,b_r\}=f_i(x_1,x_2,\ldots,x_n;c_1,c_2,\ldots,c_r),
\]
essendo le $c$ funzioni delle $a,b$, diciamo
\[
 c_i=\varphi_i(a_1,a_2,\ldots,a_r;b_1,b_2,\ldots,b_r).
\]

In generale, quando parleremo di un gruppo di trasformazioni, intenderemo che si tratti di un gruppo continuo finito, le cui trasformazioni sono contenute in un'unica serie (6). Gli altri gruppi finiti più generali che conviene considerare sono quelli nei quali le trasformazioni del gruppo, in luogo di costituire una sola serie (6), costituiscono più di tali serie ed anche un numero infinito ma discreto. Per altro lo studio di tali gruppi finiti più generali, che si dicono gruppi misti, si riporta, come dimostreremo nel prossimo capitolo, a quello dei gruppi semplici (6) combinati con ordinarii gruppi discontinui.

Ogni gruppo continuo non finito potrebbe dirsi un gruppo continuo infinito. Per altro questa denominazione è adottata da Lie solo nel caso, sul quale più tardi ritorneremo (v. \S{}36), che le trasformazioni del gruppo stesso risultino definite da un certo sistema di equazioni a derivate parziali. Del resto non ci occuperemo affatto in questo corso dei gruppi infiniti nel senso di Lie o in senso più generale, e sottintenderemo sempre che il gruppo sia finito.

Un'osservazione è subito necessaria per evitare un facile errore. Nei gruppi ordinarii di un numero finito di operazioni esiste sempre nel gruppo l'operazione identica, come pure insieme ad ogni operazione la sua inversa. Ma per i gruppi continui facciamo l'osservazione (che vale pure per i gruppi discontinui infiniti): dalla definizione di gruppo continuo non segue affatto che il gruppo debba sempre contenere la trasformazione identica, e tanto meno che insieme ad ogni trasformazione del gruppo debba contenersi la sua inversa.

Più avanti dimostreremo soltanto che, se il gruppo contiene la trasformazione identica, corrispondente ai valori
\[
 a_1^{(0)},a_2^{(0)},\ldots,a_r^{(0)}
\]
dei parametri, e se questi valori appartengono all'interno del campo ove le $f_i$ sono definite, allora restando le $a$ in un intorno sufficientemente ristretto di $(a^{(0)})$, insieme ad ogni trasformazione $S_a$ nel gruppo esisterà effettivamente la sua inversa $S_a^{-1}$.

\section*{\S{} 5. Esempi di gruppi continui finiti}

Per illustrare i concetti generali posti nel \S{} precedente addurremo alcuni semplici esempi di gruppi continui tolti dalla geometria, ove si presentano nel modo più intuitivo.

$1^{\circ}$. Se una figura si muove nello spazio senza cangiar di forma, le coordinate $x',y',z'$ di ogni punto della figura nella nuova posizione si esprimono per funzioni lineari intere delle coordinate primitive $x,y,z$. Supposti assi ortogonali, le formole si scrivono:
\[
\left\{
\begin{aligned}
 x'&=a_1x+a_2y+a_3z+a,\\
 y'&=b_1x+b_2y+b_3z+b,\\
 z'&=c_1x+c_2y+c_3z+c,
\end{aligned}\right.\tag{a}
\]
dove le nove costanti
\[
\begin{vmatrix}
 a_1&a_2&a_3\\
 b_1&b_2&b_3\\
 c_1&c_2&c_3
\end{vmatrix}
\]
sono i coefficienti di una sostituzione ortogonale destrorsa e si esprimono quindi per tre soli indipendenti, per es. gli angoli $\theta,\varphi,\psi$ euleriani. La totalità dei movimenti rigidi dello spazio, rappresentata dalle formole (a), dove
\[
\begin{vmatrix}
 a_1&a_2&a_3\\
 b_1&b_2&b_3\\
 c_1&c_2&c_3
\end{vmatrix}
\]
percorre tutte le sostituzioni ortogonali destrorse ed $a,b,c$ assumono tutti i valori possibili, costituisce quindi un gruppo continuo finito con 6 parametri essenziali. È chiaro che in questo gruppo esiste la trasformazione identica ed insieme ad ogni trasformazione la sua inversa.

$2^{\circ}$. Se a tutti i movimenti (a) associamo la simmetria rispetto, per es., al piano $xy$ data dalle formole
\[
 x'=x,
 \qquad y'=y,
 \qquad z'=-z,
\]
otteniamo una seconda serie di trasformazioni che differisce dalla (a) solo in questo, che le corrispondenti sostituzioni ortogonali
\[
\begin{vmatrix}
 a_1&a_2&a_3\\
 b_1&b_2&b_3\\
 c_1&c_2&c_3
\end{vmatrix}
\]
sono sinistrorse. Le due serie insieme ci forniscono un semplice esempio di gruppo misto.

$3^{\circ}$. I movimenti delle figure piane nel proprio piano formano un gruppo continuo a tre parametri; insieme ai ribaltamenti ci danno un gruppo misto. Il primo è rappresentato analiticamente dalle formole
\[
\left\{
\begin{aligned}
 x'&=x\cos\alpha-y\sen\alpha+a,\\
 y'&=x\sen\alpha+y\cos\alpha+b,
\end{aligned}\right.
\]
il secondo associando a queste le altre
\[
\left\{
\begin{aligned}
 x'&=x\cos\alpha+y\sen\alpha+a,\\
 y'&=x\sen\alpha-y\cos\alpha+b.
\end{aligned}\right.
\]

$4^{\circ}$. Se al gruppo totale dei movimenti dello spazio associamo una similitudine
\[
 x'=kx,
 \qquad y'=ky,
 \qquad z'=kz
\]
a parametro $k$ variabile, otteniamo un gruppo continuo a 7 parametri, il gruppo delle similitudini.

$5^{\circ}$. Combinando i movimenti dello spazio con tutte le inversioni per raggi vettori reciproci rispetto alla sfera dello spazio, si ottiene un gruppo continuo con 10 parametri, che è il gruppo delle trasformazioni conformi. Secondo un teorema di Liouville, esso contiene tutte e sole le trasformazioni dello spazio che conservano gli angoli.

$6^{\circ}$. La totalità delle collineazioni dello spazio forma un gruppo continuo con 15 parametri, il cosiddetto gruppo proiettivo dello $S_3$. In esso formano un sottogruppo a 4 parametri quelle collineazioni che lasciano fissa una conica non degenere, un sottogruppo a 6 parametri quelle che lasciano fissa una quadrica non degenere, infine un sottogruppo a 3 parametri quelle che riportano in sé una cubica gobba.

Così ancora quelle collineazioni dello spazio che lasciano fisso un complesso lineare non speciale formano un gruppo a 10 parametri. In tutti gli esempi sopra addotti le formole di trasformazione sono date da sostituzioni lineari, salvo che nell'esempio $5^{\circ}$, ove abbiamo sostituzioni quadratiche.\footnote{È questo un particolare esempio dei gruppi Cremoniani, ove le formole di trasformazione danno le nuove variabili espresse razionalmente per le antiche ed inversamente.}

La generalizzazione allo spazio a $n$ dimensioni è immediata. Così le sostituzioni lineari omogenee sopra $n$ variabili
\begin{equation*}\tag{$\beta$}
 x'_i=\sum\limits_{k=1}^n a_{ik}x_k,
\end{equation*}
cogli $n^2$ parametri $a_{ik}$, ci danno evidentemente un gruppo di Lie. Componendo la precedente colla
\[
 x''_i=\sum\limits_{k=1}^n b_{ik}x'_k
\]
si ha la sostituzione composta
\[
 x''_i=\sum\limits_{k=1}^n c_{ik}x_k,
\]
dove le $c$ si esprimono per le $a,b$ colle formole
\[
 c_{hk}=\sum\limits_i^{1,\ldots,n} a_{hi}b_{ik}.
\]
Se nella $(\beta)$ poniamo fra i parametri $a$ la relazione espressa da
\[
 a=\begin{vmatrix}
 a_{11}&a_{12}&\cdots&a_{1n}\\
 \cdots&\cdots&\cdots&\cdots\\
 a_{n1}&a_{n2}&\cdots&a_{nn}
\end{vmatrix}=1,
\]
avremo evidentemente un sottogruppo con $n^2-1$ parametri. Se invece facciamo percorrere al modulo $a$ della sostituzione $(\beta)$ solo le $\nu$ radici $\nu$-me dell'unità, $\nu$ essendo un intero fisso, avremo evidentemente un gruppo misto con $\nu$ serie di trasformazioni a $n^2-1$ parametri. Così pure formano un gruppo ad $\dfrac{n(n-1)}{2}$ parametri le sostituzioni $(\beta)$ ortogonali destrorse.

\section*{\S{} 6. Risolubilità delle equazioni $c_i=\varphi_i(a,b)$ rispetto alle $a$ o alle $b$}

Supponiamo di avere una serie $\infty^r$ di trasformazioni
\begin{equation*}\tag{6}
 x'_i=f_i(x_1,x_2,\ldots,x_n;a_1,a_2,\ldots,a_r)\qquad (i=1,2,\ldots,n),
\end{equation*}
che formano un gruppo. Dalle equazioni funzionali
\begin{equation*}\tag{7}
 f_i(x',b)=f_i(x,c),
\end{equation*}
le quali risultano identicamente soddisfatte quando per le $x'$ vi si pongano i loro valori in funzione delle $x,a$
\begin{equation*}\tag{$7^*$}
 x'_i=f_i(x,a)
\end{equation*}
e per le $c$ i loro valori in funzione delle $a,b$
\begin{equation*}\tag{8}
 c_i=\varphi_i(a,b),
\end{equation*}
si possono trarre alcune conseguenze preliminari importanti.

Dimostriamo in primo luogo che le (8) sono risolubili rispetto alle $b$, cioè il determinante funzionale
\[
 \frac{\partial(\varphi_1,\varphi_2,\ldots,\varphi_r)}{\partial(b_1,b_2,\ldots,b_r)}
\]
non può essere nullo. In caso contrario infatti si potrebbero determinare tali funzioni $\alpha_1,\alpha_2,\ldots,\alpha_r$ delle $a,b$ da soddisfare il sistema delle $r$ equazioni lineari omogenee
\begin{equation*}\tag{9}
 \sum\limits_{l=1}^r \alpha_l\frac{\partial \varphi_i}{\partial b_l}=0\qquad (i=1,2,\ldots,r).
\end{equation*}
Ora derivando l'identità (7), ove si immaginino eseguite le sostituzioni $(7^*),(8)$, rispetto alla $b_k$, otteniamo
\[
 \frac{\partial f_i(x',b)}{\partial b_k}
 =\sum\limits_{j=1}^r\frac{\partial f_i(x,c)}{\partial c_j}\frac{\partial\varphi_j}{\partial b_k};
\]
onde, moltiplicando per $\alpha_k$ e sommando da $k=1$ a $k=r$, ne dedurremo per le (9)
\[
 \sum\limits_k^{1,\ldots,r} \alpha_k\frac{\partial f_i(x',b)}{\partial b_k}=0.
\]
Se in questa identità si ponessero per le $a$ i valori di un sistema particolare qualunque che non annullino tutte le $\alpha$, queste risulterebbero delle funzioni $B_k$ delle sole $b$. Le funzioni
\[
 f_i(x',b)
\]
sarebbero quindi soluzioni di un'equazione del tipo
\[
 B_1\frac{\partial f}{\partial b_1}+B_2\frac{\partial f}{\partial b_2}+\cdots+B_r\frac{\partial f}{\partial b_r}=0,
\]
onde si concluderebbe (\S{}2) che nelle funzioni $f_i(x';b)$ gli $r$ parametri $b$ non sarebbero tutti essenziali, contro il supposto.

Per dimostrare che le relazioni (8) sono altresì risolubili rispetto alle $a$, basta osservare che sussiste il teorema seguente. Se le $\infty^r$ trasformazioni (6) costituiscono un gruppo, anche le loro trasformazioni inverse
\[
 x_i=F_i(x'_1,\ldots,x'_n;a_1,\ldots,a_r)
\]
costituiscono un gruppo a $r$ parametri. Se indichiamo infatti simbolicamente con $S_a$ la trasformazione (6) corrispondente ai valori $a$ dei parametri, poiché le (6) formano un gruppo,
\[
 S_aS_b=S_c,
\]
e quindi
\[
 S_b^{-1}S_a^{-1}=S_c^{-1},
\]
ciò che dimostra che le trasformazioni inverse formano pure un gruppo. Il numero dei parametri nel gruppo inverso è poi ancora $=r$, poiché nel passaggio da un gruppo al suo inverso il numero dei parametri non può evidentemente aumentare e quindi nemmeno diminuire.\footnote{Analiticamente si può stabilire questo risultato servendosi delle identità
\[
 f_i(F_1(x,a),\ldots,F_n(x,a);a_1,a_2,\ldots,a_r)=x_i,
\]
che, derivate rapporto ad $a_k$, danno
\[
 \frac{\partial F_i}{\partial a_k}+\sum\limits_{h=1}^n\frac{\partial f_i}{\partial x_h}\frac{\partial F_h}{\partial a_k}=0.
\]
Se $a_1,a_2,\ldots,a_r$ non fossero tutti essenziali nelle $F$, queste soddisferebbero tutte ad un'equazione della forma
\[
 A_1\frac{\partial F}{\partial a_1}+\cdots+A_r\frac{\partial F}{\partial a_r}=0
\]
colle $A$ funzioni delle $a$. Le identità precedenti ci dimostrano che vi soddisferebbero anche le $f_i$, nelle quali i parametri $a$ non sarebbero dunque tutti essenziali.}
Dopo ciò, siccome le relazioni (8) sono le stesse che pel gruppo primitivo corrispondono all'eguaglianza $S_aS_b=S_c$ e, per l'inverso, alla $S_b^{-1}S_a^{-1}=S_c^{-1}$, è chiaro che le (8) sono altresì risolubili rispetto alle $a$, c. d. d.

Abbiamo dunque stabilito il risultato. Le equazioni (8)
\[
 c_i=\varphi_i(a,b),
\]
che per una serie $\infty^r$ di trasformazioni costituenti gruppo legano i parametri $a,b$ delle trasformazioni componenti $S_a,S_b$ a quelli $c$ della composta $S_c=S_aS_b$, sono risolubili tanto rispetto alle $a$ che alle $b$, cioè i determinanti funzionali
\[
 \frac{\partial(\varphi_1,\varphi_2,\ldots,\varphi_r)}{\partial(a_1,a_2,\ldots,a_r)},
 \qquad
 \frac{\partial(\varphi_1,\varphi_2,\ldots,\varphi_r)}{\partial(b_1,b_2,\ldots,b_r)}
\]
sono ciascuno diverso da zero.


\section*{\S{} 7. Equazioni differenziali fondamentali}

Procediamo ora ad una ricerca fondamentale per la teoria dei gruppi continui finiti, a dimostrare cioè come le equazioni funzionali (7) si traducano per le funzioni $f_i$ di trasformazione che figurano in una serie $\infty^r$ (6) costituente un gruppo, in un sistema di equazioni ai differenziali totali, illimitatamente integrabile, cui debbono soddisfare le $f_i$ stesse, considerate come funzioni dei parametri $a$.

Le equazioni funzionali
\begin{equation*}\tag{7}
 f_i(x',b)=f_i(x,c)
\end{equation*}
debbono tradursi, come sappiamo, in identità rispetto alle $x,a,b$ riguardate come indipendenti, quando vi si pongano per le $x_i'$, $c_i$ i valori
\begin{equation*}\tag{6}
 x_i'=f_i(x,a),
\end{equation*}
\begin{equation*}\tag{8}
 c_\alpha=\varphi_\alpha(a,b).
\end{equation*}
Siccome le (8) sono risolubili rispetto alle $b$ (\S{}6) possiamo invece riguardare le (7) come identità nelle $x,a,c$, ove per le $x'$ si pongano i valori (6) e per le $b$ i valori in funzione delle $a,c$ tratti dalle (8). Derivando la (7) rispetto ad una qualunque delle $a$ in queste connessione deduciamo
\begin{equation*}\tag{10}
 \sum\limits_l^{1,\ldots,n}\frac{\partial f_i}{\partial x_l'}\frac{\partial x_l'}{\partial a_k}
 +\sum\limits_s^{1,\ldots,r}\frac{\partial f_i}{\partial b_s}\frac{\partial b_s}{\partial a_k}=0
 \qquad (i=1,2,\ldots,n).
\end{equation*}
Poichè il determinante funzionale
\[
 J=\frac{\partial(f_1,f_2,\ldots,f_n)}{\partial(x'_1,x'_2,\ldots,x'_n)}
\]
non è zero, possiamo risolvere le $n$ equazioni (10) lineari in
\[
 \frac{\partial x'_1}{\partial a_k},\quad
 \frac{\partial x'_2}{\partial a_k},\quad\ldots,\quad
 \frac{\partial x'_n}{\partial a_k}
\]
rispetto a queste derivate. Se indichiamo con $\Theta_{ij}$ il quoziente del complemento algebrico di $\dfrac{\partial f_j}{\partial x'_i}$ nel determinante $J$ per il determinante stesso, la risoluzione delle (10) ci dà
\[
 \frac{\partial x'_i}{\partial a_k}
 =-\sum\limits_s^{1,\ldots,r}\sum\limits_j^{1,\ldots,n}
 \Theta_{ij}\frac{\partial f_j}{\partial b_s}\frac{\partial b_s}{\partial a_k}.
\]
Ora la somma
\[
 -\sum\limits_j^{1,\ldots,n}\Theta_{ij}\frac{\partial f_j(x',b)}{\partial b_s}
\]
è una certa funzione delle $x',b$, che indichiamo con $\Phi_{si}(x',b)$, mentre la derivata $\dfrac{\partial b_s}{\partial a_k}$ è una funzione delle $a,b$, che si può calcolare dalle (8), e che indichiamo con
\[
 \frac{\partial b_s}{\partial a_k}=\psi_{sk}(a,b).
\]
Con queste notazioni le formole superiori diventano
\begin{equation*}\tag{11}
 \frac{\partial x'_i}{\partial a_k}
 =\sum\limits_s^{1,\ldots,r}\psi_{sk}(a,b)\Phi_{si}(x',b).
\end{equation*}
Queste, come le (7) da cui sono tratte per derivazione, sono identicamente soddisfatte quando vi si ponga
\[
 x'_i=f_i(x,a),
\]
poichè le $x,a,b$ sono affatto indipendenti. Nei primi membri delle (11)
\[
 \frac{\partial f_i(x,a)}{\partial a_k}
\]
non compariscono le $b$, onde si vede che anche nei secondi possono essere solo apparenti. E se diamo quindi alle $b$ un sistema fisso di valori, del resto arbitrarî,
\[
 b_i=b_i^{(0)},
\]
e poniamo
\[
 \Phi_{sj}\bigl(x',b^{(0)}\bigr)=\xi_{sj}(x'_1,x'_2,\ldots,x'_n),
\]
\[
 \psi_{sk}\bigl(a,b^{(0)}\bigr)=\psi_{sk}(a_1,a_2,\ldots,a_r),
\]
le (11) ci danno il cercato sistema di equazioni fondamentali:
\begin{equation*}\tag{A}
 \frac{\partial x'_i}{\partial a_k}
 =\sum\limits_s^{1,\ldots,r}\xi_{si}(x')\psi_{sk}(a)
 \qquad
 \begin{cases}
 i=1,2,\ldots,n,\\
 k=1,2,\ldots,r,
 \end{cases}
\end{equation*}
cui soddisfano le funzioni
\[
 x'_i=f_i(x_1,x_2,\ldots,x_n;a_1,a_2,\ldots,a_r)
 \qquad i=1,2,\ldots,n,
\]
che definiscono il nostro gruppo.

Le nostre equazioni (A) possono anche scriversi sotto la forma equivalente
\begin{equation*}\tag{$A^*$}
 dx'_j=\sum\limits_k^{1,\ldots,r}\left(\sum\limits_s^{1,\ldots,r}\xi_{sj}(x')\psi_{sk}(a)\right)da_k
 \qquad (j=1,2,\ldots,n)
\end{equation*}
e costituiscono un sistema di equazioni ai differenziali totali. Negli integrali
\begin{equation*}\tag{$\gamma$}
 x'_j=f_j(x,a)
\end{equation*}
le $n$ $x$ figurano come costanti arbitrarie; e poichè le $(\gamma)$ sono risolubili rispetto alle $x$, ne segue che il nostro sistema di equazioni ai differenziali totali (A) o $(A^*)$ è illimitatamente integrabile.

Riepiloghiamo i risultati ottenuti fin qui nel teorema:

Se le equazioni
\[
 x'_i=f_i(x_1,\ldots,x_n; a_1,a_2,\ldots,a_r)
 \qquad i=1,2,\ldots,n
\]
rappresentano un gruppo con $r$ parametri $a$ essenziali, le funzioni $f_i(x,a)$, considerate come dipendenti dalle $a$, soddisfano ad un sistema di equazioni ai differenziali totali della forma
\begin{equation*}\tag{A}
 \frac{\partial x'_i}{\partial a_k}
 =\sum\limits_s^{1,\ldots,r}\xi_{si}(x')\psi_{sk}(a)
 =\xi_{1i}(x')\psi_{1k}(a)+\xi_{2i}(x')\psi_{2k}(a)+\cdots+\xi_{ri}(x')\psi_{rk}(a),
\end{equation*}
\footnote{Si osservi la forma particolare dei secondi membri di queste equazioni come somme di $r$ prodotti di funzioni delle sole $a$ per funzioni delle sole $x'$.}
questo sistema è illimitatamente integrabile e le $f_i(x,a)$ ne danno gli integrali generali.

Come esempio si prenda il gruppo dei movimenti nel piano (\S{}5), rappresentato dalle equazioni
\[
\begin{cases}
 x'_1=x_1\cos a_3-x_2\sen a_3+a_1,\\
 x'_2=x_1\sen a_3+x_2\cos a_3+a_2,
\end{cases}
\]
dove $x_1,x_2$ rappresentano coordinate cartesiane ortogonali e $a_1,a_2,a_3$ i parametri. Le equazioni differenziali sono qui
\[
\begin{array}{lll}
 \dfrac{\partial x'_1}{\partial a_1}=1, &
 \dfrac{\partial x'_1}{\partial a_2}=0, &
 \dfrac{\partial x'_1}{\partial a_3}=a_2-x'_2,\\[0.9em]
 \dfrac{\partial x'_2}{\partial a_1}=0, &
 \dfrac{\partial x'_2}{\partial a_2}=1, &
 \dfrac{\partial x'_2}{\partial a_3}=x'_1-a_1,
\end{array}
\]
e rientrano nella forma generale (A) ponendo
\[
 \left\|\begin{array}{cc}
 \xi_{11}&\xi_{12}\\
 \xi_{21}&\xi_{22}\\
 \xi_{31}&\xi_{32}
 \end{array}\right\|
 =
 \left\|\begin{array}{cc}
 1&0\\
 0&1\\
 x'_2&-x'_1
 \end{array}\right\|,
 \qquad
 \left\|\begin{array}{ccc}
 \psi_{11}&\psi_{12}&\psi_{13}\\
 \psi_{21}&\psi_{22}&\psi_{23}\\
 \psi_{31}&\psi_{32}&\psi_{33}
 \end{array}\right\|
 =
 \left\|\begin{array}{ccc}
 1&0&a_2\\
 0&1&-a_1\\
 0&0&-1
 \end{array}\right\|.
\]


\section*{\S{} 8. Alcune proprietà delle funzioni $\psi_{sh}(a)$, $\xi_{si}(x')$}

Abbiamo stabilito il risultato di fondamentale importanza che le funzioni $f_i(x,a)$ che danno una serie $\infty^r$ di trasformazioni
\begin{equation*}\tag{6}
 x'_i=f_i(x,a)
\end{equation*}
costituenti un gruppo, soddisfano necessariamente ad un sistema di equazioni differenziali della forma (A). Diciamo però subito che le condizioni espresse dalle (A) necessarie perchè le (6) costituiscano un gruppo non sono a ciò del tutto sufficienti,\footnote{Un esempio dimostra subito la cosa. Il sistema semplice
\[
 \frac{dx'_1}{da}=\frac{x'_1}{x'_2},\qquad \frac{dx'_2}{da}=1
\]
con un solo parametro $a$, ha la forma tipica (A), ove si faccia $\xi_{11}=x'_1/x'_2$, $\xi_{12}=1$, $\psi_{11}(a)=1$; ma gli integrali generali
\[
 x'_1=x_1(x_2+a),\qquad x'_2=x_2+a
\]
non danno le trasformazioni di un gruppo.} ma si richiedono ancora alcune condizioni complementari che verranno esaminate nel prossimo capitolo. Intanto osserviamo alcune proprietà semplici ma importanti che si verificano per ogni serie $\infty^r$ di trasformazioni (6), quando siano soddisfatte le equazioni differenziali di un sistema (A), anche se quella serie non costituisse un gruppo.

In primo luogo diciamo che il determinante delle funzioni $\psi_{hk}(a)$
\[
 \psi=\left|\begin{array}{cccc}
 \psi_{11}&\psi_{12}&\cdots&\psi_{1r}\\
 \psi_{21}&\psi_{22}&\cdots&\psi_{2r}\\
 \cdots&\cdots&\cdots&\cdots\\
 \psi_{r1}&\psi_{r2}&\cdots&\psi_{rr}
 \end{array}\right|
\]
è necessariamente diverso da zero. In caso contrario infatti potremmo determinare $r$ funzioni
\[
 A_1,A_2,\ldots,A_r
\]
delle $a$, non tutte nulle, tali da soddisfare il sistema delle $r$ equazioni lineari omogenee
\[
 \sum\limits_h^{1,\ldots,r} A_h\psi_{hs}(a)=0\qquad (s=1,2,\ldots,r).
\]
Dalle equazioni differenziali (A) risulterebbero allora le seguenti
\[
 \sum\limits_h^{1,\ldots,r} A_h\frac{\partial f_i(x,a)}{\partial a_h}=0\qquad (i=1,2,\ldots,n),
\]
dunque (\S{}2) gli $r$ parametri $a$ non sarebbero essenziali, contro l'ipotesi.

Dopo ciò, se consideriamo le $r$ equazioni (A) che si ottengono tenendo fisso $i$ e ponendo $k=1,2,\ldots,r$, è chiaro che potremo risolverle appunto rispetto a
\[
 \xi_{1i},\xi_{2i},\ldots,\xi_{ri},
\]
e le formole risolutive prenderanno l'aspetto seguente:
\begin{equation*}\tag{B}
 \xi_{ji}(x')=\sum\limits_k^{1,\ldots,r}\alpha_{jk}(a)\frac{\partial x'_i}{\partial a_k},
\end{equation*}
dove con $\alpha_{jk}$ è indicato il complemento algebrico di $\psi_{jk}$ in $\psi$, diviso per $\psi$ stesso.

Segue di qui un risultato importante per il seguito e cioè che è impossibile determinare $r$ costanti non tutte nulle
\[
 e_1,e_2,\ldots,e_r
\]
che soddisfino le $n$ equazioni
\begin{equation*}\tag{12}
 \sum\limits_j^{1,\ldots,r}e_j\xi_{ji}(x')=0\qquad (i=1,2,\ldots,n).
\end{equation*}
Le (12) ci darebbero infatti per le (B) le equazioni
\[
 \sum\limits_k^{1,\ldots,r}\left(\sum\limits_j^{1,\ldots,r}e_j\alpha_{jk}(a)\right)
 \frac{\partial f_i(x,a)}{\partial a_k}=0,
\]
ed i parametri $a_1,a_2,\ldots,a_r$ non sarebbero tutti essenziali.

Formuliamo i risultati ottenuti nel teorema:
Se una serie $\infty^r$ di trasformazioni
\[
 x'_i=f_i(x,a)
\]
soddisfa ad un sistema di equazioni differenziali del tipo (A), 1. il determinante delle $\psi$ è necessariamente diverso da zero, e le (A), risolute rapporto alle $\xi$, possono porsi sotto la forma equivalente (B); 2. è impossibile determinare $r$ costanti $e_1,e_2,\ldots,e_r$ (non nulle) che soddisfino alle $n$ equazioni lineari
\[
 \sum\limits_j^{1,\ldots,r}e_j\xi_{ji}(x')=0\qquad (i=1,2,\ldots,n).
\]

\section*{\S{} 9. Gruppi ad un parametro}

Cominciamo dall'applicare questi primi principî della teoria al caso dei gruppi ad un parametro; questi formano in certa guisa gli elementi con cui si compongono i gruppi continui finiti (ed anche infiniti) e si può dire che compiono nella teoria generale dei gruppi continui un ufficio analogo a quello dei gruppi di potenze nella teoria dei gruppi discontinui.

Se la serie $\infty^1$ di trasformazioni
\begin{equation*}\tag{13}
 x'_i=f_i(x_1,x_2,\ldots,x_n;a)\qquad (i=1,2,\ldots,n)
\end{equation*}
con un parametro $a$ costituisce un gruppo, pel teorema fondamentale del \S{}7 le $x'_i$, considerate come funzioni di $a$, soddisferanno ad un sistema di equazioni differenziali ordinarie della forma
\begin{equation*}\tag{14}
 \frac{dx'_i}{da}=\psi(a)\cdot \xi_i(x'_1,x'_2,\ldots,x'_n),
 \qquad i=1,2,\ldots,n,
\end{equation*}
e ne costituiranno l'integrale generale.

Per determinare completamente le $f_i(x,a)$ basterà dunque conoscere la forma che queste $f_i$, ossia le $x'_i$, assumono per un particolare valore di $a$, diciamo per $a=a_0$.

È chiaro che il sistema integrale delle (14), cioè delle equazioni
\[
 \frac{dx'_1}{\xi_1}=\frac{dx'_2}{\xi_2}=\cdots=\frac{dx'_n}{\xi_n}=\psi(a)\,da,
\]
può porsi sotto la forma risoluta rispetto alle costanti arbitrarie
\begin{equation*}\tag{15}
\left\{
\begin{aligned}
 \Omega_1(x'_1,x'_2,\ldots,x'_n)&=C_1,\\
 \Omega_2(x'_1,x'_2,\ldots,x'_n)&=C_2,\\
 &\hspace{1.5em}\cdots\cdots\cdots\\
 \Omega_{n-1}(x'_1,x'_2,\ldots,x'_n)&=C_{n-1},\\
 \Omega_n(x'_1,x'_2,\ldots,x'_n)&=\int_{a_0}^{a}\psi(a)\,da+C_n,
\end{aligned}
\right.
\end{equation*}
dove le $\Omega_i(x'_1,x'_2,\ldots,x'_n)$ sono certe $n$ funzioni indipendenti delle $x'_1,x'_2,\ldots,x'_n$ e le $C_1,C_2,\ldots,C_n$ sono costanti arbitrarie.

Supponiamo ora che nella serie (13) sia contenuta la sostituzione identica e corrisponda al valore $a_0$ di $a$, talchè si abbia
\[
 x'_i=x_i\qquad\text{per }a=a_0.
\]
Facendo $a=a_0$ nelle (15) ne viene
\[
 C_i=\Omega_i(x_1,x_2,\ldots,x_n),
\]
ed il sistema integrale (15), cangiato per semplicità il parametro $a$ in
\[
 t=\int_{a_0}^{a}\psi(a)\,da,
\]
prende la forma
\begin{equation*}\tag{16}
\left\{
\begin{aligned}
 \Omega_i(x'_1,x'_2,\ldots,x'_n)&=\Omega_i(x_1,x_2,\ldots,x_n)\qquad (i=1,2,\ldots,n-1),\\
 \Omega_n(x'_1,x'_2,\ldots,x'_n)&=\Omega_n(x_1,x_2,\ldots,x_n)+t,
\end{aligned}
\right.
\end{equation*}
cioè è perfettamente equivalente alle (13) sotto l'ipotesi
\[
 x'_i=x_i\quad\text{per }t=0\quad (a=a_0).
\]
Le equazioni (16) manifestano nel modo più semplice la proprietà di gruppo poichè, se cangiamo le variabili $x_i$ nelle
\[
 y_i=\Omega_i(x_1,x_2,\ldots,x_n)\qquad i=1,2,\ldots,n,
\]
esse prendono la forma semplicissima
\[
 y'_1=y_1,\quad y'_2=y_2,\quad\ldots,\quad y'_{n-1}=y_{n-1},\quad y'_n=y_n+t.
\]
Questa trasformazione si dirà una traslazione dello spazio ad $n$ dimensioni $(y_1,y_2,\ldots,y_n)$.

In generale se in un gruppo qualunque $G$ si fa contemporaneamente sulle variabili primitive e sulle trasformate una medesima sostituzione $T$ ed anche, se si vuole, si cangiano i parametri in modo arbitrario, le considerazioni alla fine del \S{}1 provano che le nuove trasformazioni percorrono nuovamente un gruppo, il gruppo $T^{-1}GT$. Due tali gruppi, ottenuti l'uno dall'altro con un cangiamento di variabili e di parametri, dicesi simili.

Possiamo dunque enunciare questo primo risultato:
Se un gruppo continuo ad un parametro contiene l'identità, esso è simile ad un gruppo di traslazioni, e le sue trasformazioni sono permutabili fra loro e a due a due inverse l'una dell'altra.

Supponiamo ora, più in generale, di avere una serie $\infty^1$ di trasformazioni (13), che soddisfi ad un sistema di equazioni differenziali della forma (14), e indichiamo con
\[
 \varphi_1(x),\varphi_2(x),\ldots,\varphi_n(x)
\]
le funzioni delle $x$ in cui si riducono le $f_i(x,a)$ per $a=a_0$; il sistema (15) equivalente alle (13) si scriverà:
\[
\left\{
\begin{aligned}
 \Omega_i(x'_1,x'_2,\ldots,x'_n)&=\Omega_i(\varphi_1,\varphi_2,\ldots,\varphi_n)
 \qquad (i=1,2,\ldots,n-1),\\
 \Omega_n(x'_1,x'_2,\ldots,x'_n)&=\Omega_n(\varphi_1,\varphi_2,\ldots,\varphi_n)+t.
\end{aligned}
\right.
\]
Questa trasformazione si ottiene eseguendo, primieramente, la trasformazione
\[
 y_1=\varphi_1(x_1,x_2,\ldots,x_n),\quad\ldots,\quad y_n=\varphi_n(x_1,x_2,\ldots,x_n),
\]
indi l'altra
\[
\left\{
\begin{aligned}
 \Omega_i(x'_1,x'_2,\ldots,x'_n)&=\Omega_i(y_1,y_2,\ldots,y_n)\qquad (i=1,2,\ldots,n-1),\\
 \Omega_n(x'_1,x'_2,\ldots,x'_n)&=\Omega_n(y_1,y_2,\ldots,y_n)+t.
\end{aligned}
\right.
\]
La prima è una trasformazione fissa della serie (13) (corrispondente ad $a=a_0$), e la seconda, variando $t$, genera un gruppo della specie sopra considerata. Abbiamo dunque il teorema:

Se una serie $\infty^1$ di trasformazioni (13) soddisfa ad equazioni differenziali del tipo (14), le sue trasformazioni si ottengono componendo una trasformazione fissa della serie stessa con tutte le trasformazioni del gruppo ad un parametro (colla sostituzione identica) definito dalle equazioni differenziali (14) stesse.

\section*{\S{} 10. Trasformazione infinitesima di un $G_1$}

Ritorniamo alla considerazione dei gruppi ad un parametro contenenti l'identità e quindi contenenti insieme ad ogni trasformazione la sua inversa (\S{}9). Tali gruppi s'indicheranno d'ora innanzi col simbolo $G_1$, l'indice riferendosi al numero dei parametri. Le formole relative ad un tale gruppo $G_1$ si ottengono, come si è visto, integrando un sistema di equazioni differenziali
\[
 \frac{dx'_i}{dt}=\xi_i(x'_1,x'_2,\ldots,x'_n)
\]
colle condizioni iniziali
\[
 x'_i=x_i\qquad\text{per }t=0.
\]
Le funzioni $\xi_i(x')$ sono supposte analitiche, e se si ammette di più che il punto $(x'_1,x'_2,\ldots,x'_n)$ sia un punto regolare per queste funzioni, sappiamo che gli integrali $x'_i$ saranno sviluppabili in serie di potenze di $t$ secondo la formula
\[
 x'_i=x_i+\xi_i(x)t+\frac{t^2}{1\cdot2}\sum\limits_k^{1,\ldots,n} \xi_k\frac{\partial\xi_i}{\partial x_k}+\cdots .
\]
Possiamo scrivere questa in modo più espressivo introducendo, per una funzione qualunque $f(x_1,x_2,\ldots,x_n)$ delle variabili $x$, il simbolo d'operazione
\[
 Xf=\sum\limits_i^{1,\ldots,n}\xi_i\frac{\partial f}{\partial x_i}
\]
e indicando con $X^2f$, $X^3f,\ldots$ il risultato dell'operazione stessa ripetuto due, tre, \ldots volte; lo sviluppo precedente prende allora la forma:
\begin{equation*}\tag{17}
 x'_i=x_i+tXx_i+\frac{t^2}{1\cdot2}X^2x_i+\frac{t^3}{1\cdot2\cdot3}X^3x_i+\cdots .
\end{equation*}
Più in generale, se vogliamo sviluppare secondo potenze di $t$ una funzione qualunque $f(x'_1,x'_2,\ldots,x'_n)$, regolare nel punto $(x_1,x_2,\ldots,x_n)$, avremo
\begin{equation*}\tag{18}
\begin{aligned}
 f(x'_1,x'_2,\ldots,x'_n)&=f(x_1,x_2,\ldots,x_n)+tXf+\frac{t^2}{1\cdot2}X^2f+\cdots\\
 &\quad+\frac{t^s}{1\cdot2\cdots s}X^sf+\cdots .
\end{aligned}
\end{equation*}
Qualche volta indicheremo questo sviluppo per brevità col simbolo
\[
 f(x'_1,x'_2,\ldots,x'_n)=e^{tX}f,
\]
dove si svilupperà formalmente $e^{tX}$ colla serie esponenziale e le potenze di $X$ si considereranno qui come operanti sulla $f$.

Consideriamo ora quella trasformazione del gruppo $G_1$ che corrisponde ad un valore infinitamente piccolo $\delta t$ del parametro $t$. Questa trasformazione infinitesima del gruppo, quando si trascurino gl'infinitesimi d'ordine superiore a $\delta t$, dà alle $x$ gl'incrementi infinitesimi
\[
 \delta x_i=\xi_i\,\delta t\qquad (i=1,2,\ldots,n),
\]
e quindi ad una funzione qualunque $f$ l'incremento
\[
 \delta f=Xf\,\delta t.
\]
È chiaro che l'espressione effettiva del simbolo $Xf$ determina di per sè completamente la trasformazione infinitesima, poichè fornisce i valori delle $\xi$. Perciò il simbolo stesso $Xf$ è riguardato da Lie come simbolo della trasformazione stessa, cioè dice quindi la trasformazione infinitesima $Xf$. Naturalmente se si cambia il parametro $t$ in un nuovo $\tau$ ponendo
\[
 t=\varphi(\tau),
\]
i valori delle $\xi$ vengono moltiplicati contemporaneamente per un fattore costante, cioè pel valore di $\dfrac{dt}{d\tau}$ per $\tau=0$. Non si altera dunque una trasformazione infinitesima moltiplicandola per una costante $c$ qualunque. Si osservi ora che le equazioni differenziali
\[
 \frac{dx'_i}{dt}=\xi_i(x'_1,x'_2,\ldots,x'_n),
\]
interpretate geometricamente nello spazio a $n$ dimensioni $(x'_1,x'_2,\ldots,x'_n)$, determinano un sistema $\infty^{n-1}$ di curve in regola; queste diconsi le traiettorie del gruppo, ovvero anche, secondo Monge, le caratteristiche dell'equazione a derivate parziali
\[
 Xf=0.
\]
Possiamo quindi dire: Il simbolo $Xf$ della trasformazione infinitesima non è altro che la derivata della funzione $f$ calcolata nella direzione della traiettoria del gruppo.

Introdotto così il concetto di trasformazione infinitesima, vediamo che ogni gruppo $G_1$ contiene una trasformazione infinitesima $Xf$, determinata a meno di un fattore costante, quella alla sua volta individua perfettamente il gruppo, che viene generato dalla $Xf$ secondo le formule (17) o (18). Il gruppo stesso si indica perciò semplicemente come gruppo $Xf$.

Formuliamo ancora il teorema alla fine del \S{}9 sotto la nuova forma:
Ogni gruppo ad un parametro
\[
 x'_i=f_i(x_1,x_2,\ldots,x_n;a)
\]
individua una trasformazione infinitesima $Xf$, e le sue trasformazioni si ottengono tutte combinando una sua trasformazione fissa colle trasformazioni del gruppo $G_1$, contenente l'identità, generato da $Xf$.

\section*{\S{} 11. Cangiamento di variabili nelle trasformazioni infinitesime}

Si è visto come un gruppo $G_1$, contenente l'identità, risulta perfettamente determinato dalla sua trasformazione infinitesima $Xf$. Ora se, mediante un cangiamento di variabili indipendenti espresso dalle formole
\begin{equation*}\tag{19}
 y_i=\varphi_i(x_1,x_2,\ldots,x_n),
\end{equation*}
che fanno passare dalle antiche variabili $x$ alle nuove $y$, si trasforma il gruppo $G_1$ in un gruppo simile (\S{}9) $\Gamma_1$, si può domandare come si calcolerà la trasformazione infinitesima $Yf$ del gruppo $\Gamma_1$. È facile vedere che il nuovo simbolo $Yf$ si otterrà semplicemente esprimendo nell'antico $Xf$ le $x$ per le $y$. Se supponiamo infatti
\[
 f(x_1,x_2,\ldots,x_n)=F(y_1,y_2,\ldots,y_n),
\]
avremo
\[
 Xf=\sum\limits_\nu^{1,\ldots,n}\frac{\partial F}{\partial y_\nu}\,Xy_\nu.
\]
Espresse le $Xy_\nu$ per le $y$, si ponga
\[
 Xy_\nu=\eta_\nu(y_1,y_2,\ldots,y_n),
\]
e
\[
 YF=\sum\limits_\nu^{1,\ldots,n}\eta_\nu\frac{\partial F}{\partial y_\nu};
\]
qualunque sia $f$, si avrà identicamente
\[
 Xf=YF,
\]
poichè le $x,y$ siano legate dalle formole di trasformazione (19). Sarà quindi anche
\[
 X^2f=Y^2F,
\]
in generale
\[
 X^sf=Y^sF.
\]
In conseguenza la (18) si muta in
\[
\begin{aligned}
 F(y'_1,y'_2,\ldots,y'_n)&=F(y_1,y_2,\ldots,y_n)+tYF+\frac{t^2}{1\cdot2}Y^2F+\cdots\\
 &=e^{tY}F,
\end{aligned}
\]
formula che vale per qualunque $F$ e dimostra la proprietà enunciata.

Così ogni trasformazione finita
\[
 x'_i=x_i+tXx_i+\frac{t^2}{1\cdot2}X^2x_i+\cdots
\]
di $G_1$ si muta pel gruppo trasformato nell'altra
\[
 y'_i=y_i+tYy_i+\frac{t^2}{1\cdot2}Y^2y_i+\cdots,
\]
il valore di $t$ essendo lo stesso nelle due formole.

\section*{\S{} 12. Le serie $\infty^r$ di trasformazioni e le relative trasformazioni infinitesime}

Scopo delle ricerche che andiamo ora a fare è quello di estendere ai gruppi continui finiti a più parametri i risultati sopra conseguiti pei gruppi ad un parametro. In particolare dimostreremo che un gruppo ad $r$ parametri, che contenga l'identità e insieme ad ogni trasformazione la sua inversa, si genera mediante $r$ trasformazioni infinitesime
\[
 X_1f,\ X_2f,\ldots,\ X_rf
\]
fra di loro linearmente indipendenti, nel senso che non abbia luogo alcuna identità della forma
\[
 e_1X_1f+e_2X_2f+\cdots+e_rX_rf=0,
\]
dove le $e$ siano costanti numeriche.

Per stabilire questi risultati fondamentali è utile premettere alcune considerazioni generali sulle serie $\infty^r$ di trasformazioni in generale:
\begin{equation*}\tag{20}
 x'_i=f_i(x_1,x_2,\ldots,x_n; a_1,a_2,\ldots,a_r)
\end{equation*}
e le trasformazioni infinitesime che esse determinano. Così stabiliremo, fra altro, nuovamente le equazioni differenziali (A) del \S{}7, riconoscendone il significato intrinseco.

Prendiamo una trasformazione qualunque $S_a$\footnote{Indichiamo al solito simbolicamente con $S_a$ quella trasformazione che corrisponde ai valori $a_1,a_2,\ldots,a_r$ dei parametri in (20).} della serie (20), ed una trasformazione infinitamente vicina $S_{a+\delta a}$ della medesima serie, ove i parametri $a$ hanno ricevuto incrementi infinitesimi $\delta a$. Se osserviamo che si ha
\[
 S_{a+\delta a}=S_a(S_a^{-1}S_{a+\delta a}),
\]
\[
 S_{a+\delta a}=(S_{a+\delta a}S_a^{-1})S_a,
\]
vediamo che dalla $S_a$ si può passare all'infinitamente vicino $S_{a+\delta a}$ o eseguendo prima la $S_a$, indi la trasformazione infinitesima $S_a^{-1}S_{a+\delta a}$, od anche eseguendo per prima la trasformazione infinitesima $S_{a+\delta a}S_a^{-1}$ e poi la $S_a$. Rispetto a queste due trasformazioni infinitesime si osservi che esse non appartengono in generale alla serie (20) e che l'una è la trasformazione dell'altra per mezzo della $S_a$, come si vede dalla formula
\[
 S_a^{-1}S_{a+\delta a}=S_a^{-1}(S_{a+\delta a}S_a^{-1})S_a.
\]

Si possono scrivere facilmente le espressioni analitiche effettive di queste due trasformazioni infinitesime. Se si suppone infatti che le (20), risolute rispetto alle $x$, dànno
\begin{equation*}\tag{20*}
 x_i=F_i(x'_1,x'_2,\ldots,x'_n; a_1,a_2,\ldots,a_r),
\end{equation*}
la prima trasformazione $S_a^{-1}S_{a+\delta a}$ sarà data da
\[
\begin{aligned}
(S_a^{-1}S_{a+\delta a})\quad x'_i={}&f_i\bigl(F_1(x,a),F_2(x,a),\ldots,F_n(x,a);\\
& a_1+\delta a_1,a_2+\delta a_2,\ldots,a_r+\delta a_r\bigr),
\end{aligned}
\]
e la seconda da
\[
\begin{aligned}
(S_{a+\delta a}S_a^{-1})\quad x'_i={}&F_i\bigl(f_1(x,a+\delta a),f_2(x,a+\delta a),\ldots,f_n(x,a+\delta a);\\
& a_1,a_2,\ldots,a_r\bigr).
\end{aligned}
\]
Siccome le (20), (20*) sono inverse l'una dell'altra, si può scrivere per questo
\begin{equation*}\tag{21}
(S_a^{-1}S_{a+\delta a})\quad
 x'_i=x_i+\sum\limits_k^{1,\ldots,r}\left[\frac{\partial f_i(\xi,a)}{\partial a_k}\right]_{\xi=F(x,a)}\delta a_k,
\end{equation*}
e
\[
(S_{a+\delta a}S_a^{-1})\quad
 x'_i=x_i+\sum\limits_k^{1,\ldots,r}\delta a_k\sum\limits_\nu^{1,\ldots,n}\frac{\partial f_\nu(x,a)}{\partial a_k}
 \left[\frac{\partial F_i(\xi,a)}{\partial \xi_\nu}\right]_{\xi=f(x,a)}.
\]
L'espressione della seconda può semplificarsi osservando che dall'identità
\[
 F_i(f_1(x,a),f_2(x,a),\ldots,f_n(x,a))=x_i,
\]
derivando rapporto ad $a_k$, deduciamo
\[
 \left[\frac{\partial F_i(\xi,a)}{\partial a_k}\right]_{\xi=f(x,a)}+
 \sum\limits_\nu^{1,\ldots,n}\left[\frac{\partial F_i(\xi,a)}{\partial \xi_\nu}\right]_{\xi=f(x,a)}\frac{\partial f_\nu(x,a)}{\partial a_k}=0,
\]
onde risulta per la $S_{a+\delta a}S_a^{-1}$
\begin{equation*}\tag{21*}
(S_{a+\delta a}S_a^{-1})\quad
 x'_i=x_i-\sum\limits_k^{1,\ldots,r}\left[\frac{\partial F_i(\xi,a)}{\partial a_k}\right]_{\xi=f(x,a)}\delta a_k.
\end{equation*}
Se immaginiamo fissate le $a$ e variabili ad arbitrio le $\delta a$, le (21), (21*) coordinano alla $S_a$ due serie di trasformazioni infinitesime mediante le quali si passa da $S_a$ ad qualunque altra trasformazione infinitamente vicina nella serie (20). Ora osserviamo che la (21) si scrive, colla notazione delle trasformazioni infinitesime,
\[
 Z\Phi=\sum\limits_k^{1,\ldots,r}\delta a_k\sum\limits_i^{1,\ldots,n}\left[\frac{\partial f_i(\xi,a)}{\partial a_k}\right]_{\xi=F(x,a)}\frac{\partial\Phi}{\partial x_i},
\]
ed è quindi una combinazione lineare delle $r$ trasformazioni infinitesime
\begin{equation*}\tag{22}
 Z_k\Phi=\sum\limits_i^{1,\ldots,n}\left[\frac{\partial f_i(\xi,a)}{\partial a_k}\right]_{\xi=F(x,a)}\frac{\partial\Phi}{\partial x_i}.
\end{equation*}
Queste, per valori qualunque delle $a$, sono certo fra loro linearmente indipendenti, altrimenti gli $r$ parametri $a$ non sarebbero tutti essenziali. La stessa cosa vale naturalmente per la serie (21*), la quale non è altro che una trasformata della serie (21). Viceversa è chiaro che in una serie qualunque (21) di trasformazioni, se le $r$ trasformazioni infinitesime $Z_k$ date dalla (22) sono indipendenti, gli $r$ parametri $a$ saranno essenziali.

Si osservi infine che la trasformazione infinitesima generica (21)
\[
 Z\Phi=\sum\limits_k^{1,\ldots,r}\delta a_k Z_k\Phi,
\]
dipendendo dagli $r-1$ rapporti delle $\delta a$, percorre una serie $\infty^{r-1}$ di trasformazioni infinitesime distinte.

Riepilogando i risultati ottenuti abbiamo adunque:
Ogni serie (20) $\infty^r$ di trasformazioni coordina ad ogni sua trasformazione $S_a$ due serie (21), (21*) di $\infty^{r-1}$ trasformazioni infinitesime, mediante le quali si passa, nel modo indicato, dalla $S_a$ a tutte le trasformazioni infinitamente vicine nella serie stessa.


\section*{\S{} 13. Caso particolare di un gruppo}

Supponiamo ora che la serie $\infty^r$ data dalle (20) costituisca un gruppo. Allora due trasformazioni qualunque $S_a,S_b$ della serie si compongono in una terza trasformazione $S_c$ della serie stessa:
\begin{equation*}\tag{23}
 S_aS_b=S_c,
\end{equation*}
dove i parametri $c$ sono certe funzioni di $a,b$
\begin{equation*}\tag{24}
 c_d=\chi_d(a_1,a_2,\ldots,a_r; b_1,b_2,\ldots,b_r),\qquad d=1,2,\ldots,r,
\end{equation*}
e queste relazioni, come abbiamo dimostrato al \S{}6, sono risolubili tanto rispetto alle $a$ che alle $b$. Del resto ciò risulta ora subito dall'eguaglianza simbolica (23), osservando che, se teniamo fisse le $a$ e facciamo variare arbitrariamente le $b$, la $S_c$ percorrerà, come la $S_b$, una serie $\infty^r$ di trasformazioni, e perciò le $r$ funzioni $c$ sono indipendenti rispetto alle $b$, cioè
\[
 \frac{\partial(c_1,c_2,\ldots,c_r)}{\partial(b_1,b_2,\ldots,b_r)}\ne0;
\]
analogamente dicasi di
\[
 \frac{\partial(c_1,c_2,\ldots,c_r)}{\partial(a_1,a_2,\ldots,a_r)}.
\]
Diamo ora tanto alle $a$ che alle $b$ incrementi arbitrari $\delta a,\delta b$; basterà calcolare i corrispondenti incrementi $\delta c$ dalle (24) colle formole
\[
 \delta c_d=\sum\limits_k^{1,\ldots,r}\left(\frac{\partial\chi_d}{\partial a_k}\delta a_k+\frac{\partial\chi_d}{\partial b_k}\delta b_k\right)
\]
e si avrà
\[
 S_{a+\delta a}S_{b+\delta b}=S_{c+\delta c}.
\]
Ora prendiamo i $\delta a,\delta b$ così legati fra loro che ne risultino nulli i $\delta c$, cioè
\begin{equation*}\tag{25}
 \sum\limits_k^{1,\ldots,r}\frac{\partial\chi_\lambda}{\partial a_k}\delta a_k+
 \sum\limits_k^{1,\ldots,r}\frac{\partial\chi_\lambda}{\partial b_k}\delta b_k=0,
 \qquad \lambda=1,2,\ldots,r,
\end{equation*}
e allora risulterà
\[
 S_{a+\delta a}S_{b+\delta b}=S_c=S_aS_b,
\]
ossia
\[
 S_a^{-1}S_{a+\delta a}\, S_{b+\delta b}S_b^{-1}=1,
\]
ciò che dimostra che la $S_{b+\delta b}S_b^{-1}$ è l'inversa di $S_a^{-1}S_{a+\delta a}$. Ma evidentemente l'inversa di $S_{b+\delta b}S_b^{-1}$ si ottiene cangiando i segni di tutte le $\delta b$, e perciò come conseguenza delle (24) abbiamo
\begin{equation*}\tag{26}
 S_a^{-1}S_{a+\delta a}=S_{b-\delta b}S_b^{-1}.
\end{equation*}
Si osservi ora che nelle (25) possiamo prendere ad arbitrio le $\delta a$ e trarne le $\delta b$, o inversamente; ne segue che la trasformazione a sinistra nella (26) può identificarsi con una qualunque della serie (21), e così quella a destra con una qualunque della serie (21*). In questo caso dunque le due serie di trasformazioni coincidono. Ma di più se teniamo fisse le $b$ e facciamo variare ad arbitrio le $a$, la (26) dimostra che la prima serie di trasformazioni infinitesime (21), coordinata alla $S_a$, rimane la stessa al variare di $S_a$, e altrettanto naturalmente accade della seconda serie (21*).

Concludiamo adunque: Ad ogni gruppo di trasformazioni a $r$ parametri è coordinata una serie perfettamente determinata di $\infty^{r-1}$ trasformazioni infinitesime, per le quali moltiplicando sia a destra che a sinistra la $S_a$ si ottengono tutte le trasformazioni del gruppo infinitamente vicine alla $S_a$.

\section*{\S{} 14. Nuova deduzione delle equazioni differenziali fondamentali (A) \S{}7}

Traduciamo ora in formole l'eguaglianza simbolica (26), sostituendo per le trasformazioni infinitesime dei due membri le loro espressioni effettive date dalle (21), (21*), ed avremo:
\begin{equation*}\tag{27}
 \sum\limits_k^{1,\ldots,r}\left[\frac{\partial f_i(\xi,a)}{\partial a_k}\right]_{\xi=F(x,a)}\delta a_k
 =
 \sum\limits_j^{1,\ldots,r}\left[\frac{\partial F_i(\xi,b)}{\partial b_j}\right]_{\xi=f(x,b)}\delta b_j,
 \qquad i=1,2,\ldots,n,
\end{equation*}
relazioni che sono una conseguenza delle (25). Ora risolvendo le dette equazioni (25) rapporto alle $\delta b$ scriviamo
\begin{equation*}\tag{28}
 \delta b_j=\sum\limits_k^{1,\ldots,r}\psi_{jk}(a,b)\delta a_k,
\end{equation*}
e le $\psi_{jk}$ avranno precisamente il significato loro attribuito al \S{}7, poichè rappresentano le derivate delle $b$ rapporto alle $a$, quando si tengono fisse le $c$. Sostituendo i valori (28) nelle (27), e ricordando che le $\delta a$ sono indipendenti, ne dedurremo
\begin{equation*}\tag{29}
 \left[\frac{\partial f_i(\xi,a)}{\partial a_k}\right]_{\xi=F(x,a)}
 =
 \sum\limits_j^{1,\ldots,r}\left[\frac{\partial F_i(\xi,b)}{\partial b_j}\right]_{\xi=f(x,b)}\psi_{jk}(a,b).
\end{equation*}
Se introduciamo ora la notazione
\begin{equation*}\tag{a}
 \left[\frac{\partial F_i(\xi,b)}{\partial b_j}\right]_{\xi=f(x,b)}=\Phi_{ji}(x,b),
\end{equation*}
poi nella (29) cangiamo $x_i$ in $x'_i=f_i(x,a)$, otteniamo le equazioni seguenti:
\begin{equation*}\tag{30}
 \frac{\partial x'_i}{\partial a_k}=\sum\limits_j^{1,\ldots,r}\Phi_{ji}(x',b)\psi_{jk}(a,b).
\end{equation*}
Queste sono le equazioni stesse (11) che abbiamo già stabilito al \S{}7, poichè non solo le $\psi_{jk}(a,b)$ hanno il significato loro attribuito al \S{}7, ma anche le $\Phi_{ji}(x',b)$, come risulta dalle osservazioni seguenti. Si ha identicamente
\begin{equation*}\tag{b}
 \left[F_i(\xi,b)\right]_{\xi=f(x,b)}=x_i,
\end{equation*}
e quindi, derivando rapporto a $b_j$,
\[
 \left[\frac{\partial F_i(\xi,b)}{\partial b_j}\right]_{\xi=f(x,b)}
 =-\sum\limits_\lambda^{1,\ldots,n}\left[\frac{\partial F_i(\xi,b)}{\partial\xi_\lambda}\right]_{\xi=f(x,b)}
 \frac{\partial f_\lambda(x,b)}{\partial b_j},
\]
onde per la (a)
\[
 \Phi_{ji}(x,b)= -\sum\limits_\lambda^{1,\ldots,n}\left[\frac{\partial F_i(\xi,b)}{\partial\xi_\lambda}\right]_{\xi=f(x,b)}
 \frac{\partial f_\lambda(x,b)}{\partial b_j}.
\]
Basterà dunque dimostrare, per provare l'asserita coincidenza, che si ha
\[
 \left[\frac{\partial F_i(\xi,b)}{\partial\xi_\lambda}\right]_{\xi=f(x,b)}=\Theta_{\lambda i},
\]
essendo $\Theta_{\lambda i}$ il quoziente del complemento algebrico di $\dfrac{\partial f_\lambda}{\partial x_i}$ nel determinante funzionale
\[
 J=\frac{\partial(f_1,f_2,\ldots,f_n)}{\partial(x_1,x_2,\ldots,x_n)},
\]
diviso per $J$ stesso. Ma ciò risulta subito dalle (b) derivate rapporto ad una $x$ qualunque, per es. $x_\ell$, il che dà\footnote{Si ricordi il significato di: $\varepsilon_{\ell i}=0$ per $\ell\ne i$, $\varepsilon_{\ell i}=1$ per $\ell=i$.}
\[
 \sum\limits_\lambda^{1,\ldots,n}\left[\frac{\partial F_i(\xi,b)}{\partial\xi_\lambda}\right]_{\xi=f(x,b)}
 \frac{\partial f_\lambda(x,b)}{\partial x_\ell}=\varepsilon_{\ell i}.
\]
Ottenute così nuovamente le (30), che coincidono colle (11) \S{}7, basterà dare ancora alle $b$ valori fissi qualunque per dedurne le equazioni differenziali fondamentali (A) \S{}7
\begin{equation*}\tag{A}
 \frac{\partial x'_i}{\partial a_k}=\sum\limits_s^{1,\ldots,r}\xi_{si}(x')\psi_{sk}(a),
 \qquad
 \begin{cases}
 i=1,2,\ldots,n,\\
 k=1,2,\ldots,r.
 \end{cases}
\end{equation*}
Queste appaiono, colla nuova deduzione, sotto un nuovo aspetto come la traduzione in formole dell'equazione simbolica (26) alla quale equivalgono.

\section*{\S{} 15. Le $r$ trasformazioni infinitesime $X_\lambda f=\sum\limits_i^{1,\ldots,n} \xi_{\lambda i}\dfrac{\partial f}{\partial x_i}$ $(\lambda=1,2,\ldots,r)$}

Se supponiamo soltanto che la serie $\infty^r$ di trasformazioni (20) verifichi il sistema di equazioni differenziali (A), la qual cosa avverrà certamente in particolare se la serie costituisce un gruppo, potremo porre sotto altra forma le $r$ trasformazioni infinitesime del \S{}12
\[
 Z_1\Phi,
 Z_2\Phi,
 \ldots,
 Z_r\Phi,
\]
mediante le quali si combinano linearmente tutte le trasformazioni infinitesime (21) che fanno passare da una trasformazione $S_a$ all'infinitamente vicina $S_{a+\delta a}$.

Essendo soddisfatte le (A), gli incrementi $\delta x_i$ che ricevono le $x$ per una qualunque trasformazione infinitesima (21) sono dati da
\[
\begin{aligned}
 \delta x_i
 &=\sum\limits_k^{1,\ldots,r}\left[\frac{\partial f_i(\xi,a)}{\partial a_k}\right]_{\xi=F(x,a)}\delta a_k\\
 &=\sum\limits_k^{1,\ldots,r}\delta a_k\sum\limits_s^{1,\ldots,r}\xi_{si}(x)\psi_{sk}(a).
\end{aligned}
\]
La più generale trasformazione infinitesima in discorso si scrive dunque:
\[
 Xf=\sum\limits_{s,k}^{1,\ldots,r}\delta a_k\psi_{sk}(a)\sum\limits_i^{1,\ldots,n}\xi_{si}\frac{\partial f}{\partial x_i},
\]
ovvero posto
\[
 \sum\limits_k^{1,\ldots,r}\psi_{sk}(a)\delta a_k=\lambda_s,
\]
\[
 Xf=\sum\limits_s^{1,\ldots,r}\lambda_s\sum\limits_i^{1,\ldots,n}\xi_{si}\frac{\partial f}{\partial x_i},
\]
dove le $\lambda$ sono costanti arbitrarie. Se poniamo
\begin{equation*}\tag{31}
 X_1f=\sum\limits_i^{1,\ldots,n}\xi_{1i}\frac{\partial f}{\partial x_i},\quad
 X_2f=\sum\limits_i^{1,\ldots,n}\xi_{2i}\frac{\partial f}{\partial x_i},\quad \ldots,\quad
 X_rf=\sum\limits_i^{1,\ldots,n}\xi_{ri}\frac{\partial f}{\partial x_i},
\end{equation*}
vediamo che la $Xf$ è una combinazione lineare a coefficienti costanti di queste $r$ trasformazioni infinitesime fondamentali
\[
 X_1f,
 X_2f,
 \ldots,
 X_rf.
\]
Queste sono linearmente indipendenti, come risulta da quanto abbiamo già detto al \S{}12, od anche dalle osservazioni alla fine del \S{}8, dalle quali segue che non sussiste alcuna identità della forma
\[
 e_1X_1f+e_2X_2f+\cdots+e_rX_rf=0
\]
colle $e$ costanti.

Dunque: Se una serie $\infty^r$ di trasformazioni soddisfa le equazioni differenziali (A), si passa da una trasformazione qualunque $S_a$ della serie ad una infinitamente vicina $S_{a+\delta a}$, componendo la $S_a$ con una trasformazione infinitesima $\sum\limits_s^{1,\ldots,r}\lambda_sX_sf$, combinazione lineare a coefficienti $\lambda$ costanti delle $r$ trasformazioni infinitesime fisse, linearmente indipendenti,
\[
 X_sf=\sum\limits_i^{1,\ldots,n}\xi_{si}\frac{\partial f}{\partial x_i}\qquad(s=1,2,\ldots,r).
\]
Se si ricorda (\S{}10) che una trasformazione infinitesima genera un gruppo ad un parametro, è ben naturale domandare ora se la proprietà sopra dimostrata per la trasformazione infinitesima $\sum\limits_s^{1,\ldots,r}\lambda_sX_sf$ vale anche per le trasformazioni finite del gruppo $G_1$ ad un parametro da questa generato, se cioè componendo $S_a$ con una trasformazione finita di $G_1$ si ottiene ancora una trasformazione della serie primitiva.

La risposta è affermativa, come andiamo ora a dimostrare, e da questo risultato si dedurranno conseguenze di fondamentale importanza per la teoria dei gruppi.

\section*{\S{} 16. Teorema fondamentale per le serie di trasformazioni che soddisfano le (A)}

Per dimostrare quanto abbiamo sopra asserito, cominciamo dall'osservare che essendo per valori generici delle $a$ diverso da zero il determinante delle $\psi$, possiamo, come già abbiamo fatto al \S{}8, risolvere le (A) rispetto alle $\xi$, e scriverle sotto la forma equivalente
\begin{equation*}\tag{B}
 \xi_{ji}(x')=\sum\limits_k^{1,\ldots,r}\alpha_{jk}(a)\frac{\partial x'_i}{\partial a_k}.
\end{equation*}
La trasformazione infinitesima $\sum\limits_s^{1,\ldots,r}\lambda_sX_sf$ scritta per disteso è data da
\[
 \delta x_i=\delta t\sum\limits_s^{1,\ldots,r}\lambda_s\xi_{si}(x),\qquad (i=1,2,\ldots,n),
\]
e calcolata dalle (A) si scrive
\[
 \delta x_i=\sum\limits_{s,k}^{1,\ldots,r}\xi_{si}(x)\psi_{sk}(a)\delta a_k,
\]
onde paragonando si ha
\[
 \sum\limits_s^{1,\ldots,r}\xi_{si}\left(\lambda_s\delta t-\sum\limits_k^{1,\ldots,r}\psi_{sk}(a)\delta a_k\right)=0
 \qquad(i=1,2,\ldots,n),
\]
e perciò\footnote{Altrimenti le $X_1f,X_2f,\ldots,X_rf$ non sarebbero indipendenti.}
\[
 \sum\limits_k^{1,\ldots,r}\psi_{sk}(a)\delta a_k=\lambda_s\delta t.
\]
Risolvendo rapporto alle $\delta a_k$, abbiamo le formole
\begin{equation*}\tag{c}
 \delta a_k=\delta t\sum\limits_s^{1,\ldots,r}\alpha_{sk}(a)\lambda_s,
\end{equation*}
colle quali sono da calcolarsi gl'incrementi delle $a$ nel passaggio dalla $S_a$ alla $S_{a+\delta a}$ che risulta dalla composizione di $S_a$ colla trasformazione infinitesima $\sum\limits_s^{1,\ldots,r}\lambda_sX_sf$.

Fissiamo per le $a$ un sistema di valori particolari
\[
 a_1^{(0)},a_2^{(0)},\ldots,a_r^{(0)}
\]
tali soltanto che per essi il determinante delle $\psi_{jk}(a^{(0)})$ e quindi quello delle $\alpha_{jk}(a^{(0)})$ non sia nullo. Fissati poi per le $\lambda$ valori costanti qualunque, si determinino le $a$ in funzione di $t$ secondo le (c), dalle equazioni differenziali
\begin{equation*}\tag{32}
 \frac{da_k}{dt}=\sum\limits_s^{1,\ldots,r}\lambda_s\alpha_{sk}(a)
\end{equation*}
colle condizioni iniziali
\[
 a_k=a_k^{(0)}\qquad\text{per }t=0.
\]
La trasformazione $S_a$ corrispondente al variare di $t$ percorrerà altrettante trasformazioni della serie (20) nell'intorno di $S_{a^{(0)}}$. È anzi facile vedere che, facendo poi variare le $\lambda$, la $S_a$ percorrerà tutte le trasformazioni nell'intorno di $S_{a^{(0)}}$; invero per la forma delle (32) è evidente che gl'integrali $a$ dipenderanno dalle variabili
\[
 u_1=\lambda_1t,
 \quad u_2=\lambda_2t,
 \ldots,
 \quad u_r=\lambda_rt;
\]
e poichè il determinante funzionale delle $a$ rapporto alle $u$ si riduce per $t=0$ appunto ad $\lvert\alpha_{ik}(a^{(0)})\rvert\ne0$, segue che possiamo inversamente esprimere le $u_1,u_2,\ldots,u_r$ per $a_1,a_2,\ldots,a_r$, cosicchè prendendo convenientemente le $\lambda$ e $t$ i parametri $a$ riceveranno valori arbitrarî.

Ora se, come sopra, teniamo fisse le $\lambda$ e facciamo variare $t$, la $S_a$ percorrerà una serie $\infty^1$ di trasformazioni entro la serie completa
\[
 x'_i=f_i(x,a).
\]
Tale che le $x'_i$, considerate come funzioni di $t$, soddisferanno, a causa delle equazioni fondamentali (A) e delle (32), le equazioni differenziali
\[
 \frac{dx'_i}{dt}=\sum\limits_s^{1,\ldots,r}\xi_{si}(x')\sum\limits_k^{1,\ldots,r}\psi_{sk}(a)\sum\limits_j^{1,\ldots,r}\lambda_j\alpha_{jk}(a),
\]
ossia
\[
 \frac{dx'_i}{dt}=\sum\limits_s^{1,\ldots,r}\lambda_s\xi_{si}(x').
\]
Secondo il risultato ottenuto pei gruppi ad un parametro alla fine del \S{}9, la trasformazione $S_a$ nasce quindi eseguendo prima la trasformazione fissa $S_{a^{(0)}}$, poi una trasformazione finita del gruppo $G_1$ contenente l'identità generato dalla trasformazione infinitesima $\sum\limits_j^{1,\ldots,r}\lambda_jX_jf$, e ciò è appunto quello che volevamo provare.

Possiamo enunciare dunque il teorema.
Se una serie $\infty^r$ di trasformazioni soddisfa ad un sistema di equazioni differenziali del tipo (A), ed $S_{a^{(0)}}$ è una sua particolare trasformazione per la quale il determinante delle $\psi_{hk}(a^{(0)})$ non è nullo, ogni sua trasformazione nell'intorno di $S_{a^{(0)}}$ si ottiene componendo la $S_{a^{(0)}}$ con una trasformazione finita del gruppo generato da una combinazione lineare arbitraria, a coefficienti costanti, delle $r$ trasformazioni infinitesime indipendenti
\[
 X_\lambda f=\sum\limits_i^{1,\ldots,n}\xi_{\lambda i}\frac{\partial f}{\partial x_i}.
\]
Viceversa comunque si prendano le $\lambda$, la composizione della $S_{a^{(0)}}$ con una trasformazione finita del gruppo $\sum\limits_j^{1,\ldots,r}\lambda_jX_jf$ dà luogo ad una trasformazione della serie stessa.

\chapter*{Capitolo II}
\addcontentsline{toc}{chapter}{Capitolo II}

\begin{center}
\begin{minipage}{0.88\linewidth}
\centering
I tre teoremi fondamentali della teoria dei gruppi.\par
I gruppi misti.
\end{minipage}
\end{center}

\section*{\S{} 17. Il primo teorema fondamentale}

L'intera teoria dei gruppi continui è stata fondata da Lie sopra tre successivi teoremi, di cui andiamo ora ad occuparci.

Il primo di essi, al quale già altrove abbiamo accennato, dà in certo modo l'inversione del teorema finale al \S{}7, stabilendo che le equazioni differenziali del tipo (A), insieme ad alcune condizioni secondarie, sono caratteristiche per le serie costituenti un gruppo.

Per le serie ad un parametro abbiamo già stabilito il teorema in discorso ed ora vogliamo estenderlo al caso di una serie ad $r$ parametri
\begin{equation*}\tag{1}
 x'_i=f_i(x_1,x_2,\ldots,x_n; a_1,a_2,\ldots,a_r),\qquad i=1,2,\ldots,n,
\end{equation*}
che supponiamo soddisfino ad un sistema di equazioni differenziali del tipo (A)
\begin{equation*}\tag{A}
 \frac{\partial x'_i}{\partial a_k}=\sum\limits_s^{1,\ldots,r}\xi_{si}(x')\psi_{sk}(a),
 \qquad
 \begin{cases}
 i=1,2,\ldots,n,\\
 k=1,2,\ldots,r.
 \end{cases}
\end{equation*}
Basta per ciò applicare i risultati conseguiti nel paragrafo precedente. Cominciamo dal supporre che la serie (1) contenga la trasformazione identica, la quale corrisponda ai valori $a_i=a_i^{(0)}$ dei parametri, e supponiamo inoltre che per questi valori dei parametri che danno $S_{a^{(0)}}=1$ risulti diverso da zero il determinante delle $\psi$. Applicando il teorema alla fine del \S{}16, prendendo per $S_a$ la $S_{a^{(0)}}=1$, vediamo intanto che la serie (1), almeno in un certo intorno dell'identità, non è altro che l'insieme di tutte le trasformazioni degli $\infty^{r-1}$ gruppi $G_1$ generati dalle trasformazioni infinitesime
\[
 \sum\limits_i^{1,\ldots,r} \lambda_i X_i f.
\]
Se poi applichiamo un'altra volta il teorema prendendo per $S_a$ una trasformazione nell'intorno di $S_{a^{(0)}}=1$, e con $S_e$ indichiamo una trasformazione di uno dei gruppi $G_1$, che è pure una trasformazione qualunque della serie (1), risulterà
\[
 S_aS_e=S_c,
\]
dove $S_c$ è una nuova trasformazione della serie. Così è appunto dimostrato che, sotto le condizioni supposte, la serie (1) costituisce un gruppo ed anzi un gruppo a coppie di trasformazioni inverse, poichè tale proprietà appartiene, pel \S{}9, ai gruppi $G_1$ che compongono il nostro gruppo a $r$ parametri.

Riassumendo i risultati ora ottenuti con quelli direttamente stabiliti al \S{}7, possiamo dare al primo teorema fondamentale l'enunciato seguente.

Se una serie $\infty^r$ di trasformazioni
\[
 x'_i=f_i(x_1,x_2,\ldots,x_n; a_1,a_2,\ldots,a_r),\qquad i=1,2,\ldots,n,
\]
forma un gruppo, le $x'_i$, considerate come funzioni delle $a$, soddisfano un sistema di equazioni differenziali della forma
\begin{equation*}\tag{A}
 \frac{\partial x'_i}{\partial a_k}=\sum\limits_s^{1,\ldots,r}\xi_{si}(x')\psi_{sk}(a),
 \qquad
 \begin{cases}
 i=1,2,\ldots,n,\\
 k=1,2,\ldots,r,
 \end{cases}
\end{equation*}
dove il determinante delle $\psi_{sk}(a)$ non è identicamente nullo e le $r$ trasformazioni infinitesime
\[
 X_kf=\sum\limits_i^{1,\ldots,n} \xi_{ki}(x)\frac{\partial f}{\partial x_i}\qquad (k=1,2,\ldots,r)
\]
sono linearmente indipendenti.

Viceversa, se una serie $\infty^r$ di trasformazioni soddisfa ad un sistema di equazioni del tipo (A) e contiene la trasformazione identica come corrispondente ai valori $a_i=a_i^{(0)}$ dei parametri, e di più il determinante $|\psi_{hk}(a^{(0)})|$ non è nullo, la serie stessa è un gruppo a coppie di trasformazioni inverse e coincide colla totalità delle trasformazioni contenute nei gruppi ad un parametro $G_1$ generati dalle $\infty^{r-1}$ combinazioni lineari
\[
 \sum\limits_i^{1,\ldots,r} \lambda_i X_i f.
\]
Nell'ipotesi dell'ultima parte del teorema il gruppo (a coppie di trasformazioni inverse) si indicherà con $G_r$ e si dirà generato dalle $r$ trasformazioni infinitesime
\[
 X_1f,\ X_2f,\ldots,\ X_rf.
\]

\section*{\S{} 18. Casi ulteriori}

Per poter applicare il teorema superiore nella sua seconda parte bisogna che il determinante
\[
 |\psi_{hk}(a^{(0)})|
\]
pei valori $a^{(0)}$ dei parametri che danno l'identità sia diverso da zero. Ma supponiamo ora soltanto che la serie $\infty^r$ di trasformazioni (1) soddisfi ad un sistema (A), e sia $S_{\bar a}$ una particolare trasformazione della serie tale che
\[
 |\psi_{hk}(\bar a)|\ne0.
\]
La serie
\[
 S_{\bar a}^{-1}S_a
\]
contiene $\infty^r$ trasformazioni e, per $a_i=\bar a_i$, dà l'identità. D'altronde le formole effettive per la serie $S_{\bar a}^{-1}S_a$ essendo
\[
 x'_i=f_i\bigl(F_1(x,\bar a),F_2(x,\bar a),\ldots,F_n(x,\bar a); a_1,a_2,\ldots,a_r\bigr),
\]
queste nuove funzioni $x'_i$ soddisfano ancora il sistema (A), perchè dalle funzioni primitive differiscono solo per aver cambiato le $x_i$ in $F_i(x,\bar a)$. Dunque a questa serie $S_{\bar a}^{-1}S_a$ è applicabile il teorema fondamentale, onde abbiamo il risultato:

Se una serie $\infty^r$ di trasformazioni (1) soddisfa ad un sistema (A), e la $S_{\bar a}$ è una tale trasformazione che il determinante $|\psi_{hk}(\bar a)|$ non sia nullo, le $\infty^r$ trasformazioni $S_{\bar a}^{-1}S_a$, variando le $a$, costituiscono un gruppo $G_r$, generato dalle $r$ trasformazioni infinitesime
\[
 X_kf=\sum\limits_i^{1,\ldots,n} \xi_{ki}\frac{\partial f}{\partial x_i}\qquad (k=1,2,\ldots,r).
\]
Se supponiamo ora che la serie $\infty^r$ di trasformazioni (1) costituisca un gruppo, ne segue di per sè che esse soddisfano ad un sistema di equazioni del tipo (A); e presa nel gruppo una trasformazione fissa $S_{\bar a}$ tale che $|\psi_{hk}(\bar a)|\ne0$, le $\infty^r$ trasformazioni $S_{\bar a}^{-1}S_a$ costituiscono un gruppo $G_r$ a coppie di trasformazioni inverse, generato dalle trasformazioni infinitesime $X_1f,X_2f,\ldots,X_rf$. Ora se nel gruppo oltre la $S_{\bar a}$ figura l'inversa $S_{\bar a}^{-1}$, la detta serie $S_{\bar a}^{-1}S_a$ coincide col gruppo stesso. Si ha quindi l'importante teorema.

Ogni gruppo a $r$ parametri, le cui trasformazioni siano due a due inverse,\footnote{Per quanto si è visto, basta che insieme ad una certa trasformazione $S_{\bar a}$ esista nel gruppo l'inversa $S_{\bar a}^{-1}$, purchè sia diverso da zero il determinante $|\psi_{hk}(\bar a)|$; allora il teorema del testo risulta applicabile.} è un gruppo $G_r$ generato da $r$ trasformazioni infinitesime
\[
 X_1f,\ X_2f,\ldots,\ X_rf,
\]
le cui espressioni si ottengono nel noto modo stabilendo il sistema di equazioni differenziali (A) cui soddisfano le trasformazioni del gruppo.

I gruppi $G_r$ a coppie di trasformazioni inverse sono quelli che quasi esclusivamente si presentano, e la semplicità della loro teoria consiste principalmente in ciò che il loro studio può farsi sulle $r$ trasformazioni infinitesime che li generano. È di questi gruppi soltanto che intendiamo sempre parlare in seguito, ove non si avverta esplicitamente il contrario.

\section*{\S{} 19. Il sistema completo di equazioni a derivate parziali equivalente al sistema (A)}

I gruppi sui quali si concentrerà tutto l'interesse dei nostri studi sono, come si è detto, i gruppi $G_r$ generati da $r$ trasformazioni infinitesime
\[
 X_kf=\sum\limits_i^{1,\ldots,n}\xi_{ki}\frac{\partial f}{\partial x_i}\qquad (k=1,2,\ldots,r)
\]
linearmente indipendenti. Ora queste $r$ trasformazioni
\[
 X_1f,\ X_2f,\ldots,\ X_rf
\]
per poter generare un gruppo, non possono essere prese ad arbitrio, ma debbono soddisfare a certe speciali condizioni, poichè altrimenti il complesso di gruppi $G_1$ ad un parametro
\[
 \lambda_1X_1f+\lambda_2X_2f+\cdots+\lambda_rX_rf
\]
formerebbe soltanto una serie $\infty^r$ di trasformazioni senza il carattere di gruppo. Nella espressione di queste condizioni necessarie e sufficienti per la generazione di un gruppo $G_r$ consiste il secondo teorema fondamentale, che Lie chiamò il teorema principale (\emph{Hauptsatz}) come il più importante dei tre. Cominciamo dal ricercare queste condizioni come necessarie. Supposto che la serie (1) costituisca un gruppo $G_r$, dovranno le $f_i(x,a)$, considerate come funzioni delle $a$, soddisfare ad un sistema di equazioni ai differenziali totali (A), che deve essere illimitatamente integrabile (\S{}7). Le condizioni d'integrabilità del sistema (A), che ora andiamo a sviluppare, danno appunto le condizioni cercate.

Sappiamo che l'integrazione di un sistema di equazioni ai differenziali totali come il sistema (A) è perfettamente equivalente a quella di un sistema completo di equazioni lineari omogenee per un'unica funzione incognita delle variabili indipendenti e delle funzioni incognite primitive. Nel caso nostro si tratta di una funzione incognita
\[
 F(x'_1,x'_2,\ldots,x'_n; a_1,a_2,\ldots,a_r)
\]
delle $n+r$ variabili $x',a$ che deve soddisfare alle $r$ equazioni lineari omogenee
\begin{equation*}\tag{2}
 \frac{\partial F}{\partial a_k}+\sum\limits_{j=1}^r\sum\limits_{i=1}^n\xi_{ji}(x')\psi_{jk}(a)\frac{\partial F}{\partial x'_i}=0,
 \qquad k=1,2,\ldots,r.
\end{equation*}
Perchè il sistema sia illimitatamente integrabile, è necessario e sufficiente che il sistema (2) sia completo. Ora per il calcolo delle condizioni d'integrabilità è più utile scrivere le (2) sotto la forma seguente. Nel determinante non nullo delle $\psi$ indichiamo, come al solito, con $\alpha_{jk}(a)$ il complemento algebrico di $\psi_{jk}$ diviso per $\psi$ stesso. Moltiplicando la (2) per $\alpha_{jk}$ e sommando rispetto a $k$ da $1$ ad $r$, e dando a $j$ i suoi $r$ valori, sostituiremo alle (2) il sistema equivalente
\begin{equation*}\tag{2*}
 \sum\limits_k^{1,\ldots,r}\alpha_{jk}(a)\frac{\partial F}{\partial a_k}+\sum\limits_i^{1,\ldots,n}\xi_{ji}(x')\frac{\partial F}{\partial x'_i}=0,
 \qquad j=1,2,\ldots,r.
\end{equation*}
Delle due somme che costituiscono il primo membro, la seconda, i cui coefficienti contengono solo le $x'$, non è altro che l'espressione $X_jf$, ove alle $x$ si sostituiscano le $x'$, e che indicheremo perciò con
\[
 X'_jf=\sum\limits_i^{1,\ldots,n}\xi_{ji}(x')\frac{\partial F}{\partial x'_i}.
\]
Anche i coefficienti della prima somma dipendono dalle sole variabili $a$ rispetto alle quali sono prese le derivate di $F$. Introducendo il simbolo
\begin{equation*}\tag{3}
 A_jF=\sum\limits_k^{1,\ldots,r}\alpha_{jk}(a)\frac{\partial F}{\partial a_k},
\end{equation*}
il sistema completo (2) da integrare, generato da quello delle equazioni differenziali fondamentali (A), si scrive
\begin{equation*}\tag{C'}
 A_jF+X'_jF=0,\qquad j=1,2,\ldots,r.
\end{equation*}

\section*{\S{} 20. Condizioni d'integrabilità}

Affinchè il sistema (C') sia completo debbono tutte le operazioni alternate
\[
 (A_j+X'_j,\ A_k+X'_k),\qquad j,k=1,2,\ldots,r,
\]
comporsi linearmente coi primi membri delle (C') stesse. Si osserva che nelle $A$ non entrano le $x'$ e nelle $X'$ non entrano le $a$; queste condizioni si scrivono
\begin{equation*}\tag{4}
 (A_j,A_k)+(X'_j,X'_k)=\sum\limits_s^{1,\ldots,r} \gamma_{jks}(A_sF+X'_sF),
\end{equation*}
dove le $\gamma_{jks}$ indicano convenienti funzioni delle $x'$ e delle $a$, che poi, come ora si vedrà, si riducono a costanti.

Intanto, dovendo le (4) sussistere per una funzione qualunque $F$ delle $x'$ e quindi anche per una funzione qualsiasi delle sole $a$ o delle sole $x'$, avremo separatamente
\begin{equation*}\tag{5}
 (A_j,A_k)=\sum\limits_s^{1,\ldots,r} \gamma_{jks} A_sF,
\end{equation*}
\begin{equation*}\tag{6}
 (X'_j,X'_k)=\sum\limits_s^{1,\ldots,r} \gamma_{jks} X'_sF.
\end{equation*}
Paragonando nelle (5) i coefficienti di $\partial F/\partial a_i$ da una parte e dall'altra abbiamo
\[
 A_j(\alpha_{ki})-A_k(\alpha_{ji})=\sum\limits_s^{1,\ldots,r} \gamma_{jks}\alpha_{si},\qquad i=1,2,\ldots,r.
\]
Queste possono considerarsi come $r$ equazioni lineari nelle
\[
 \gamma_{jk1},\gamma_{jk2},\ldots,\gamma_{jkr},
\]
e poichè il determinante delle $\alpha$ non è zero, possono risolversi rispetto alle $\gamma_{jks}$ ($s=1,2,\ldots,r$). Segue dunque intanto che le $\gamma_{jks}$ non contengono le $x'$. Similmente, se nella (6) paragoniamo i coefficienti di $\partial F/\partial x_i$ da una parte e dall'altra, abbiamo
\[
 X'_j(\xi_{ki})-X'_k(\xi_{ji})=\sum\limits_s^{1,\ldots,r}\gamma_{jks}\xi_{si}(x'),
\]
onde, derivando rispetto ad una qualunque $a$,
\[
 \sum\limits_s^{1,\ldots,r}\frac{\partial\gamma_{jks}}{\partial a_q}X'_sF=0;
\]
dunque le derivate $\partial\gamma_{jks}/\partial a_q$ debbono essere tutte nulle, giacchè le
\[
 X'_1F,\ X'_2F,\ldots,\ X'_rF
\]
sono linearmente indipendenti. Se ne conclude che le $\gamma_{jks}$, non contenendo nè le $x'$ nè le $a$, sono costanti numeriche; queste indichiamo con $c_{jks}$.

Dunque: Le condizioni d'illimitata integrabilità delle equazioni fondamentali (A) si scrivono
\begin{equation*}\tag{$\alpha$}
 (X'_j,X'_k)=\sum\limits_s^{1,\ldots,r} c_{jks}X'_sF,
\end{equation*}
\begin{equation*}\tag{$\beta$}
 (A_j,A_k)=\sum\limits_s^{1,\ldots,r} c_{jks}A_sF,
\end{equation*}
essendo le $c_{jks}$ costanti numeriche.

\section*{\S{} 21. Le condizioni $(\alpha),(\beta)$ come sufficienti}

Supponiamo ora inversamente che le $r$ trasformazioni infinitesime linearmente indipendenti nelle $x'_1,x'_2,\ldots,x'_n$
\[
 X'_jF=\sum\limits_i^{1,\ldots,n}\xi_{ji}(x')\frac{\partial F}{\partial x'_i}
\]
soddisfino le $(\alpha)$ e le $r$ trasformazioni infinitesime nelle $a$,
\[
 A_jF=\sum\limits_k^{1,\ldots,r}\alpha_{jk}(a)\frac{\partial F}{\partial a_k},
\]
pure linearmente indipendenti, soddisfino le $(\beta)$. Allora il sistema (C') è completo, ed il sistema equivalente (2), essendo risoluto rispetto alle derivate
\[
 \frac{\partial F}{\partial a_1},\ \frac{\partial F}{\partial a_2},\ldots,\frac{\partial F}{\partial a_r},
\]
è Jacobiano di forma normale. Fissato un sistema particolare di valori delle $a$
\[
 a_1^{(0)},a_2^{(0)},\ldots,a_r^{(0)},
\]
tale soltanto che il determinante $|\psi_{ik}(a^{(0)})|$ non sia nullo, consideriamo le soluzioni principali
\[
 F_1,F_2,\ldots,F_n
\]
del sistema (2), che per $a_i=a_i^{(0)}$ si riducono rispettivamente a
\[
 x'_1,x'_2,\ldots,x'_n.
\]
Se risolviamo il sistema di equazioni
\[
 F_i(x'_1,x'_2,\ldots,x'_n; a_1,a_2,\ldots,a_r)=x_i,\qquad i=1,2,\ldots,n,
\]
rispetto alle $x'$ ed otteniamo
\begin{equation*}\tag{1'}
 x'_i=f_i(x_1,x_2,\ldots,x_n; a_1,a_2,\ldots,a_r),
\end{equation*}
queste funzioni $f_i(x,a)$ soddisfano, come sappiamo, al sistema (A) di equazioni ai differenziali totali. Poichè inoltre la serie $\infty^r$ di trasformazioni contiene l'identità, che corrisponde ai valori $a^{(0)}$ delle $a$, e $|\psi_{ik}(a^{(0)})|\ne0$, basta invocare la seconda parte del primo teorema fondamentale (\S{}17) per vedere che le (1') definiscono un gruppo $G_r$ generato dalle $r$ trasformazioni infinitesime $X_1f,X_2f,\ldots,X_rf$.

\section*{\S{} 22. Nuova dimostrazione del primo teorema fondamentale}

Invece di ricorrere, come ultimamente abbiamo fatto, alla seconda parte del primo teorema fondamentale, possiamo qui dedurre dalle considerazioni del paragrafo precedente una nuova dimostrazione di questo teorema dovuta a Schur.

Supponiamo infatti: 1\textsuperscript{o} che la serie $\infty^r$ di trasformazioni
\begin{equation*}\tag{a}
 x'_i=f_i(x_1,x_2,\ldots,x_n; a_1,a_2,\ldots,a_r)
\end{equation*}
soddisfi ad un sistema di equazioni differenziali della forma
\begin{equation*}\tag{b}
 \frac{\partial x'_i}{\partial a_k}=\sum\limits_s^{1,\ldots,r}\xi_{si}(x')\psi_{sk}(a);
\end{equation*}
2\textsuperscript{o} che nella (a) sia contenuta l'identità, per $a_k=a_k^{(0)}$, e il determinante $|\psi_{ik}(a^{(0)})|$ non sia nullo.

Allora, siccome le (b) formano un sistema illimitatamente integrabile, le trasformazioni infinitesime
\[
 X_jf,\qquad A_jf
\]
soddisfaranno necessariamente le $(\alpha),(\beta)$. Ora cerchiamo di determinare $r$ funzioni incognite $c_1,c_2,\ldots,c_r$ di $r$ variabili $b_1,b_2,\ldots,b_r$ dal sistema di equazioni ai differenziali totali
\begin{equation*}\tag{c}
 \frac{\partial c_\lambda}{\partial b_k}=\sum\limits_s^{1,\ldots,r}\alpha_{\lambda s}(c_1,c_2,\ldots,c_r)\psi_{sk}(b_1,b_2,\ldots,b_r),
 \qquad \lambda,k=1,2,\ldots,r.
\end{equation*}
Questo sistema (c) è illimitatamente integrabile, perchè il corrispondente sistema di equazioni lineari omogenee alle derivate parziali per una funzione incognita $\Theta(c_1,c_2,\ldots,c_r;b_1,b_2,\ldots,b_r)$ si scrive
\[
 -\frac{\partial\Theta}{\partial b_k}+\sum\limits_{\lambda,s}^{1,\ldots,r}\alpha_{\lambda s}(c)\psi_{sk}(b)\frac{\partial\Theta}{\partial c_\lambda}=0,
 \qquad k=1,2,\ldots,r;
\]
moltiplicando la precedente per $\alpha_{jk}(b)$ e sommando da $k=1$ a $k=r$, si ha il sistema equivalente
\[
 \sum\limits_k^{1,\ldots,r}\alpha_{jk}(b)\frac{\partial\Theta}{\partial b_k}+\sum\limits_\lambda^{1,\ldots,r}\alpha_{j\lambda}(c)\frac{\partial\Theta}{\partial c_\lambda}=0.
\]
Ora questo è un sistema completo a causa delle $(\beta)$.

Ciò posto integriamo le (c) sotto le condizioni iniziali
\[
 c_i=a_i\quad\text{per}\quad b_k=a_k^{(0)},
\]
e siano
\begin{equation*}\tag{d}
 c_i=\theta_i(a_1,a_2,\ldots,a_r; b_1,b_2,\ldots,b_r)
\end{equation*}
gl'integrali corrispondenti. Se poniamo
\begin{equation*}\tag{e}
 x_i''=f_i(x_1,x_2,\ldots,x_n;c_1,c_2,\ldots,c_r),
\end{equation*}
avremo per le (b)
\[
 \frac{\partial x_i''}{\partial b_k}
 =\sum\limits_\lambda^{1,\ldots,r} \frac{\partial x_i''}{\partial c_\lambda}\frac{\partial c_\lambda}{\partial b_k}
 =\sum\limits_{\lambda,s,t}^{1,\ldots,r}\xi_{si}(x'')\psi_{s\lambda}(c)\alpha_{\lambda t}(c)\psi_{tk}(b).
\]
Ma
\[
 \sum\limits_\lambda^{1,\ldots,r}\psi_{s\lambda}(c)\alpha_{\lambda t}(c)=\varepsilon_{st}=\begin{cases}0,&s\ne t,\\1,&s=t,\end{cases}
\]
dunque
\[
 \frac{\partial x_i''}{\partial b_k}=\sum\limits_s^{1,\ldots,r}\xi_{si}(x'')\psi_{sk}(b).
\]
Queste sono, per le $x_i''$, considerate come funzioni delle $b$, precisamente le equazioni (b); e poichè per $b_k=a_k^{(0)}$ si riducono, per le (e), a
\[
 f_i(x_1,x_2,\ldots,x_n; a_1,a_2,\ldots,a_r)=x_i',
\]
ne risulta (per il teorema d'unicità)
\[
 x_i''=f_i(x_1',x_2',\ldots,x_n'; b_1,b_2,\ldots,b_r).
\]
Sono dunque soddisfatte le equazioni funzionali
\[
 f_i\bigl(f_1(x,a),f_2(x,a),\ldots,f_n(x,a); b_1,b_2,\ldots,b_r\bigr)
 =f_i(x_1,x_2,\ldots,x_n; c_1,c_2,\ldots,c_r),
\]
caratteristiche di un gruppo.

Si osserverà che in questa dimostrazione di Schur del primo teorema fondamentale le formole che legano i parametri $c$ della trasformazione composta $S_c=S_bS_a$ a quelli $a,b$ dei componenti si ottengono integrando le equazioni ai differenziali totali (c), senza che sia necessario conoscere le equazioni finite (a) del gruppo. Su questa osservazione più tardi tratteremo dei così detti gruppi parametrici.

\section*{\S{} 23. Il secondo teorema fondamentale (teorema principale)}

Affinchè le trasformazioni infinitesime
\[
 X_1f,\ X_2f,\ldots,\ X_rf
\]
linearmente indipendenti generino un gruppo $G_r$ è necessario e sufficiente, come si è visto alla fine del \S{}21, che siano soddisfatte le condizioni (2) e (3). Ma, poichè il gruppo è già perfettamente determinato, osia esiste, dalle espressioni delle sole $X$, è naturale domandarsi se non bastino già le prime condizioni
\begin{equation*}\tag{2}
 (X_i,X_k)=\sum\limits_s^{1,\ldots,r} c_{iks}X_s f
\end{equation*}
per assicurarne l'esistenza; per dedurre cioè che gli $\infty^{r-1}$ gruppi $G_1$ generati dalle combinazioni lineari a coefficienti costanti delle $X$
\[
 \sum\limits_k^{1,\ldots,r} a_kX_kf
\]
costituiscono un gruppo $G_r$. La risposta alla domanda è affermativa, e per dimostrarlo basterà stabilire, secondo il \S{}21, che se le $r$ trasformazioni infinitesime linearmente indipendenti
\begin{equation*}\tag{2*}
 X_kf=\sum\limits_i^{1,\ldots,n}\xi_{ki}(x)\frac{\partial f}{\partial x_i}
\end{equation*}
soddisfano le (2), si possono sempre trovare altre $r$ trasformazioni infinitesime linearmente indipendenti
\[
 A_if=\sum\limits_k^{1,\ldots,r}\alpha_{ik}(a_1,a_2,\ldots,a_r)
        \frac{\partial f}{\partial a_k},
 \qquad i=1,2,\ldots,r,
\]
in $r$ variabili indipendenti $a_1,a_2,\ldots,a_r$, tali che sussistano anche le
\begin{equation*}\tag{3}
 (A_i,A_k)=\sum\limits_s^{1,\ldots,r}c_{iks}A_s f.
\end{equation*}
Per trovare le cercate $A_if$ si può procedere nel modo seguente. Scriviamo $r$ volte le espressioni $X$ in $r$ sistemi di variabili indipendenti
\[
 (x^{(1)}_1,x^{(1)}_2,\ldots,x^{(1)}_n),\quad
 (x^{(2)}_1,x^{(2)}_2,\ldots,x^{(2)}_n),\quad \ldots,
 (x^{(r)}_1,x^{(r)}_2,\ldots,x^{(r)}_n),
\]
ponendo
\[
 X^{(\lambda)}_i\Phi=
 \sum\limits_k^{1,\ldots,n}\xi_{ik}(x^{(\lambda)}_1,x^{(\lambda)}_2,\ldots,x^{(\lambda)}_n)
       \frac{\partial\Phi}{\partial x^{(\lambda)}_k},
 \qquad i,\lambda=1,2,\ldots,r,
\]
e poniamo
\begin{equation*}\tag{8}
 W_i\Phi=X_i^{(1)}\Phi+X_i^{(2)}\Phi+\cdots+X_i^{(r)}\Phi,
\end{equation*}
essendo $\Phi$ una funzione qualunque delle $rn$ variabili $x^{(\lambda)}_k$. Prima di tutto è facile dimostrare che il sistema delle $rn$ equazioni lineari omogenee nelle derivate parziali della $\Phi$
\begin{equation*}\tag{9}
 W_1\Phi=0,
 \quad W_2\Phi=0,
 \quad \ldots,
 \quad W_r\Phi=0
\end{equation*}
è costituito da equazioni tutte indipendenti, cioè non può sussistere un'identità della forma
\[
 \chi_1W_1\Phi+\chi_2W_2\Phi+\cdots+\chi_rW_r\Phi=0,
\]
essendo le $\chi$ convenienti funzioni delle $rn$ variabili $x^{(\lambda)}$. Una tale identità darebbe infatti, per ogni $\lambda$,
\begin{equation*}\tag{10}
 \chi_1\xi_{1k}(x^{(\lambda)})+
 \chi_2\xi_{2k}(x^{(\lambda)})+
 \cdots+
 \chi_r\xi_{rk}(x^{(\lambda)})=0,
 \qquad k=1,2,\ldots,n.
\end{equation*}
Se le supponiamo compatibili, nel primo gruppo di $n$ equazioni ottenuto facendo $\lambda=1$ debbono essere solo $r-\nu$ equazioni indipendenti. Le $n$ equazioni seguenti ottenute prendendo $\lambda=2$ non possono essere tutte conseguenze delle precedenti; altrimenti potremmo soddisfare le $n$ equazioni
\[
 \sum\limits_i^{1,\ldots,r}\chi_i\xi_{ik}(x^{(\lambda)}_1,x^{(\lambda)}_2,\ldots,x^{(\lambda)}_n)=0
\]
coi valori delle $\chi$ indipendenti dalle $x^{(2)}$ e quindi le $r$ espressioni
\[
 X_1^{(2)}\Phi,
 X_2^{(2)}\Phi,
 \ldots,
 X_r^{(2)}\Phi
\]
non sarebbero indipendenti. Dunque il numero delle equazioni del secondo gruppo indipendenti fra loro e dalle precedenti è certo $\nu_1$. Nello stesso modo si vede che ciascuno degli $r$ gruppi fornisce almeno un'equazione indipendente dalle precedenti; si arriva così ad $r$ equazioni indipendenti, onde devono annullarsi tutte le $\chi$.

Ora il sistema di $r$ equazioni (9) è un sistema completo, perchè, avendosi in generale
\[
 (X_i^{(\lambda)},X_k^{(\lambda)})=\sum\limits_s^{1,\ldots,r}c_{iks}X_s^{(\lambda)}\Phi,
\]
si ha anche evidentemente
\[
 (W_i,W_k)=\sum\limits_s^{1,\ldots,r}c_{iks}W_s\Phi.
\]
Esso ammette dunque $rn-r$ soluzioni indipendenti
\[
 u_1,u_2,\ldots,u_{rn-r}.
\]
Facciamo un cambiamento delle variabili indipendenti, prendendo per nuove variabili le
\[
 u_1,u_2,\ldots,u_{rn-r}
\]
insieme con $r$ nuove funzioni indipendenti fra loro e dalle $u$, che indicheremo con
\[
 a_1,a_2,\ldots,a_r.
\]
Se trasformiamo le $W_i$ nelle nuove variabili e indichiamo le operazioni trasformate con $\bar W_i\Phi$, siccome le $u$ sono soluzioni delle (9), avremo
\[
 \bar W_i\Phi=
 W_i(a_1)\frac{\partial\Phi}{\partial a_1}+
 W_i(a_2)\frac{\partial\Phi}{\partial a_2}+\cdots+
 W_i(a_r)\frac{\partial\Phi}{\partial a_r}.
\]
Le $r$ nuove $\bar W_i\Phi$ sono ancora linearmente indipendenti e soddisfano le relazioni
\[
 (\bar W_i,\bar W_k)=\sum\limits_s^{1,\ldots,r}c_{iks}\bar W_s\Phi.
\]
Ma poichè nelle $\bar W_i$ non figurano le derivate della $\Phi$ rispetto alle $u$, è chiaro che la proprietà ora osservata sussiste ancora se, per $u_1,u_2,\ldots,u_{rn-r}$, poniamo delle costanti arbitrarie. Chiamando allora
\[
 A_i\Phi=\sum\limits_k^{1,\ldots,r}\alpha_{ik}(a_1,a_2,\ldots,a_r)\frac{\partial\Phi}{\partial a_k}
\]
le espressioni che così otteniamo dalle $\bar W_i$, avremo evidentemente costruito le espressioni $A_i\Phi$ domandate.

Abbiamo così dimostrato in tutte le due parti il secondo teorema fondamentale che si enuncia:

Ogni gruppo $G_r$ a coppie di trasformazioni inverse è generato da $r$ trasformazioni infinitesime indipendenti
\[
 X_1f,
 X_2f,
 \ldots,
 X_rf,
\]
che soddisfano a relazioni della forma
\begin{equation*}\tag{2}
 (X_i,X_j)=\sum\limits_s^{1,\ldots,r} c_{ijs}X_s f,
 \qquad i,j=1,2,\ldots,r,
\end{equation*}
essendo le $c_{ijs}$ costanti numeriche. Viceversa, se le $r$ trasformazioni infinitesime indipendenti $X_if$ soddisfano relazioni della forma (2), esse generano un tale gruppo $G_r$, il quale è l'insieme di tutti i gruppi $G_1$ a un parametro generati dalle trasformazioni infinitesime
\[
 \sum\limits_k^{1,\ldots,r} a_kX_kf.
\]

\section*{\S{} 24. Le relazioni fra le costanti $c_{iks}$}

Nelle relazioni (2), che legano le $r$ trasformazioni infinitesime
\[
 X_1f,
 X_2f,
 \ldots,
 X_rf
\]
generatrici di un gruppo $G_r$, le quantità $c_{ijs}$ sono $r^3$ costanti numeriche. Ora è assai importante osservare che queste $r^3$ costanti non sono indipendenti, bensì legate da certe relazioni, che ora andiamo a stabilire.

In primo luogo, avendosi identicamente
\[
 (X_i,X_k)+(X_k,X_i)=0,
\]
ne risulta
\[
 \sum\limits_s^{1,\ldots,r}(c_{iks}+c_{kis})X_s f=0,
\]
e poichè le $X_sf$ sono linearmente indipendenti segue intanto
\[
 c_{iks}+c_{kis}=0.
\]
Queste sono le prime relazioni cercate.

In secondo luogo, se $i,k,j$ sono tre indici qualunque, sussiste l'identità di Jacobi
\begin{equation*}\tag{11}
 ((X_i,X_k),X_j)+((X_k,X_j),X_i)+((X_j,X_i),X_k)=0,
\end{equation*}
che si dimostra subito nel modo seguente. Si ha
\[
 ((X_i,X_k),X_j)=X_iX_kX_jf-X_kX_iX_jf-X_jX_iX_kf+X_jX_kX_if,
\]
e siccome il terzo termine diviene dal primo per scambio circolare degli indici, e medesimamente il quarto dal secondo, così è chiaro che, applicando due volte la permutazione ciclica $(ikj)$ e sommando i risultati, si ottiene l'identità Jacobiana (11).

Ora in questa identità poniamo per le espressioni alternate $(X_i,X_k)$ le loro espressioni in funzione lineare delle $X$ date dalle (2). Troviamo dapprima
\[
 ((X_i,X_k),X_j)=
 \left(\sum\limits_s^{1,\ldots,r}c_{iks}X_sf, X_jf\right)
 =\sum\limits_s^{1,\ldots,r}c_{iks}(X_s,X_j)
 =\sum\limits_s^{1,\ldots,r}\sum\limits_a^{1,\ldots,r}c_{iks}c_{sja}X_af,
\]
e per ciò la (11) diventa
\[
 \sum\limits_a^{1,\ldots,r}\left\{\sum\limits_s^{1,\ldots,r}(c_{iks}c_{sja}+c_{kjs}c_{sia}+c_{jis}c_{ska})\right\}X_af=0.
\]
Essendo le $X_af$ indipendenti, segue che le $c$ debbono soddisfare le equazioni
\begin{equation*}\tag{I}
\begin{cases}
 c_{iks}+c_{kis}=0,\\[2mm]
 \displaystyle\sum\limits_s^{1,\ldots,r}
  \left(c_{iks}c_{sja}+c_{kjs}c_{sia}+c_{jis}c_{ska}\right)=0,
\end{cases}
\end{equation*}
per tutti i valori da $1$ a $r$ degli indici $i,k,j,a$.

Abbiamo quindi l'importante risultato: le costanti $c_{iks}$ delle relazioni
\[
 (X_i,X_k)=\sum\limits_s^{1,\ldots,r}c_{iks}X_sf
\]
che legano le $r$ trasformazioni infinitesime $X_1f,X_2f,\ldots,X_rf$ generatrici di un gruppo $G_r$ soddisfano, per tutti i valori degli indici da $1$ a $r$, alle equazioni (I). Si osservi che queste relazioni sono indipendenti affatto dal numero delle variabili su cui opera il gruppo.

Le $r^3$ costanti $c_{iks}$ hanno la massima importanza per la natura del gruppo, e diverranno, per una ragione che apparirà dagli studi ulteriori, le costanti di composizione o strutturali del gruppo.

Affinchè le $r^3$ costanti $c_{iks}$ appartengano come costanti di composizione ad un gruppo $G_r$, è dunque necessario che soddisfino le (I). Ora tali condizioni, come andiamo ora a dimostrare, sono altresì sufficienti: in questa proprietà consiste il terzo teorema fondamentale della teoria.

\section*{\S{} 25. Proprietà del simbolo $\lvert UV\rvert$}

Per dimostrare che le relazioni (I) fra le costanti $c_{iks}$ sono sufficienti perchè esista un gruppo $G_r$ che ammetta le $c_{iks}$ come costanti di composizione, basterà costruire in un conveniente numero $m$ di variabili $r$ trasformazioni infinitesime linearmente indipendenti
\[
 X_1f,
 X_2f,
 \ldots,
 X_rf
\]
tali che risulti
\[
 (X_i,X_k)=\sum\limits_s^{1,\ldots,r}c_{iks}X_sf.
\]
Consideriamo dapprima $r$ variabili indipendenti
\[
 y_1,y_2,\ldots,y_r
\]
ed essendo $U,V$ due funzioni qualunque delle $y$, formiamo l'espressione bilineare nelle derivate delle $U,V$ data da
\[
 \sum\limits_{i,k,s}^{1,\ldots,r}c_{iks}y_s
       \frac{\partial U}{\partial y_i}\frac{\partial V}{\partial y_k},
\]
che indichiamo col simbolo $\lvert UV\rvert$:
\begin{equation*}\tag{12}
 \lvert UV\rvert=
 \sum\limits_{i,k,s}^{1,\ldots,r}c_{iks}y_s
       \frac{\partial U}{\partial y_i}\frac{\partial V}{\partial y_k}.
\end{equation*}
Notiamo che, avendosi
\[
 \lvert U,y_k\rvert=\sum\limits_{i,s}^{1,\ldots,r}c_{iks}y_s\frac{\partial U}{\partial y_i},
\]
possiamo scrivere la (12) anche così:
\begin{equation*}\tag{12*}
 \lvert UV\rvert=\sum\limits_k^{1,\ldots,r}\lvert U,y_k\rvert
                 \frac{\partial V}{\partial y_k}.
\end{equation*}
Ora, per la prima relazione (I) cui soddisfano le $c_{iks}$, vedo subito che il simbolo $\lvert UV\rvert$ possiede la prima proprietà delle espressioni alternate, espressa dalla formola
\begin{equation*}\tag{13}
 \lvert UV\rvert+\lvert VU\rvert=0.
\end{equation*}
In secondo luogo diciamo che se $U,V,W$ sono tre funzioni qualunque delle $y$, sussiste altresì l'identità analoga a quella di Jacobi
\begin{equation*}\tag{14}
 \lvert\lvert U,V\rvert,W\rvert+
 \lvert\lvert V,W\rvert,U\rvert+
 \lvert\lvert W,U\rvert,V\rvert=0.
\end{equation*}
Per questo incominciamo dal dimostrare che se la (14) vale per due date $U,V$ con $W=y_j$, varrà pure colle medesime $U,V$ e $W$ qualunque. E infatti si ha per la (12*)
\[
 \lvert U,W\rvert=\sum\limits_i^{1,\ldots,r}\lvert U,y_i\rvert\frac{\partial W}{\partial y_i}=AW,
 \qquad
 \lvert V,W\rvert=\sum\limits_i^{1,\ldots,r}\lvert V,y_i\rvert\frac{\partial W}{\partial y_i}=BW,
\]
indicando con $AW,BW$ le relative operazioni del primo ordine eseguite su $W$; quindi
\[
\begin{split}
&\lvert U,\lvert V,W\rvert\rvert-
 \lvert V,\lvert U,W\rvert\rvert
 =A(BW)-B(AW)                                                    \\
&\quad=\sum\limits_i^{1,\ldots,r}\{A\lvert V,y_i\rvert-B\lvert U,y_i\rvert\}\frac{\partial W}{\partial y_i}
 =\sum\limits_i^{1,\ldots,r}\{\lvert U,\lvert V,y_i\rvert\rvert-
                 \lvert V,\lvert U,y_i\rvert\rvert\}
       \frac{\partial W}{\partial y_i}.
\end{split}
\]
Ma sussistendo per ipotesi la (14) per $W=y_i$, si ha
\[
 \lvert U,\lvert V,y_i\rvert\rvert-
 \lvert V,\lvert U,y_i\rvert\rvert=
 \lvert\lvert U,V\rvert,y_i\rvert,
\]
e quindi la precedente diviene
\[
 \lvert U,\lvert V,W\rvert\rvert-
 \lvert V,\lvert U,W\rvert\rvert
 =\sum\limits_i^{1,\ldots,r}\lvert\lvert U,V\rvert,y_i\rvert\frac{\partial W}{\partial y_i}
 =\lvert\lvert U,V\rvert,W\rvert,
\]
che è precisamente la (14).

Basta ora osservare che facendo $U=y_i$, $V=y_k$, $W=y_j$, la (14) sussiste perchè si ha
\[
 \lvert y_i,y_k\rvert=\sum\limits_s^{1,\ldots,r}c_{iks}y_s
\]
e quindi
\[
 \lvert\lvert y_i,y_k\rvert,y_j\rvert
 =\sum\limits_s^{1,\ldots,r} c_{iks}\lvert y_s,y_j\rvert
 =\sum\limits_s^{1,\ldots,r} c_{iks}\sum\limits_a^{1,\ldots,r} c_{sja}y_a;
\]
l'identità da verificare segue dunque dalle relazioni quadratiche (I) cui soddisfano le $c_{iks}$.

Ciò posto, siccome la (14) vale per
\[
 U=y_i,
 \quad V=y_k,
 \quad W=y_j,
\]
essa vale anche per
\[
 U=y_i,
 \quad V=y_k,
 \quad W\hbox{ qualunque},
\]
indi per
\[
 U=y_i,
 \quad V\hbox{ qualunque},
 \quad W\hbox{ qualunque},
\]
e infine per $U,V,W$ qualunque. C. D. D.

Se poniamo
\[
 A_if=\lvert y_i,f\rvert\qquad(i=1,2,\ldots,r),
\]
per l'identità (14) avremo
\[
 (A_i,A_k)=(\lvert y_i,f\rvert,\lvert y_k,f\rvert)=
 \sum\limits_s^{1,\ldots,r} c_{iks}\lvert y_s,f\rvert,
\]
cioè
\[
 (A_i,A_k)=\sum\limits_s^{1,\ldots,r} c_{iks}A_s f.
\]
Le $r$ trasformazioni infinitesimali
\[
 A_1f,A_2f,\ldots,A_rf
\]
nelle $r$ variabili $y_1,y_2,\ldots,y_r$ soddisfano dunque alle relazioni di composizione richieste. Se esse risultano anche linearmente indipendenti, lo scopo è già raggiunto.

Ma ciò sempre non accade; lo dimostra l'esempio semplice in cui tutte le $c_{iks}$ siano nulle: allora tutte le $A_if$ si annullano identicamente.

\section*{\S{} 26. Cangiamento di variabili}

Per ottenere in tutti i casi le $r$ trasformazioni richieste linearmente indipendenti, giova innanzi tutto eseguire un dato cambiamento di variabili indipendenti, dalle $y_1,y_2,\ldots,y_r$ nelle $Y_1,Y_2,\ldots,Y_r$, che tutti i simboli $\lvert Y_i,Y_k\rvert$ abbiano il valore zero od uno. Ciò si ottiene nel modo seguente.

In primo luogo, se tutte le $c_{iks}$ sono nulle, basta prendere $Y_i=y_i$; in caso contrario poniamo, per esempio, che non tutte le $c_{1ks}$ siano nulle e prendiamo una soluzione $X_1$ dell'equazione a derivate parziali
\[
 \lvert y_1,f\rvert=\sum\limits_{k,s}^{1,\ldots,r}c_{1ks}y_s\frac{\partial f}{\partial y_k}=1,
\]
e ponendo $P_1=y_1$, la coppia di funzioni $P_1,X_1$ soddisfi alla condizione
\[
 \lvert P_1,X_1\rvert=1.
\]
Consideriamo ora le due equazioni simultanee
\begin{equation*}\tag{15}
 A_1f=\lvert P_1,f\rvert=0,
 \qquad
 A_2f=\lvert X_1,f\rvert=0,
\end{equation*}
le quali sono certo indipendenti perchè $f=P_1$ soddisfa la prima e non la seconda. Esse formano inoltre un sistema completo perchè
\[
 (A_1,A_2)=\lvert P_1,\lvert X_1,f\rvert\rvert-
 \lvert X_1,\lvert P_1,f\rvert\rvert=
 \lvert\lvert P_1,X_1\rvert,f\rvert=
 \lvert1,f\rvert=0;
\]
ammettono dunque $r-2$ soluzioni indipendenti
\[
 y'_1,y'_2,\ldots,y'_{r-2}.
\]
Queste sono anche indipendenti da $P_1,X_1$, perchè supponendole legate da una relazione
\[
 \varphi(P_1,X_1,y'_1,y'_2,\ldots,y'_{r-2})=0
\]
le operazioni $A_1,A_2$ eseguite sulla $\varphi$ darebbero
\[
 \frac{\partial\varphi}{\partial X_1}=0,
 \qquad
 \frac{\partial\varphi}{\partial P_1}=0.
\]
Inoltre, applicando l'identità (14) col fare
\[
 U=y'_i,
 \quad V=y'_k,
 \quad W=\begin{cases}P_1,\\ X_1,
 \end{cases}
\]
vediamo che $\lvert y'_i,y'_k\rvert$ è altresì una soluzione del sistema (15), e quindi è una funzione di $y'_1,y'_2,\ldots,y'_{r-2}$ soltanto. Così le $r$ funzioni indipendenti
\[
 P_1,
 \quad X_1,
 \quad y'_1,y'_2,\ldots,y'_{r-2}
\]
soddisfano alle condizioni
\[
 \lvert P_1,X_1\rvert=1,
 \quad \lvert P_1,y'_i\rvert=0,
 \quad \lvert X_1,y'_i\rvert=0,
 \quad
 \lvert y'_i,y'_k\rvert=\varphi_{ik}(y'_1,y'_2,\ldots,y'_{r-2}).
\]
Se tutte le $\varphi_{ik}$ fossero nulle avremmo già conseguito il richiesto cambiamento di variabili. In caso contrario procedendo sopra
\[
 y'_1,y'_2,\ldots,y'_{r-2}
\]
come prima per $y_1,y_2,\ldots,y_r$, troveremo $r-2$ funzioni indipendenti delle $y'$
\[
 P_2,
 X_2,
 y''_1,y''_2,\ldots,y''_{r-4}
\]
tali che
\[
 \lvert P_2,X_2\rvert=1,
 \quad \lvert P_2,y''_i\rvert=0,
 \quad \lvert X_2,y''_i\rvert=0,
 \quad
 \lvert y''_i,y''_k\rvert=\varphi''_{ik}(y''_1,y''_2,\ldots,y''_{r-4}).
\]
Così procedendo, troveremo $r=2m+q$ funzioni indipendenti di $y_1,y_2,\ldots,y_r$
\[
 P_1,P_2,\ldots,P_m,
 \quad X_1,X_2,\ldots,X_m,
 \quad Y_1,Y_2,\ldots,Y_q
\]
tali che si abbia
\begin{equation*}\tag{16}
 \lvert P_i,P_k\rvert=0,
 \quad \lvert P_i,X_k\rvert=\varepsilon_{ik},
 \quad \lvert X_i,X_k\rvert=0,
 \quad \lvert Y_i,Y_k\rvert=0.
\end{equation*}
Assumendo come nuove variabili le $P$ e le $X$, si raggiunge lo scopo proposto.

\section*{\S{} 27. Il terzo teorema fondamentale}

Esprimendo il simbolo $\lvert UV\rvert$ nelle nuove variabili $P_i,X_i$ avremo
\[
\begin{split}
\lvert UV\rvert={}&
 \sum\limits_{i,k}^{1,\ldots,m+q}\lvert P_i,P_k\rvert
        \frac{\partial U}{\partial P_i}\frac{\partial V}{\partial P_k}
 +\sum\limits_{i,k}^{1,\ldots,m}\lvert X_i,X_k\rvert
        \frac{\partial U}{\partial X_i}\frac{\partial V}{\partial X_k}        \\
&+\sum\limits_i^{1,\ldots,m+q}\sum\limits_k^{1,\ldots,m}
 \lvert P_i,X_k\rvert
 \left(
  \frac{\partial U}{\partial P_i}\frac{\partial V}{\partial X_k}
 -\frac{\partial U}{\partial X_k}\frac{\partial V}{\partial P_i}
 \right).
\end{split}
\]
E quindi, per le (16),
\[
 \lvert UV\rvert=
 \sum\limits_i^{1,\ldots,m}
 \left(
  \frac{\partial U}{\partial P_i}\frac{\partial V}{\partial X_i}
 -\frac{\partial U}{\partial X_i}\frac{\partial V}{\partial P_i}
 \right).
\]
Indicando con
\[
 \Omega_k=\Omega_k(P,X)
\]
le formole di trasformazione delle variabili, le $r$ operazioni
\[
 A_k f=\lvert y_k,f\rvert\qquad (k=1,2,\ldots,r)
\]
costruite al \S{}25, espresse per le nuove variabili, divengono
\begin{equation*}\tag{17}
 A_kf=\sum\limits_i^{1,\ldots,m}
 \left(
  \frac{\partial\Omega_k}{\partial P_i}\frac{\partial f}{\partial X_i}
 -\frac{\partial\Omega_k}{\partial X_i}\frac{\partial f}{\partial P_i}
 \right).
\end{equation*}

Naturalmente anche nelle nuove variabili le $A_i f$ soddisfano alle relazioni
\[
 (A_i,A_k)=\sum\limits_s^{1,\ldots,r} c_{iks}A_s f;
\]
però il numero delle linearmente indipendenti fra le $A_i f$ è rimasto ancora lo stesso e non abbiamo fino ad ora conseguito alcun vantaggio sulla trattazione primitiva. Per altro la forma delle (17) suggerisce un artificio semplice per conseguire lo scopo. Alle $r=2m+q$ variabili
\[
 P_1,P_2,\ldots,P_{m+q};\quad X_1,X_2,\ldots,X_m,
\]
aggiungiamo $q$ nuove variabili indipendenti
\[
 X_{m+1},X_{m+2},\ldots,X_{m+q}
\]
e formiamo
\begin{equation*}\tag{18}
 B_k f=\sum\limits_i^{1,\ldots,m+q}
 \left(
  \frac{\partial\Omega_k}{\partial P_i}\frac{\partial f}{\partial X_i}
 -\frac{\partial\Omega_k}{\partial X_i}\frac{\partial f}{\partial P_i}
 \right),
\end{equation*}
ossia, siccome le $\Omega$ non contengono $X_{m+1},\ldots,X_{m+q}$,
\begin{equation*}\tag{18*}
 B_k f=A_k f+\sum\limits_{\alpha}^{1,\ldots,q}
 \frac{\partial\Omega_k}{\partial P_{m+\alpha}}
 \frac{\partial f}{\partial X_{m+\alpha}}.
\end{equation*}
Dico che le $r$ trasformazioni infinitesime
\[
 B_1f,\ B_2f,\ldots,\ B_rf
\]
nelle $2m+2q$ variabili $P,X$ soddisfano alle relazioni di composizione
\[
 (B_i,B_k)=\sum\limits_s^{1,\ldots,r} c_{iks}B_sf
\]
e sono linearmente indipendenti.

La prima cosa risulta osservando che dalla (18*), non comparendo $X_{m+1},\ldots,X_{m+q}$ nei coefficienti delle $A_k$, si ha
\[
 (B_i,B_k)=(A_i,A_k)+\sum\limits_{\alpha}^{1,\ldots,q}
 \left\{A_i\left(\frac{\partial\Omega_k}{\partial P_{m+\alpha}}\right)
       -A_k\left(\frac{\partial\Omega_i}{\partial P_{m+\alpha}}\right)
 \right\}
 \frac{\partial f}{\partial X_{m+\alpha}}.
\]
Ma si trova subito dalle (17)
\[
\begin{split}
&A_i\left(\frac{\partial\Omega_k}{\partial P_{m+\alpha}}\right)
 -A_k\left(\frac{\partial\Omega_i}{\partial P_{m+\alpha}}\right)       \\
&\quad=\frac{\partial}{\partial P_{m+\alpha}}
 \sum\limits_{\nu}^{1,\ldots,m}
 \left(
  \frac{\partial\Omega_i}{\partial P_\nu}
  \frac{\partial\Omega_k}{\partial X_\nu}
 -\frac{\partial\Omega_i}{\partial X_\nu}
  \frac{\partial\Omega_k}{\partial P_\nu}
 \right)                                                   \\
&\quad=\frac{\partial}{\partial P_{m+\alpha}}\lvert y_i,y_k\rvert
 =\frac{\partial}{\partial P_{m+\alpha}}\sum\limits_s^{1,\ldots,r}c_{iks}\Omega_s
 =\sum\limits_s^{1,\ldots,r} c_{iks}\frac{\partial\Omega_s}{\partial P_{m+\alpha}}.
\end{split}
\]
E quindi
\[
\begin{split}
 (B_i,B_k)&=\sum\limits_s^{1,\ldots,r}c_{iks}A_s f
 +\sum\limits_s^{1,\ldots,r}c_{iks}\sum\limits_{\alpha}^{1,\ldots,q}
 \frac{\partial\Omega_s}{\partial P_{m+\alpha}}
 \frac{\partial f}{\partial X_{m+\alpha}}                 \\
 &=\sum\limits_s^{1,\ldots,r}c_{iks}B_s f,
\end{split}
\]
come si voleva.

Ammesso ora che fra le $B_if$ sussista una relazione
\[
 \sum\limits_i^{1,\ldots,r} e_iB_i f=0
\]
colle $e$ costanti, facendo in questa successivamente
\[
 f=P_1,P_2,\ldots,P_{m+q};\quad X_1,X_2,\ldots,X_{m+q},
\]
ne verrebbe
\[
 \sum\limits_i^{1,\ldots,r} e_i\frac{\partial\Omega_i}{\partial X_\lambda}=0,
 \qquad
 \sum\limits_i^{1,\ldots,r} e_i\frac{\partial\Omega_i}{\partial P_\lambda}=0,
 \qquad \lambda=1,2,\ldots,m+q,
\]
e quindi
\[
 \sum\limits_i^{1,\ldots,r} e_i\Omega_i=\sum\limits_i^{1,\ldots,r} e_iy_i=\hbox{costante},
\]
ciò che, essendo le $y$ indipendenti, dà
\[
 e_1=e_2=\cdots=e_r=0.
\]
Così sono effettivamente costruite le $r$ trasformazioni richieste e possiamo infine formulare il teorema fondamentale della teoria dei gruppi, che si enuncia:

Se un gruppo $G_r$, in un numero qualunque $n$ di variabili $x_1,x_2,\ldots,x_n$, possiede le $r$ trasformazioni infinitesime
\[
 X_i f=\sum\limits_k^{1,\ldots,n}\xi_{ik}(x_1,x_2,\ldots,x_n)\frac{\partial f}{\partial x_k},
\]
le costanti $c_{iks}$ delle relazioni di composizione
\[
 (X_i,X_k)=\sum\limits_s^{1,\ldots,r} c_{iks}X_s f
\]
soddisfano alle relazioni
\[
\left\{
\begin{array}{l}
 c_{iks}+c_{kis}=0,\\[0.4em]
 \displaystyle\sum\limits_s^{1,\ldots,r}
 \left(c_{iks}c_{sja}+c_{kjs}c_{sia}+c_{jis}c_{ska}\right)=0.
\end{array}
\right.
\tag{I}
\]
Viceversa se le $r^3$ costanti $c_{iks}$ soddisfano le (I), si possono sempre trovare, in un numero conveniente di variabili, $r$ trasformazioni infinitesime linearmente indipendenti per le quali le costanti di composizione abbiano i valori $c_{iks}$.

\section*{\S{} 28. Le equazioni di Maurer per le $\psi_{ik}$}

I tre teoremi che abbiamo dimostrato nel presente capitolo sono di fondamentale importanza per la teoria dei gruppi, ed è utile conoscere diverse dimostrazioni. Già al \S{}22 abbiamo esposto una nuova dimostrazione del primo teorema ed ora ci proponiamo di fare altrettanto per gli altri due. Le nuove dimostrazioni sono dovute ancora a Schur e noi le esporremo colle modificazioni e semplificazioni apportatevi da Engel e da Lie stesso.

Dimostriamo perciò che se le $r^3$ costanti $c_{ijs}$ soddisfano le relazioni (I), si possono trovare in $r$ variabili
\[
 a_1,a_2,\ldots,a_r
\]
$r$ trasformazioni linearmente indipendenti
\begin{equation*}\tag{19}
 A_i f=\sum\limits_k^{1,\ldots,r}\alpha'_{ik}(a_1,a_2,\ldots,a_r)
        \frac{\partial f}{\partial a_k}
\end{equation*}
che soddisfino alle relazioni di composizione
\begin{equation*}\tag{20}
 (A_i,A_j)=\sum\limits_s^{1,\ldots,r}c_{ijs}A_sf.
\end{equation*}
Con ciò verrà dimostrato nuovamente il terzo teorema fondamentale ed insieme anche il secondo, come risulta subito dal \S{}21.

Si tratta dunque di trovare $r^2$ funzioni $\alpha'_{ik}$ di $a_1,a_2,\ldots,a_r$, il cui determinante sia $\ne0$, ciò che assicura l'indipendenza lineare delle (19), e che soddisfino alle relazioni (20)
\[
 A_i(\alpha'_{j\mu})-A_j(\alpha'_{i\mu})=\sum\limits_s^{1,\ldots,r} c_{ijs}\alpha'_{s\mu},
\]
ossia alle seguenti
\begin{equation*}\tag{21}
 \sum\limits_k^{1,\ldots,r}
 \left(\alpha'_{ik}\frac{\partial\alpha'_{j\mu}}{\partial a_k}
      -\alpha'_{jk}\frac{\partial\alpha'_{i\mu}}{\partial a_k}\right)
 =\sum\limits_s^{1,\ldots,r} c_{ijs}\alpha'_{s\mu}.
\end{equation*}
In armonia colle notazioni già usate al \S{}19 indichi $\psi_{hk}$ il quoziente del complemento algebrico di $\alpha'_{kh}$ nel determinante delle $\alpha'$ per il determinante stesso. Le equazioni (21) possono cambiarsi in equazioni equivalenti per le $\psi$ nel modo seguente.

Moltiplichiamo la (21) per $\psi_{\nu\mu}$, essendo $i,j,\nu$ indici qualunque da $1$ a $r$, e sommiamo rispetto a $\mu$; otteniamo
\begin{equation*}\tag{22}
 \sum\limits_{k,\mu}^{1,\ldots,r}
 \left(\alpha'_{ik}\psi_{\nu\mu}\frac{\partial\alpha'_{j\mu}}{\partial a_k}
      -\alpha'_{jk}\psi_{\nu\mu}\frac{\partial\alpha'_{i\mu}}{\partial a_k}\right)
 =c_{ij\nu}.
\end{equation*}
Ma si ha
\[
 \sum\limits_\mu^{1,\ldots,r}\psi_{\nu\mu}\alpha'_{j\mu}=\varepsilon_{\nu j},
 \qquad
 \sum\limits_\mu^{1,\ldots,r}\psi_{\nu\mu}\alpha'_{i\mu}=\varepsilon_{\nu i},
\]
e quindi derivando
\[
 \sum\limits_\mu^{1,\ldots,r}\psi_{\nu\mu}\frac{\partial\alpha'_{j\mu}}{\partial a_k}
 =-\sum\limits_\mu^{1,\ldots,r}\alpha'_{j\mu}\frac{\partial\psi_{\nu\mu}}{\partial a_k},
\]
\[
 \sum\limits_\mu^{1,\ldots,r}\psi_{\nu\mu}\frac{\partial\alpha'_{i\mu}}{\partial a_k}
 =-\sum\limits_\mu^{1,\ldots,r}\alpha'_{i\mu}\frac{\partial\psi_{\nu\mu}}{\partial a_k},
\]
onde la (22) diventa
\[
 \sum\limits_{k,\mu}^{1,\ldots,r}\alpha'_{jk}\alpha'_{i\mu}\frac{\partial\psi_{\nu\mu}}{\partial a_k}
 -\sum\limits_{k,\mu}^{1,\ldots,r}\alpha'_{ik}\alpha'_{j\mu}\frac{\partial\psi_{\nu\mu}}{\partial a_k}=c_{ij\nu};
\]
e permutando nella seconda somma gli indici di sommazione $k,\mu$ abbiamo
\[
 \sum\limits_{k,\mu}^{1,\ldots,r}\alpha'_{jk}\alpha'_{i\mu}
 \left(\frac{\partial\psi_{\nu\mu}}{\partial a_k}
      -\frac{\partial\psi_{\nu k}}{\partial a_\mu}\right)=c_{ij\nu}.
\]
In fine moltiplicando quest'ultima per $\psi_{i\lambda}\psi_{j\nu}$ e sommando rispetto a $i,j$, si ottengono le equazioni cercate
\begin{equation*}\tag{II}
 \frac{\partial\psi_{k\lambda}}{\partial a_\nu}
 -\frac{\partial\psi_{k\nu}}{\partial a_\lambda}
 =\sum\limits_{i,j}^{1,\ldots,r}c_{ijk}\psi_{i\lambda}\psi_{j\nu}.
\end{equation*}
Come le (II) si deducono dalle (22), così, invertendo il calcolo ora eseguito, si vede che le (22) seguono dalle (II), alle quali equivalgono dunque perfettamente.

Le equazioni (II) furono stabilite da Maurer. La questione proposta si riduce dunque a trovare $r^2$ funzioni $\psi_{ik}$ delle $a_1,a_2,\ldots,a_r$, a determinante non nullo, e che soddisfino le equazioni (II) di Maurer.

\section*{\S{} 29. Dimostrazione di Schur del secondo e terzo teorema fondamentale}

Con un opportuno artificio si può provare, come Schur dimostrò, un sistema di $r^2$ funzioni $\psi$ che soddisfino le (II), integrando soltanto un sistema di equazioni differenziali lineari ordinarie a coefficienti costanti; ciò che richiede solo, come è ben noto, operazioni algebriche, la risoluzione della così detta equazione caratteristica. Così non soltanto si provano le cercate trasformazioni infinitesime
\[
 A_1f,A_2f,\ldots,A_rf,
\]
ma si dimostra che occorrono per ciò soltanto operazioni algebriche.

A tal scopo operiamo nelle (II) un cangiamento di variabili e di funzioni incognite ponendo
\[
 a_1=\lambda_1t,
 \quad a_2=\lambda_2t,
 \quad \ldots,
 \quad a_r=\lambda_rt,
\]
\[
 \psi_{ik}(a_1,a_2,\ldots,a_r)=\theta_{ik}(\lambda_1,\lambda_2,\ldots,\lambda_r;t);
\]
con ciò le (II)
\[
 \frac{\partial\psi_{s\alpha}}{\partial a_\beta}
 -\frac{\partial\psi_{s\beta}}{\partial a_\alpha}
 =\sum\limits_{i,j}^{1,\ldots,r}c_{ijs}\psi_{i\alpha}\psi_{j\beta}
\]
diventano
\begin{equation*}\tag{II*}
 \frac{\partial\theta_{s\alpha}}{\partial\lambda_\beta}
 -\frac{\partial\theta_{s\beta}}{\partial\lambda_\alpha}
 =\sum\limits_{i,j}^{1,\ldots,r}c_{ijs}\theta_{i\alpha}\theta_{j\beta}.
\end{equation*}
Ora si integri il sistema di equazioni differenziali lineari ordinarie rispetto a $t$
\begin{equation*}\tag{23}
 \frac{\partial\theta_{s\alpha}}{\partial t}
 =\varepsilon_{s\alpha}+\sum\limits_{i,j}^{1,\ldots,r}c_{ijs}\lambda_j\theta_{i\alpha},
 \qquad s,\alpha=1,2,\ldots,r,
\end{equation*}
colle condizioni iniziali
\begin{equation*}\tag{24}
 \theta_{s\alpha}=0\qquad\hbox{per }t=0,
\end{equation*}
dove nella integrazione si riguarderanno le $\lambda$ come costanti (parametri). Dalla forma di queste equazioni differenziali lineari a coefficienti costanti risulta subito che le funzioni $\theta_{s\alpha}$ conterranno le variabili $\lambda_1,\lambda_2,\ldots,\lambda_r,t$ appunto nei prodotti
\[
 \lambda_1t,\quad \lambda_2t,\quad \ldots,
 \quad \lambda_rt;
\]
ed ora dimostreremo che le $\theta_{s\alpha}$ così determinate soddisferanno le (II*), e per ciò le $\psi_{s\alpha}$ soddisferanno le (II) come si voleva.

Pongasi infatti
\begin{equation*}\tag{25}
 V_{s\alpha\beta}=\frac{\partial\theta_{s\alpha}}{\partial\lambda_\beta}
 -\frac{\partial\theta_{s\beta}}{\partial\lambda_\alpha}
 -\sum\limits_{i,j}^{1,\ldots,r}c_{ijs}\theta_{i\alpha}\theta_{j\beta}.
\end{equation*}
E si derivi rispetto a $t$ osservando le (23); si avrà
\[
\begin{split}
\frac{\partial V_{s\alpha\beta}}{\partial t}={}&
\frac{\partial}{\partial\lambda_\beta}
 \left[\varepsilon_{s\alpha}+\sum\limits_{i,j}^{1,\ldots,r}c_{ijs}\lambda_j\theta_{i\alpha}\right]
-\frac{\partial}{\partial\lambda_\alpha}
 \left[\varepsilon_{s\beta}+\sum\limits_{i,j}^{1,\ldots,r}c_{ijs}\lambda_j\theta_{i\beta}\right]                 \\
&-\sum\limits_{i,j}^{1,\ldots,r}c_{ijs}\theta_{j\beta}
 \left[\varepsilon_{i\alpha}+\sum\limits_{k,v}^{1,\ldots,r}c_{kvi}\lambda_v\theta_{k\alpha}\right]
-\sum\limits_{i,j}^{1,\ldots,r}c_{ijs}\theta_{i\alpha}
 \left[\varepsilon_{j\beta}+\sum\limits_{k,v}^{1,\ldots,r}c_{kvj}\lambda_v\theta_{k\beta}\right].
\end{split}
\]
Ossia
\[
\begin{split}
\frac{\partial V_{s\alpha\beta}}{\partial t}={}&
 \sum\limits_i^{1,\ldots,r} c_{i\beta s}\theta_{i\alpha}
-\sum\limits_i^{1,\ldots,r} c_{i\alpha s}\theta_{i\beta}
+\sum\limits_{i,j}^{1,\ldots,r}c_{ijs}\lambda_j
 \left(\frac{\partial\theta_{i\alpha}}{\partial\lambda_\beta}
      -\frac{\partial\theta_{i\beta}}{\partial\lambda_\alpha}\right)                         \\
&-\sum\limits_{i,j}^{1,\ldots,r}c_{ijs}\theta_{j\beta}\varepsilon_{i\alpha}
-\sum\limits_{i,j}^{1,\ldots,r}c_{ijs}\theta_{i\alpha}\varepsilon_{j\beta} \\
&-\sum\limits_{i,j,k,v}^{1,\ldots,r}c_{ijs}c_{kvi}\lambda_v\theta_{j\beta}\theta_{k\alpha}
-\sum\limits_{i,j,k,v}^{1,\ldots,r}c_{ijs}c_{kvj}\lambda_v\theta_{i\alpha}\theta_{k\beta}.
\end{split}
\]
Le somme semplici si distruggono per la prima delle relazioni (I) per le $c$ e, ponendo per
\[
 \frac{\partial\theta_{i\alpha}}{\partial\lambda_\beta}
 -\frac{\partial\theta_{i\beta}}{\partial\lambda_\alpha}
\]
l'espressione data dalla (25), resta
\[
\begin{split}
\frac{\partial V_{s\alpha\beta}}{\partial t}={}&
 \sum\limits_{i,j}^{1,\ldots,r}c_{ijs}\lambda_jV_{i\alpha\beta}
-\sum\limits_{i,j,k,v}^{1,\ldots,r}c_{ijs}c_{kvi}\lambda_v\theta_{j\beta}\theta_{k\alpha}  \\
&-\sum\limits_{i,j,k,v}^{1,\ldots,r}c_{ijs}c_{kvj}\lambda_v\theta_{i\alpha}\theta_{k\beta}.
\end{split}
\]
E permutando nella seconda delle somme quadruple gli indici $j,v$, e nell'ultima eseguendo sugli indici la sostituzione ciclica $(ikv)$, si vede che le due somme quadruple si distruggono. Dunque le $V_{s\alpha\beta}$ definite dalla (25) soddisfano alle equazioni lineari omogenee
\[
 \frac{\partial V_{s\alpha\beta}}{\partial t}
 =\sum\limits_{i,j}^{1,\ldots,r}c_{ijs}\lambda_jV_{i\alpha\beta}.
\]
D'altronde per $t=0$ si annullano, a causa delle condizioni iniziali (24), e per conseguenza sono sempre nulle; ciò che dimostra sussistere le (II*). Così dunque abbiamo dimostrato:

Dato un sistema di $r^3$ costanti $c_{iks}$, che verifichino le (I), si possono trovare, con sole operazioni algebriche, $r$ trasformazioni infinitesime indipendenti
\[
 A_i f=\sum\limits_k^{1,\ldots,r}\alpha_{ik}(a_1,a_2,\ldots,a_r)\frac{\partial f}{\partial a_k}
\]
in $r$ variabili $a_1,a_2,\ldots,a_r$ che soddisfino alle relazioni di composizione
\[
 (A_i,A_k)=\sum\limits_s^{1,\ldots,r} c_{iks}A_s f
\]
e generino quindi un gruppo $G_r$.

\section*{\S{} 30. Composizione di due serie $\infty^r$ di trasformazioni}

Coi tre teoremi di cui ci siamo occupati nel presente capitolo abbiamo posti i fondamenti per la teoria generale dei gruppi continui finiti $G_r$ a coppie di trasformazioni inverse, le cui equazioni finite sono contenute in un'unica serie $\infty^r$
\[
 x'_i=f_i(x_1,x_2,\ldots,x_n;a_1,a_2,\ldots,a_r).
\]
Già al \S{}4 abbiamo indicato che si possono considerare gruppi più generali a $r$ parametri le cui trasformazioni, in luogo di formare una sola serie $\infty^r$, costituiscono più di tali serie od anche un numero infinito ma discreto. Ora vogliamo dimostrare quanto fu allora asserito, che lo studio di tali gruppi misti si riporta ai gruppi $G_r$ combinati con ordinarii gruppi discontinui.

Per questo converrà che riprendiamo, estendendole, le considerazioni dei \S\S{}12--14, ed innanzi tutto ci occupiamo della questione seguente.

Date due serie, ciascuna $\infty^r$, di trasformazioni
\begin{align*}
 x'_i&=f_i(x_1,x_2,\ldots,x_n;a_1,a_2,\ldots,a_r), \tag{26}\\
 x''_i&=\varphi_i(x'_1,x'_2,\ldots,x'_n;b_1,b_2,\ldots,b_r), \tag{27}
\end{align*}
se componiamo una qualunque trasformazione $S_a$ della prima con una qualunque $T_b$ della seconda, nella trasformazione composta $S_aT_b$
\[
 (S_aT_b)\quad
 x''_i=\varphi_i(f_1(x,a),\ldots,f_n(x,a);b_1,\ldots,b_r)
\]
il numero dei parametri essenziali è al massimo $2r$ e al minimo $r$\begingroup\renewcommand\thefootnote{(*)}\footnotemark\endgroup.\begingroup\renewcommand\thefootnote{(*)}\footnotetext{Ciò è evidente perchè, tenendo per es. fisse le $b$ e facendo variare le $a$, la $S_aT_b$ percorre $\infty^r$ trasformazioni.}\endgroup{} Domandiamo: quali relazioni debbono passare fra le due serie $\infty^r$ $S_a,T_b$ affinchè la serie composta $S_aT_b$ abbia ancora soltanto $r$ parametri essenziali?

Indicando con $c_1,c_2,\ldots,c_r$ gli $r$ parametri essenziali della serie composta potremo scrivere
\begin{equation*}\tag{28}
 x''_i=\psi_i(x_1,x_2,\ldots,x_n;c_1,c_2,\ldots,c_r),
\end{equation*}
dove i parametri $c$ sono funzioni delle $a$ e delle $b$
\begin{equation*}\tag{29}
 c_\eta=\theta_\eta(a_1,\ldots,a_r;b_1,\ldots,b_r).
\end{equation*}
Dalla osservazione stessa fatta al principio del \S{}13 risulta che queste relazioni (28) sono risolubili sia rispetto alle $a$ che alle $b$, cioè i determinanti funzionali
\[
 \frac{\partial(\theta_1,\theta_2,\ldots,\theta_r)}{\partial(a_1,a_2,\ldots,a_r)},
 \qquad
 \frac{\partial(\theta_1,\theta_2,\ldots,\theta_r)}{\partial(b_1,b_2,\ldots,b_r)}
\]
saranno ciascuno diverso da zero. Ora diamo alle $a,b$ tali incrementi infinitesimi $\delta a,\delta b$ da soddisfare alle condizioni $\delta c_\eta=0$, cioè
\begin{equation*}\tag{30}
 \sum\limits_k^{1,\ldots,r}\frac{\partial\theta_\eta}{\partial a_k}\delta a_k
 +\sum\limits_k^{1,\ldots,r}\frac{\partial\theta_\eta}{\partial b_k}\delta b_k=0,
\end{equation*}
nella qual cosa, per quanto è detto sopra, si potranno prendere ad arbitrio le $\delta a$ e calcolare le $\delta b$ inversamente. Precisamente come al \S{}13 ne seguirà
\begin{equation*}\tag{31}
 S_a^{-1}S_{a+\delta a}=T_bT_{b+\delta b}^{-1}.
\end{equation*}
Fissando ad arbitrio le $a$ e facendo variare comunque le $\delta a$, il primo membro della (31) percorre una serie $\infty^{r-1}$ di trasformazioni infinitesime; lo stesso dicasi del secondo membro ove si fissino le $b$ e si facciano variare ad arbitrio le $\delta b$.

L'eguaglianza (31) esprime quindi (cfr. \S{}13) che queste due serie $\infty^{r-1}$ di trasformazioni infinitesime rimangono sempre le stesse qualunque siano i valori attribuiti alle $a$ o alle $b$.

Siccome adunque le trasformazioni della serie (26) infinitamente vicine alla $S_a$ si ottengono combinando la $S_a$ con una trasformazione infinitesima della detta serie fissa, è applicabile ancora il ragionamento del \S{}14 e ne risulta che la serie (26) soddisfa ad un sistema di equazioni differenziali del tipo
\[
 \frac{\partial x'_i}{\partial a_k}=\sum\limits_j^{1,\ldots,r}\eta_{jk}(a)\xi_{ji}(x'),
\]
per le quali il determinante delle $\eta$ non è identicamente nullo. Se diamo quindi alle $a$ un particolare sistema di valori $a^{(0)}$, tale soltanto che sia
\[
 \lvert\eta_{jk}(a^{(0)})\rvert\ne0,
\]
segue dal teorema alla fine del \S{}16 che le trasformazioni
\[
 S_{a^{(0)}}^{-1}S_a
\]
formano il gruppo $G_r$, a coppie di trasformazioni inverse, generato dalle $r$ trasformazioni infinitesime
\[
 X_j f=\sum\limits_i^{1,\ldots,n}\xi_{ji}(x_1,x_2,\ldots,x_n)\frac{\partial f}{\partial x_i}.
\]
D'altronde se fissiamo una trasformazione qualunque $T_{b^{(0)}}$ della serie (27), possiamo soddisfare alla relazione
\[
 S_aT_b=S_{a^{(0)}}T_{b^{(0)}},
\]
cioè alle
\begin{equation*}\tag{32}
 \theta_\lambda(a,b)=\theta_\lambda(a^{(0)},b^{(0)})
 \qquad (\lambda=1,2,\ldots,r),
\end{equation*}
sia prendendo ad arbitrio $S_a$ e calcolandone $T_b$ inversamente. Se scriviamo la (32) così:
\[
 S_{a^{(0)}}^{-1}S_a=T_{b^{(0)}}T_b^{-1},
\]
vediamo che $T_{b^{(0)}}T_b^{-1}$ percorre le stesse trasformazioni come $S_{a^{(0)}}^{-1}S_a$, cioè le trasformazioni del gruppo $G_r$. Indicando queste ultime col simbolo $E_\lambda$, avremo
\[
 S_a=S_{a^{(0)}}E_\lambda,
\]
e poichè l'inversa $T_bT_{b^{(0)}}^{-1}$ di $T_{b^{(0)}}T_b^{-1}$ è pure una $E$,
\[
 T_b=E_\mu T_{b^{(0)}}.
\]
Viceversa se $E_\lambda$ percorre le trasformazioni di un gruppo $G_r$ ed $S_{a^{(0)}},T_{b^{(0)}}$ sono due trasformazioni fisse arbitrarie, è chiaro che le due serie $\infty^r$
\[
 S_{a^{(0)}}E_\lambda,
 \qquad
 E_\lambda T_{b^{(0)}}
\]
godono della voluta proprietà, perchè
\[
 (S_{a^{(0)}}E_\lambda)(E_\mu T_{b^{(0)}})
 =S_{a^{(0)}}(E_\lambda E_\mu)T_{b^{(0)}}
 =S_{a^{(0)}}E_\nu T_{b^{(0)}}
\]
contiene i soli $r$ parametri della trasformazione composta $E_\nu$. Abbiamo dunque stabilito il teorema:

Affinchè due serie $\infty^r$ di trasformazioni $S_a,T_b$ diano luogo ad una serie composta $S_aT_b$ con soli $r$ parametri essenziali, è necessario e sufficiente che esista un gruppo $G_r$, a coppie di trasformazioni inverse, tali che, indicando con $E_\lambda$ una trasformazione variabile in $G_r$ e con $A,B$ due trasformazioni fisse, l'una della prima, l'altra della seconda serie, ogni trasformazione della prima serie sia data da
\[
 S_a=AE_\lambda
\]
ed ognuna della seconda da
\[
 T_b=E_\mu B.
\]
Se supponiamo in particolare che le due serie $\infty^r$ coincidano, si dovrà avere
\[
 S_{a^{(0)}}E_\lambda=E_\mu S_{a^{(0)}},
\]
cioè
\[
 S_{a^{(0)}}^{-1}E_\mu S_{a^{(0)}}=E_\lambda,
\]
onde abbiamo il teorema:

Se una serie $\infty^r$ di trasformazioni $S_a$ gode della proprietà che la trasformazione composta di due qualunque della serie abbia solo $r$ parametri essenziali, la serie stessa ha la forma
\[
 S_a\equiv AE_a,
\]
dove $E_a$ percorre le trasformazioni di un gruppo $G_r$ ed $A$ è una trasformazione fissa qualunque nella serie, che deve soddisfare alla condizione
\[
 A^{-1}G_rA=G_r,
\]
cioè sarà permutabile col gruppo.

\section*{\S{} 31. Serie $\infty^r$ di trasformazioni nelle quali i prodotti di $m$ trasformazioni sono nella serie}

Passiamo ora a risolvere la questione seguente. Quali condizioni deve verificare una serie $\infty^r$ di trasformazioni perchè componendo $m$ trasformazioni qualunque della serie nasca sempre una trasformazione della serie stessa, nè ciò accada componendone un numero minore?

Osserviamo in primo luogo che la trasformazione composta di due della serie dovrà contenere solo $r$ parametri, perchè se ne contenesse di più questi non potrebbero più sparire componendovi le successive. Applicando il secondo teorema del paragrafo precedente, vediamo intanto che le trasformazioni della nostra serie avranno certo la forma
\[
 S_a=AE_a
\]
del detto teorema. Ora componendo due qualunque trasformazioni della serie
\[
 AE_\lambda,
 \qquad
 AE_\mu,
\]
essendo $AE_\mu A^{-1}=E_\nu$ e $AE_\mu=E_\nu A$, si ha
\[
 (AE_\lambda)(AE_\mu)=A^2E_\nu.
\]
In generale componendo $s$ trasformazioni della serie si avrà una trasformazione della forma
\[
 A^sE_a.
\]
La nostra ipotesi porta adunque che si abbia
\[
 A^mE_\lambda=AE_\mu,
\]
cioè
\[
 A^{m-1}=E_\nu,
\]
vale a dire che la potenza $(m-1)^{\mathrm{ma}}$ di $A$ appartenga a $G_r$, ma nessuna altra potenza d'ordine inferiore. Concludiamo che: Ogni serie della specie voluta si ottiene colle sostituzioni $E_\lambda$ di un gruppo $G_r$ a coppie di trasformazioni inverse, una trasformazione fissa $A$ permutabile col gruppo $G_r$, e tale che la potenza $A^{m-1}$ di $A$, ma

nessuna potenza d'ordine inferiore, appartenga a $G_r$.

È necessario ricordare che per l'esatta applicazione di questo teorema conviene limitarsi a considerare quelle trasformazioni della serie che giacciono in un intorno convenientemente ristretto della $A$, poichè la validità del teorema alla fine del \S{}16 sul quale ci siamo appoggiati è soggetto appunto ad una tale limitazione.

Se si applica il teorema superiore al caso $m=2$, cioè nell'ipotesi che la serie $S_a$ costituisca un gruppo $\Gamma$, senza supporre che in $\Gamma$ vi sia l'identità nè coppie di trasformazioni inverse, vediamo che in tal caso le trasformazioni di $\Gamma$ appartengono ad un gruppo $G_r$ a coppie di trasformazioni inverse. Se ne conclude che le equazioni di qualsiasi gruppo finito a $r$ parametri si ottengono da quelle di un gruppo $G_r$ a coppie di trasformazioni inverse con un cangiamento di parametri e per prolungamento analitico. Così apparisce l'opportunità, nella teoria dei gruppi continui finiti, di limitarsi alla considerazione dei gruppi $G_r$ a coppie di trasformazioni inverse.

\section*{\S{} 32. I gruppi misti}

Supponiamo ora di avere un numero dapprima finito, diciamo $m$, di serie di trasformazioni contenenti ciascuna $r$ parametri
\[
 x'_i=f_i^{(k)}(x_1,x_2,\ldots,x_n; a_1^{(k)},a_2^{(k)},\ldots,a_r^{(k)}),
 \qquad
 \begin{gathered}
 i=1,2,\ldots,n,\\
 k=1,2,\ldots,m,
 \end{gathered}
\]
e che queste $m$ serie, insieme considerate, formino un gruppo misto. Considerando ciascuna di queste serie per sè e paragonandole due a due saranno applicabili i teoremi del \S{}30. Quindi, indicando con $S^1,S^2,\ldots,S^m$ $m$ trasformazioni fisse prese nelle rispettive serie e con
\[
 S_{a^{(1)}}^1,
 \quad S_{a^{(2)}}^2,
 \quad \ldots,
 \quad S_{a^{(m)}}^m
\]
trasformazioni variabili nelle serie stesse, risulta che le $m$ serie
\[
 (S^1)^{-1}S_{a^{(1)}}^1,
 \quad (S^2)^{-1}S_{a^{(2)}}^2,
 \quad \ldots,
 \quad (S^m)^{-1}S_{a^{(m)}}^m
\]
coincidono in un unico gruppo $G_r$ a coppie di trasformazioni inverse, gruppo permutabile con tutte le trasformazioni delle serie. Indicando con $E_a$ una trasformazione variabile in $G_r$, le $m$ serie saranno date dalle trasformazioni
\begin{equation*}\tag{33}
 S^1E_a,
 \quad S^2E_a,
 \quad \ldots,
 \quad S^mE_a.
\end{equation*}
Ora se si moltiplicano tutte queste trasformazioni per una qualunque di esse si riprodurranno solamente trasformazioni formanti un gruppo; ma, di più, essendo le $S$ permutabili col gruppo $G_r$, le $m$ serie non faranno che permutarsi fra loro. Inverso, il prodotto $S_iS_k$, appartenendo alla serie stessa, ha necessariamente la forma
\[
 S_iS_k=S_\ell E_{\lambda'},
\]
e si ha
\[
\begin{split}
 (S_iE_\lambda)(S_kE_\mu)
  &=S_iE_\lambda S_kE_\mu
   =S_iS_kE_\lambda E_\mu       \\
  &=S_\ell E_{\lambda'}E_\lambda E_\mu
   =S_\ell E_\alpha.
\end{split}
\]
Dunque le trasformazioni (33) formano un gruppo. Il gruppo misto in discorso dipende così da un gruppo continuo $G_r$ e da $m$ trasformazioni fisse $S_1,S_2,\ldots,S_m$ permutabili con $G_r$ e tali che il prodotto di due qualunque di esse sia uguale ad una terza, moltiplicata per una trasformazione del gruppo $G_r$.

Gli elementi dalla formazione di un gruppo misto dipendono dai gruppi $G_r$, avendo in primo luogo le trasformazioni $S_i$ permutabili con $G_r$, del quale problema ci occuperemo in seguito. In secondo luogo bisogna scegliere $m$ di tali trasformazioni $S_1,S_2,\ldots,S_m$ in guisa che rispetto al gruppo $G_r$ diano luogo ad un gruppo finito ordinario. Se il numero delle serie $\infty^r$ è composto, il gruppo misto è soltanto discreto ma infinito; il risultato è del tutto analogo, ma allora le trasformazioni moltiplicatrici
\[
 S_1,
 \quad S_2,
 \quad S_3,
 \quad \ldots
\]
sono in numero infinito e costituiscono un gruppo discontinuo. In ogni caso vediamo, come si è asserito al \S{}4, che questi gruppi misti partecipano insieme della natura dei gruppi continui $G_r$ e di quella dei gruppi discontinui.

Apparentemente i gruppi misti che abbiamo qui considerati sono di natura speciale perchè le serie che li compongono hanno lo stesso numero $r$ di parametri; ma in realtà vediamo facilmente che a questo riduce il caso più generale. Supponiamo infatti che le $m$ serie componenti il gruppo abbiano rispettivamente
\[
 r_1,r_2,\ldots,r_m
\]
parametri e sia $r$ il massimo dei numeri $r_i$. Combinando una serie con $r_i$ parametri con un'altra a $r_j$ parametri si ottiene un'unica serie il numero dei cui parametri essenziali è al massimo $r_i+r_j$ e, supposto $r_i\ge r_j$, al minimo $r_i$. Perciò quelle fra le $m$ serie che hanno il numero massimo $r$ di parametri costituiscono già di per sè un gruppo misto $\Gamma$ della specie sopra considerata, ed in particolare vi è in $\Gamma$ un sottogruppo $G_r$ (una particolare serie)
\[
 \text{a)}\qquad x'_i=f_i(x_1,x_2,\ldots,x_n; a_1,a_2,\ldots,a_r)
\]
contenente l'identità. Ma allora se noi consideriamo le sostituzioni di un'altra delle serie con $r_i<r$ parametri,
\[
 \text{b)}\qquad x''_i=\varphi_i(x'_1,x'_2,\ldots,x'_n; b_1,b_2,\ldots,b_{r_i}),
\]
la serie composta ha $r$ parametri ed è quindi in $\Gamma$. Ma in questa serie è contenuta completamente la b), perchè in a) vi è l'identità. Si può dunque senz'altro sopprimere la serie b) come già contenuta in un'altra di $\Gamma$.

Dalle ricerche ultime risulta come lo studio dei gruppi continui finiti si riduca in sostanza unicamente a quello dei gruppi $G_r$ generati da $r$ trasformazioni infinitesime
\[
 X_1f,
 \quad X_2f,
 \quad \ldots,
 \quad X_rf,
\]
e contenenti l'identità e con ogni trasformazione la sua inversa. In seguito ci occuperemo quindi ormai solo di questi gruppi $G_r$, sui quali si concentra tutto l'interesse.

\chapter*{Capitolo III}
\addcontentsline{toc}{chapter}{Capitolo III}

\begin{center}
\begin{minipage}{0.88\linewidth}
\centering
Equazioni di definizione di un gruppo $G_r$.\par
Ordine delle trasformazioni infinitesime.
\end{minipage}
\end{center}

\section*{\S{} 33. Eliminazione di costanti arbitrarie da un sistema di funzioni}

Ogni gruppo $G_r$ è perfettamente definito dalle sue $r$ trasformazioni infinitesime
\[
 X_k f=\sum\limits_i^{1,\ldots,n}\xi_{ki}(x_1,x_2,\ldots,x_n)\frac{\partial f}{\partial x_i},
\]
soddisfacenti alle condizioni di composizione
\[
 (X_kX_l)=\sum\limits_s^{1,\ldots,r}c_{kls}X_sf.
\]
La più generale trasformazione infinitesima di un tale $G_r$ è data da
\[
 Xf=\sum\limits_k^{1,\ldots,r}e_kX_kf
   =\sum\limits_k^{1,\ldots,r}\sum\limits_i^{1,\ldots,n}
     e_k\xi_{ki}\frac{\partial f}{\partial x_i},
\]
dove $e_1,e_2,\ldots,e_r$ sono $r$ costanti arbitrarie. I coefficienti $\xi_i$ di una tal trasformazione $Xf$ essendo dati da
\[
 \xi_i=\sum\limits_k^{1,\ldots,r}e_k\xi_{ki}
 \qquad (i=1,2,\ldots,n),
\]
contengono (linearmente ed omogeneamente) le $r$ costanti arbitrarie $e_i$.

Fissate le $\xi_i$, è individuato il gruppo $G_r$ e quindi le espressioni delle $\xi_i$. Ora invece di dare esplicitamente le $\xi_i$, si possono anche definire come le soluzioni più generali di un certo sistema di equazioni simultanee alle derivate parziali, che hanno grande importanza in molte ricerche e diconsi le equazioni di definizione del gruppo.

Per arrivare a costruire il sistema in discorso dobbiamo premettere alcune osservazioni generali sulla eliminazione di costanti da un gruppo di funzioni di un numero qualunque di variabili indipendenti.

Siano $z_1,z_2,\ldots,z_m$ $m$ funzioni di $n$ variabili indipendenti $x_1,x_2,\ldots,x_n$, contenenti inoltre $r$ parametri (o costanti) $a_1,a_2,\ldots,a_r$, che potremo supporre e supporremo essenziali, sia
\begin{equation*}\tag{1}
 z_i=f_i(x_1,x_2,\ldots,x_n; a_1,a_2,\ldots,a_r),
 \qquad i=1,2,\ldots,m.
\end{equation*}
Dovendo considerare nel seguito le derivate dei varî ordini delle $z$, è opportuno introdurre una notazione abbreviata per indicarle, e noi porremo
\[
 \frac{\partial z_i}{\partial x_k}=z'_{i;k},
 \qquad
 \frac{\partial^2 z_i}{\partial x_k\partial x_h}=z'_{i;kh},
 \qquad \ldots,
\]
in generale
\[
 \frac{\partial^s z_i}{\partial x_{k_1}\partial x_{k_2}\cdots\partial x_{k_s}}
   =z^{(s)}_{i;k_1k_2\ldots k_s}.
\]
Ciò posto aggiungiamo alle (1) le equazioni che ne risultano con una prima derivazione
\begin{equation*}\tag{$1_1$}
 z'_{i;k}=f'_{i;k}(x,a)
 \qquad
 \begin{gathered}
 i=1,2,\ldots,m,\\
 k=1,2,\ldots,n,
 \end{gathered}
\end{equation*}
indi quelle che se ne ottengono con una seconda derivazione
\begin{equation*}\tag{$1_2$}
 z'_{i;k_1k_2}=f'_{i;k_1k_2}(x,a),
 \qquad \ldots,
\end{equation*}
e così continuiamo, protraendo la derivazione finchè sia possibile l'eliminazione degli $r$ parametri del sistema delle equazioni
\[
 (1),(1_1),(1_2),\ldots,
\]
ciò che necessariamente dovrà accadere, come ora andiamo a precisare. Riprendendo a tale scopo le notazioni del \S{}3, indichiamo con $\mu_0$ la caratteristica della matrice funzionale delle (1), cioè delle $f_i$ rispetto ai parametri, con $\mu_1$ quella del sistema $(1),(1_1)$, ecc., sia delle funzioni $f_i,f_{i;k},f_{i;kh}$ e così di seguito. Per quanto abbiamo visto al \S{}3, i numeri
\[
 \mu_0,\mu_1,\mu_2,\ldots,
\]
che sono tutti $\le r$, formano una serie crescente fino al massimo $=r$ (perchè gli $r$ parametri $a$ sono per ipotesi essenziali); e da quel punto in poi la serie delle $\mu$ è stazionaria. Indichiamo con $\mu_{s-1}$ la prima $\mu=r$ ed avremo
\[
 \mu_0<\mu_1<\cdots<\mu_{s-2}<\mu_{s-1}=r.
\]

Il sistema di equazioni
\[
 (1),(1_1),(1_2),\ldots,(1_{s-1})
\]
avendo la matrice funzionale dei secondi membri rispetto agli $r$ parametri $a$ appunto la caratteristica $r$, è certo risolubile rispetto a questi parametri $a$, che risulteranno così espressi per le $x$, le $z$ e le loro derivate fino all'ordine $s-1$. Questi valori delle $a$ sostituiti nelle equazioni stesse $(1),(1_1),\ldots,(1_{s-1})$ danno luogo ad un numero $q$ di equazioni in termini finiti fra le $x$, le $z$ e le loro derivate fino all'ordine $s-1$:
\begin{equation*}\tag{$\alpha$}
 W_j(x,z,z'_{i;k},\ldots,z^{(s-1)}_{i;k_1\ldots k_{s-1}})=0,
 \qquad j=1,2,\ldots,q.
\end{equation*}
Questo numero $q$ eguaglia la differenza fra il numero totale delle equazioni $(1),(1_1),\ldots,(1_{s-1})$ ed il numero $r$. Se indichiamo in generale con $\varepsilon_k$ il numero delle derivate delle $z$ d'ordine $k$\begingroup\renewcommand\thefootnote{(*)}\footnotemark\endgroup{} e per simmetria poniamo $\varepsilon_0=m$, avremo
\begingroup\renewcommand\thefootnote{(*)}\footnotetext{Evidentemente è $\varepsilon_k=m\dfrac{n(n+1)\cdots(n+k-1)}{k!}$.}\endgroup{}
\begin{equation*}\tag{2}
 q=\varepsilon_0+\varepsilon_1+\cdots+\varepsilon_{s-1}-r.
\end{equation*}

Consideriamo altresì le successive equazioni $(1_s)$, cioè
\[
 z^{(s)}_{i;kk_1\ldots k_s}=f^{(s)}_{i;kk_1\ldots k_s}(x,a),
\]
e poniamo al posto delle $a$ le loro espressioni per le $x$, le $z$ e le derivate fino all'ordine $s-1$; avremo così equazioni del tipo
\begin{equation*}\tag{$\beta$}
 z^{(s)}_{i;kk_1\ldots k_s}
 =\psi^{(i,kk_1\ldots k_s)}(x,z,z'_{i;k},\ldots,z^{(s-1)}_{i;k_1\ldots k_{s-1}}),
\end{equation*}
mediante le quali tutte le derivate parziali d'ordine $s$ delle $z$ sono espresse per quelle d'ordine inferiore, delle $z$ e delle $x$. Il sistema di equazioni a derivate parziali $(\alpha),(\beta)$ è appunto quello che ci proponevamo di costruire.

\section*{\S{} 34. Il sistema $(\alpha),(\beta)$ come sistema di Mayer. Proprietà di questi sistemi}

Riconosciamo subito alcune importanti proprietà di questo sistema ricordando che esso è soddisfatto dalle funzioni date
\[
 z_i=f_i(x_1,x_2,\ldots,x_n; a_1,\ldots,a_r)
\]
contenenti $r$ costanti essenziali arbitrarie. Se assumiamo, oltre le $z_i$, come altrettante nuove funzioni incognite, le derivate delle $z$ fino all'ordine $s-1$, potremo sostituire al sistema $(\alpha),(\beta)$ un sistema perfettamente equivalente di equazioni ai differenziali totali per le $\varepsilon_0+\varepsilon_1+\cdots+\varepsilon_{s-1}$ funzioni incognite
\[
 z_i,
 \quad z'_{i;k},
 \quad z'_{i;k_1k_2},
 \quad \ldots,
 \quad z^{(s-1)}_{i;k_1\ldots k_{s-1}},
\]
combinate colle $q$ equazioni in termini finiti fra queste che risultano dalle $(\alpha)$.

Siccome, per la teoria delle equazioni ai differenziali totali, il sistema integrale generale può contenere al massimo
\[
 \varepsilon_0+\cdots+\varepsilon_{s-1}-q=r
\]
costanti arbitrarie e ne avrà precisamente $r$ solo quando il sistema differenziale sia illimitatamente integrabile. Ma questo massimo $r$ è effettivamente raggiunto dalle espressioni $z_i=f_i(x,a)$ (e le loro derivate sino all'ordine $s-1$), onde si conclude che queste espressioni ci danno l'integrale generale del sistema differenziale, il quale deve essere illimitatamente integrabile.

Dopo ciò vediamo che il sistema di equazioni a derivate parziali $(\alpha),(\beta)$ gode delle proprietà seguenti:

a) Il sistema permette di esprimere tutte le derivate parziali d'ordine $s$ delle funzioni incognite $z$ per le derivate d'ordine inferiore, per le $z$ e le $x$ stesse. È impossibile invece esprimere tutte le derivate d'ordine $s-1$ per quelle d'ordine inferiore\begingroup\renewcommand\thefootnote{(*)}\footnotemark\endgroup.
\begingroup\renewcommand\thefootnote{(*)}\footnotetext{Perchè ciò fosse possibile infatti, bisognerebbe che dalle equazioni $(\alpha),(1),\ldots,(1_{s-2})$ si potessero eliminare le $r$ $a$ per sostituirle nelle $(1_{s-1})$; ma ciò non è perchè $\mu_{s-2}<r$.}\endgroup{}

b) Derivando le $(\alpha),(\beta)$ e sostituendovi per le derivate d'ordine $s$ le espressioni $(\beta)$, ed eliminando quelle d'ordine $s+1$, si ottengono sole equazioni che sono conseguenze delle $(\alpha),(\beta)$ stesse\begingroup\renewcommand\thefootnote{(**)}\footnotemark\endgroup. In altro modo si può dire che le conseguenze differenziali del sistema $(\alpha),(\beta)$ fino all'ordine $s+1$ sono già contenute nel sistema stesso.
\begingroup\renewcommand\thefootnote{(**)}\footnotetext{Ciò segue dall'illimitata integrabilità del sistema d'equazioni ai differenziali totali.}\endgroup{}

In generale un sistema di equazioni simultanee a derivate parziali, che goda delle proprietà a), b), dicesi un sistema di Mayer, perchè la sua integrazione si riduce a quella di un sistema di equazioni ai differenziali totali illimitatamente integrabile, sistemi che Mayer specialmente ha considerato. È infatti, se si eliminano come altrettante nuove funzioni incognite le derivate prime, seconde, ecc. fino all'ordine $s-1$ delle $z$, si può sostituire un sistema di equazioni ai differenziali totali che risulta, per la proprietà b), illimitatamente integrabile. Il numero effettivo di costanti arbitrarie che comporta l'integrale generale è sempre dato da
\[
 r=\varepsilon_0+\varepsilon_1+\cdots+\varepsilon_{s-1}-q.
\]

Dagli sviluppi del \S{} precedente risulta in generale che: Se un sistema di equazioni simultanee alle derivate parziali è tale che nella espressione del suo integrale generale entri solo un numero finito di costanti arbitrarie, esso è equivalente ad un sistema di Mayer.

Spesse volte è opportuno sostituire alle equazioni $(\alpha)$ di un sistema di Mayer un sistema equivalente ottenuto nel modo seguente. Può darsi che fra le equazioni del sistema $(\alpha)$ ve ne siano di quelle contenenti solo le $x,z$, ovvero che vi siano di tali equazioni conseguenze del sistema $(\alpha)$. Poniamo che esistano precisamente $\nu_0$ relazioni fra le sole $x,z$ conseguenza del sistema $(\alpha)$. Per riconoscere il valore di $\nu_0$ basta il calcolo della caratteristica di una matrice, cioè della matrice funzionale di $W_1,W_2,\ldots,W_q$ rispetto alle
\[
 z'_{i;k},
 \quad z'_{i;k_1k_2},
 \quad \ldots,
 \quad z^{(s-1)}_{i;k_1\ldots k_{s-1}}.
\]
Se questa caratteristica è data da $\sigma_0$, che è necessariamente $\le q$, sarà
\[
 \nu_0=q-\sigma_0.\begingroup\renewcommand\thefootnote{(*)}\footnotemark\endgroup
\]
\begingroup\renewcommand\thefootnote{(*)}\footnotetext{L'essere $\sigma_0$ quella caratteristica significa che fra le equazioni $W_1,\ldots,W_q$, le $x$ e le $z$ sussistono $q-\sigma_0$ relazioni indipendenti, le quali per $W_1=0,\ldots,W_q=0$ danno altrettante relazioni fra le $x$ e le $z$.}\endgroup{}
Similmente indichiamo con $\nu_1$ il numero delle relazioni fra le
\[
 x,z,z'_{i;k}
\]
che seguono dalle $(\alpha)$, non computando le $\nu_0$ precedenti fra le $x,z$; con $\nu_2$ il numero delle relazioni conseguenti dalle $(\alpha)$ fra le $x,z$ e le derivate prime e seconde. I numeri
\[
 \nu_0,\nu_1,\nu_2,\ldots,\nu_{s-1}
\]
si calcoleranno come il primo $\nu_0$, osservando i valori delle caratteristiche di matrici funzionali, e si avrà naturalmente
\[
 \nu_0+\nu_1+\nu_2+\cdots+\nu_{s-1}=q,
\]
e quindi
\[
 r=(\varepsilon_0-\nu_0)+(\varepsilon_1-\nu_1)+\cdots+(\varepsilon_{s-1}-\nu_{s-1}).
\]
Al sistema $(\beta)$ potremo sostituire un sistema equivalente formato $1^{\circ}$ da $\nu_0$ equazioni che esprimono $\nu_0$ delle $z$ per le rimanenti $z$ e per le $x$; $2^{\circ}$ da $\nu_1$ equazioni che danno $\nu_1$ delle derivate prime delle $z$ espresse per le rimanenti $z'$, per le $z$ e per le $x$ ecc.; in generale $\nu_k$ delle derivate d'ordine $k$ saranno espresse per le rimanenti $\varepsilon_k-\nu_k$ e per quelle d'ordine inferiore $(k=0,1,2,\ldots,s-1)$. È chiaro che sarà sempre
\[
 \nu_k<\varepsilon_k,
\]
perchè se tutte le $\varepsilon_k$ derivate d'ordine $k$ potessero esprimersi per quelle d'ordine inferiore, avremmo, per $k<s$, ciò contraddirebbe alla prima a) delle proprietà del sistema di Mayer.

\section*{\S{} 35. Le equazioni di definizione di un gruppo $G_r$}

Applichiamo ora le osservazioni generali dei due \S\S{} precedenti alla formazione di quel sistema di equazioni a derivate parziali (del tipo di Mayer) cui soddisfano i coefficienti $\xi_i$ della più generale trasformazione infinitesima
\[
 Xf=\sum\limits_i^{1,\ldots,n}\xi_i\frac{\partial f}{\partial x_i}
\]
del gruppo $G_r$ generato dalle $r$ trasformazioni infinitesime
\[
 X_k f=\sum\limits_i^{1,\ldots,n}\xi_{ki}(x_1,\ldots,x_n)\frac{\partial f}{\partial x_i}.
\]
Avendosi
\begin{equation*}\tag{3}
 \xi_i=e_1\xi_{1i}+e_2\xi_{2i}+\cdots+e_r\xi_{ri},
 \qquad i=1,2,\ldots,n,
\end{equation*}
le $\xi$ contengono linearmente ed omogeneamente le $r$ costanti essenziali arbitrarie $e_1,e_2,\ldots,e_r$. Secondo quanto si è detto in generale al \S{}33, si associeranno alle (3) le equazioni
\begin{equation*}\tag{$3_1$}
 \frac{\partial\xi_i}{\partial x_k}
  =e_1\frac{\partial\xi_{1i}}{\partial x_k}
  +e_2\frac{\partial\xi_{2i}}{\partial x_k}
  +\cdots
  +e_r\frac{\partial\xi_{ri}}{\partial x_k},
 \qquad i,k=1,2,\ldots,n.
\end{equation*}


\end{document}
